% GENERATED by bin/curate/gen_tutorials_guide.py -- DO NOT EDIT BY HAND.
% Assembled from each case's controlDict description + header comment.

\subsection{adsorber}
\grouptheory{adsorber}
\clearpage
\tutentry{batch09\_adsorber\_co2}{tutorials/batch/adsorber/batch09\_adsorber\_co2}
\begin{tutscheme}
\begin{tikzpicture}[>={Stealth[length=2mm]}, every node/.style={font=\scriptsize}]
  \node[draw, rounded corners=2pt, align=center, minimum width=1.9cm, minimum height=0.85cm, fill=black!4] (adsorber) at (0.0,-0.0) {\textbf{adsorber}\\[-1pt]{\tiny batchAdsorber}};
\end{tikzpicture}
\end{tutscheme}
\tutwhat{Closed isothermal CO2 uptake on zeolite 13X (LDF, solid film): the vessel relaxes to the inventory equilibrium p\_inf = 7229.13 Pa, q\_inf = 1.9708 mol/kg}
\begin{tutnarrative}
closed, isothermal gas--solid uptake vessel.

2 mol of pure CO2 in a 0.01 m3 ideal-gas headspace over 1 kg of regenerated zeolite 13X at 298.15 K. Linear-driving-force kinetics on the solid film (dq/dt = k$\ast$(q$\ast$ - q), k = 0.05 1/s, equipment data, this case's sample) against the competitive-Langmuir equilibrium record of the pair catalogue (here a single pair, CO2/zeolite13X: q\_max = 5.0 mol/kg, b\_298 = 9.0 1/bar, dH\_ads = -45 kJ/mol).

The final state is FIXED BY INVENTORY, not by kinetics: with n\_gas = p V/(R T) and q = q\_max b (p/1e5)/(1 + b p/1e5),

p V/(R T) + m\_ads $\ast$ q\_max $\ast$ b $\ast$ (p/1e5)/(1 + b$\ast$(p/1e5)) = n0

whose root at 298.15 K is (a2\_vectors.py, 10 digits)

p\_inf = 7229.128209 Pa q\_inf = 1.970838025 mol/kg n\_gas = 0.02916197467 mol uptake = 1.970838025 mol (98.54 \%) tau (linearised at the equilibrium) = 0.4768 s

The run must reproduce these to the golden tolerances -- an independent analytic anchor, not a mirror of the code. The ODE state is ONLY q; n\_gas is algebraic (n\_tot - m\_ads$\ast$q), so the per-species closure column closure\_worst\_rel in trajectory.csv is exact by construction (<= 1e-13). Log: the unit announces the ideal-gas assumption, the initialLoading default, the RK4 stiffness bound and the isosteric-KPI label at start-up.

\end{tutnarrative}
\begin{tutgolden}
\goldrow{kpi}{adsorber}{P\_final}{7229.128209}
\goldrow{kpi}{adsorber}{Q\_ads\_isosteric\_kJ}{-88.68771114}
\goldrow{kpi}{adsorber}{T\_final}{298.15}
\goldrow{kpi}{adsorber}{closure\_worst\_rel}{0}
\goldrow{kpi}{adsorber}{n\_CO2\_final}{2.916197467e-05}
\goldrow{kpi}{adsorber}{q\_CO2\_final}{1.970838025}
\goldrow{kpi}{adsorber}{t\_end}{600}
\goldrow{kpi}{adsorber}{totalMoles\_final}{2.916197467e-05}
\goldrow{kpi}{adsorber}{uptake\_CO2\_mol}{1.970838025}
\goldrow{kpi}{campaign}{element\_C\_closure\_rel}{0}
\goldrow{kpi}{campaign}{element\_O\_closure\_rel}{0}
\goldrow{kpi}{campaign}{element\_worst\_closure\_rel}{0}
\goldrow{kpi}{campaign}{energy\_balance\_available}{0}
\goldrow{kpi}{campaign}{energy\_records\_invalid}{0}
\goldmore{12}
\end{tutgolden}

\clearpage
\tutentry{batch10\_adsorber\_co2\_n2}{tutorials/batch/adsorber/batch10\_adsorber\_co2\_n2}
\begin{tutscheme}
\begin{tikzpicture}[>={Stealth[length=2mm]}, every node/.style={font=\scriptsize}]
  \node[draw, rounded corners=2pt, align=center, minimum width=1.9cm, minimum height=0.85cm, fill=black!4] (adsorber) at (0.0,-0.0) {\textbf{adsorber}\\[-1pt]{\tiny batchAdsorber}};
\end{tikzpicture}
\end{tutscheme}
\tutwhat{Closed isothermal competitive CO2/N2 uptake on zeolite 13X: extended-Langmuir site competition drives the equilibrium selectivity S = 71.43}
\begin{tutnarrative}
COMPETITIVE binary uptake, closed vessel.

0.5 mol CO2 + 1.5 mol N2 in a 0.01 m3 ideal-gas headspace over 1 kg of regenerated zeolite 13X at 298.15 K. Both species carry an isotherm record (extended Langmuir: CO2 q\_max 5.0 / b\_298 9.0 1/bar; N2 3.5 / 0.18 1/bar) and both carry a declared LDF constant -- so both integrate, and the strong adsorbate suppresses the weak one through the SHARED site denominator 1 + Sum\_j b\_j p\_j.

The final state solves the 2-species inventory system n\_gas\_i + m\_ads $\ast$ q\_i(p\_CO2, p\_N2) = n0\_i, p\_i = n\_gas\_i R T / V whose root (a2\_vectors.py, Newton residual 0):

CO2: p\_inf = 0.01611568794 bar q = 0.4934990047 mol/kg N2 : p\_inf = 1.802652166 bar q = 0.7728183086 mol/kg P\_total = 181876.7854 Pa selectivity S = (q\_CO2/q\_N2)/(p\_CO2/p\_N2) = 71.42857143 = (q\_max$\ast$b)\_CO2/(q\_max$\ast$b)\_N2 (extended-Langmuir identity)

The golden pins the COMPETITIVE path: get the shared denominator wrong and every one of these numbers moves. closure\_worst\_rel stays exact by construction (state = q only; n\_gas algebraic).

\end{tutnarrative}
\begin{tutgolden}
\goldrow{kpi}{adsorber}{P\_final}{181876.7854}
\goldrow{kpi}{adsorber}{Q\_ads\_isosteric\_kJ}{-36.1181845604}
\goldrow{kpi}{adsorber}{T\_final}{298.15}
\goldrow{kpi}{adsorber}{closure\_worst\_rel}{0}
\goldrow{kpi}{adsorber}{n\_CO2\_final}{6.500995276e-06}
\goldrow{kpi}{adsorber}{n\_N2\_final}{7.271816914e-04}
\goldrow{kpi}{adsorber}{q\_CO2\_final}{0.4934990047}
\goldrow{kpi}{adsorber}{q\_N2\_final}{0.7728183086}
\goldrow{kpi}{adsorber}{t\_end}{600}
\goldrow{kpi}{adsorber}{totalMoles\_final}{0.000733682686703}
\goldrow{kpi}{adsorber}{uptake\_CO2\_mol}{0.4934990047}
\goldrow{kpi}{adsorber}{uptake\_N2\_mol}{0.7728183086}
\goldrow{kpi}{campaign}{element\_C\_closure\_rel}{0}
\goldrow{kpi}{campaign}{element\_N\_closure\_rel}{0}
\goldmore{16}
\end{tutgolden}

\clearpage
\tutentry{batch11\_adsorber\_henry}{tutorials/batch/adsorber/batch11\_adsorber\_henry}
\begin{tutscheme}
\begin{tikzpicture}[>={Stealth[length=2mm]}, every node/.style={font=\scriptsize}]
  \node[draw, rounded corners=2pt, align=center, minimum width=1.9cm, minimum height=0.85cm, fill=black!4] (adsorber) at (0.0,-0.0) {\textbf{adsorber}\\[-1pt]{\tiny batchAdsorber}};
\end{tikzpicture}
\end{tutscheme}
\tutwhat{Closed isothermal Henry-limit uptake (linear isotherm): the coupled gas+solid system is LINEAR and the run must match q(t) = q\_inf$\ast$(1 - exp(-k\_eff t)) to 1e-8}
\begin{tutnarrative}
the EXACT-solution gate (G1a of the A2 spec).

1 mol of CO2 in a 0.01 m3 ideal-gas headspace over 1 kg of a SYNTHETIC Henry-limit sorbent (linear isotherm q$\ast$ = K\_H p, K\_H = 0.4 mol/(kg bar) at 298.15 K -- a teaching record, vendored case-local). With a linear isotherm and one species the coupled vessel is LINEAR and closes in analytic form:

alpha = K\_H[mol/(kg Pa)] $\ast$ R$\ast$T/V = 0.9915828118 1/kg k\_eff = k $\ast$ (1 + alpha$\ast$m\_ads) = 0.09957914059 1/s q\_inf = alpha$\ast$n0 / (1 + alpha$\ast$m) = 0.4978868094 mol/kg q(t) = q\_inf $\ast$ (1 - exp(-k\_eff$\ast$t)) (q0 = 0)

q(5) = 0.1952670621 q(10) = 0.3139520079 q(30) = 0.4727835287 q(60) = 0.4966211107 mol/kg

The pedagogical point the closed vessel adds over the fixed-q$\ast$ packet formula: k\_eff > k -- gas depletion ACCELERATES the approach to equilibrium (the packet's q$\ast$ = const law is the V\_gas $\to$ infinity limit of this same solution). The trajectory rows at t = 5/10/30/60 s must match the table above to 1e-8 mol/kg (RK4 dt = 0.01 s); the golden pins the t = 60 s state.

\end{tutnarrative}
\begin{tutgolden}
\goldrow{kpi}{adsorber}{P\_final}{124785.4636}
\goldrow{kpi}{adsorber}{Q\_ads\_isosteric\_kJ}{-9.93242221329}
\goldrow{kpi}{adsorber}{T\_final}{298.15}
\goldrow{kpi}{adsorber}{closure\_worst\_rel}{0}
\goldrow{kpi}{adsorber}{n\_CO2\_final}{5.033788893e-04}
\goldrow{kpi}{adsorber}{q\_CO2\_final}{0.4966211107}
\goldrow{kpi}{adsorber}{t\_end}{60}
\goldrow{kpi}{adsorber}{totalMoles\_final}{5.033788893e-04}
\goldrow{kpi}{adsorber}{uptake\_CO2\_mol}{0.4966211107}
\goldrow{kpi}{campaign}{element\_C\_closure\_rel}{0}
\goldrow{kpi}{campaign}{element\_O\_closure\_rel}{0}
\goldrow{kpi}{campaign}{element\_worst\_closure\_rel}{0}
\goldrow{kpi}{campaign}{energy\_balance\_available}{0}
\goldrow{kpi}{campaign}{energy\_records\_invalid}{0}
\goldmore{12}
\end{tutgolden}

\clearpage
\tutentry{batch12\_adsorber\_refused\_P}{tutorials/batch/adsorber/batch12\_adsorber\_refused\_P}
\begin{tutscheme}
\begin{tikzpicture}[>={Stealth[length=2mm]}, every node/.style={font=\scriptsize}]
  \node[draw, rounded corners=2pt, align=center, minimum width=1.9cm, minimum height=0.85cm, fill=black!4] (adsorber) at (0.0,-0.0) {\textbf{adsorber}\\[-1pt]{\tiny batchAdsorber}};
\end{tikzpicture}
\end{tutscheme}
\tutwhat{REFUSAL lesson: declared initial.P (10 bar) contradicts n$\ast$R$\ast$T/V (4.958 bar) -- the honest cross-check FATALs; P of a closed ideal-gas headspace is a result, not an input}
\begin{tutnarrative}
the REFUSAL lesson (deliberate, must FAIL).

Identical to batch09\_adsorber\_co2 EXCEPT initial.P declares 10 bar while n$\ast$R$\ast$T/V = 4.9579 bar. P of a closed ideal-gas headspace is a RESULT (n, T, V fix it -- Duhem); a declared P is honoured only as an honest cross-check, and a relative deviation > 1e-3 REFUSES, printing both numbers. Expected FATAL (exit != 0):

batchAdsorber 'adsorber': declared P inconsistent with n$\ast$R$\ast$T/V -- initial.P = 1e+06 Pa (10 bar) but n$\ast$R$\ast$T/V = 495791 Pa (4.95791 bar), relative deviation 1.01698 > 1e-3. ...

Marked .expect-nonconvergence: counted EXPECTED-FAIL by runTests, never green, never recordable. Base case (the passing twin):

2 mol of pure CO2 in a 0.01 m3 ideal-gas headspace over 1 kg of regenerated zeolite 13X at 298.15 K. Linear-driving-force kinetics on the solid film (dq/dt = k$\ast$(q$\ast$ - q), k = 0.05 1/s, equipment data, this case's sample) against the competitive-Langmuir equilibrium record of the pair catalogue (here a single pair, CO2/zeolite13X: q\_max = 5.0 mol/kg, b\_298 = 9.0 1/bar, dH\_ads = -45 kJ/mol).

The final state is FIXED BY INVENTORY, not by kinetics: with n\_gas = p V/(R T) and q = q\_max b (p/1e5)/(1 + b p/1e5),

p V/(R T) + m\_ads $\ast$ q\_max $\ast$ b $\ast$ (p/1e5)/(1 + b$\ast$(p/1e5)) = n0

whose root at 298.15 K is (a2\_vectors.py, 10 digits)

p\_inf = 7229.128209 Pa q\_inf = 1.970838025 mol/kg n\_gas = 0.02916197467 mol uptake = 1.970838025 mol (98.54 \%) tau (linearised at the equilibrium) = 0.4768 s

The run must reproduce these to the golden tolerances -- an independent analytic anchor, not a mirror of the code. The ODE state is ONLY q; n\_gas is algebraic (n\_tot - m\_ads$\ast$q), so the per-species closure column closure\_worst\_rel in trajectory.csv is exact by construction (<= 1e-13). Log: the unit announces the ideal-gas assumption, the initialLoading default, the RK4 stiffness bound and the isosteric-KPI label at start-up.

\end{tutnarrative}

\clearpage
\tutentry{batch13\_breakthrough\_co2}{tutorials/batch/adsorber/batch13\_breakthrough\_co2}
\begin{tutscheme}
\begin{tikzpicture}[>={Stealth[length=2mm]}, every node/.style={font=\scriptsize}]
  \node[draw, rounded corners=2pt, align=center, minimum width=1.9cm, minimum height=0.85cm, fill=black!4] (bed) at (0.0,-0.0) {\textbf{bed}\\[-1pt]{\tiny fixedBedAdsorber}};
\end{tikzpicture}
\end{tutscheme}
\tutwhat{CO2 15\%/He breakthrough on zeolite 13X -- 1-D isothermal fixed bed (A3): t\_st = 3040.5 s announced pre-run, t\_b(5\%) = 2934.1 s on the second-order van Leer scheme (2890.7 s under first-order upwind), ledger closure at machine level}
\begin{tutnarrative}
CO2/He breakthrough on a zeolite-13X bed.

THE A3 anchor case (spec 2026-07-12): a 0.5 m isothermal fixed bed of zeolite 13X, fed 15\% CO2 in helium at u = 0.05 m/s superficial, 1 bar, 298 K. Helium carries NO isotherm record on 13X $\to$ announced inert CARRIER, transported but never adsorbed, closed by difference from the constant c\_tot = P/(RT). At 15\% CO2 this is NOT a dilute approximation: the constant-(u, P, T) closure implies He at the net uptake rate, and the campaign books it as a NAMED INLET STREAM of the bed, `carrierMakeUp` ("carrier implied by the constant-(u, P, T) assumption"), so mass and element balances close and the stream is the explanation; its size is the assumption's size (KPIs carrier\_fabricated\_mol, carrier\_makeup\_rel); the genuinely clean transport anchor is batch14. The engine announces, BEFORE integrating, where the front must arrive:

q$\ast$(c\_in) = 5$\ast$9$\ast$0.15/(1 + 9$\ast$0.15) = 2.872340426 mol/kg R\_f = eps + rho\_b q$\ast$ / c\_in = 304.0512509 t\_st = (L/u) R\_f = 3040.512509 s u\_sh = u/R\_f = 1.644459605e-4 m/s

The front this run computes (convectionScheme vanLeer, the stated default since C375; nCells = 100):

t\_b(5\%) = 2934.069 s t\_50 = 3035.101 s t\_b(95\%) = 3165.565 s 5-95 \% width 231.5 s; unused bed length L(1 - t\_b/t\_st) = 0.0175 m INT(1 - c\_out/c\_in)dt = t\_st to 3e-8 (telescopic FV conservation) ledger mass closure at round-off (IN/OUT rows of the SAME ODE)

That is the PHYSICAL front within the stated tolerance: the mesh- converged front (Richardson, 800/1600 cells) has t\_b(5\%) = 2939.26 s, width 219.6 s, unused length 0.0167 m, so at 100 cells van Leer is 0.18 \% early on t\_b and 5.4 \% wide (gate check\_fixed\_bed\_scheme; the width converges at order \textasciitilde{}2). THE COMPARISON A STUDENT CAN REPRODUCE: declare `convectionScheme upwind;` and the first-order scheme adds its numerical dispersion v dz/2 = 3.1e-4 m2/s, 1.25 times the physical Dax/eps = 2.5e-4 (the run announces it), and the front becomes the grid's: t\_b(5\%) = 2890.669 s, t\_50 = 3036.448 s, t\_b(95\%) = 3204.166 s, width 313.5 s, unused length 0.0246 m (the A3 anchors of anchors\_a3.py, FV N = 100, observed order \textasciitilde{}1.03-1.05 -- batch15 is that mesh study). The stoichiometric time does not move with either scheme or mesh.

Numerics: conservative FV, nCells = 100 DECLARED, Danckwerts inlet as the imposed face flux u$\ast$c\_in, outlet dc/dz = 0. u, P, T are DECLARED constants of this phase, and the price is said OUT LOUD: the real velocity drops \textasciitilde{}15\% across the uptake zone, so the constant-u closure FABRICATES carrier (He) at exactly the net uptake rate (9.19 mol here) -- announced in the run header, pinned as the KPI carrier\_fabricated\_mol and booked as the `carrierMakeUp` inlet stream of the campaign ledger; Ergun (`flowModel ergun`) needs no such stream, and the energy balance of this closure stays UNAVAILABLE (named).

Log: run header prints the hypotheses, the derived rho\_p, Pe = 250, the Gershgorin step bound term by term, the scheme's numerical dispersion and the stiffness verdict; the initial profile lands in 0/bed.profile (the later <t>/bed.profile files are erased by the instant writer -- a defect recorded in docs/design/fixed-bed-numerics-c375.md); trajectory.csv carries c\_out, y\_out, qbar and the per-species ledger closure.

\end{tutnarrative}
\begin{tutgolden}
\goldrow{kpi}{bed}{E\_adsorption\_seg0\_kJ}{-413.617021276}
\goldrow{kpi}{bed}{E\_adsorption\_total\_kJ}{-413.617021276}
\goldrow{kpi}{bed}{P\_final}{100000}
\goldrow{kpi}{bed}{Pe}{250}
\goldrow{kpi}{bed}{T\_final}{298}
\goldrow{kpi}{bed}{c\_out\_over\_cin\_final\_CO2}{0.999999999943}
\goldrow{kpi}{bed}{carrier\_fabricated\_mol}{9.19148936169}
\goldrow{kpi}{bed}{integral\_anchor\_CO2}{3040.51243456}
\goldrow{kpi}{bed}{mass\_closure\_CO2}{6.86742033269e-14}
\goldrow{kpi}{bed}{n\_CO2\_final}{1.21079552058e-05}
\goldrow{kpi}{bed}{n\_He\_final}{6.86117461666e-05}
\goldrow{kpi}{bed}{qbar\_CO2\_final}{2.87234042553}
\goldrow{kpi}{bed}{retention\_factor\_CO2}{304.051250948}
\goldrow{kpi}{bed}{t\_50\_CO2}{3035.10088974}
\goldmore{44}
\end{tutgolden}

\clearpage
\tutentry{batch14\_transport\_only}{tutorials/batch/adsorber/batch14\_transport\_only}
\begin{tutscheme}
\begin{tikzpicture}[>={Stealth[length=2mm]}, every node/.style={font=\scriptsize}]
  \node[draw, rounded corners=2pt, align=center, minimum width=1.9cm, minimum height=0.85cm, fill=black!4] (bed) at (0.0,-0.0) {\textbf{bed}\\[-1pt]{\tiny fixedBedAdsorber}};
\end{tikzpicture}
\end{tutscheme}
\tutwhat{Pure transport (k = 0) on the batch13 bed, N = 200: t\_50 vs the Ogata-Banks analytic ADE solution (gate G1) under fixed RK4 with the Gershgorin bound printed term by term}
\begin{tutnarrative}
the k = 0 transport gate (A3 spec, gate G1).

Same bed, same feed, same mesh family as batch13 -- but the LDF constant is DECLARED ZERO: the solid is frozen and CO2 crosses the bed as a pure tracer. What is left is exactly the advection-dispersion equation with a Danckwerts inlet, and THAT has an analytic yardstick (Ogata-Banks, semi-infinite ADE; the finite-domain difference is O(1/Pe) = 0.4\% at Pe = 250):

u\_i = u/eps = 0.125 m/s D\_i = Dax/eps = 2.5e-4 m2/s c/c0 (L, 3.5 s) = 0.07340274358 c/c0 (L, 4.0 s) = 0.5178057707 c/c0 (L, 4.5 s) = 0.913645218 t\_50 analytic = 3.984074227 s

FV yardsticks pinned by this case (anchors\_a3.py, N = 200):

t\_50 numeric = 3.975611088 s (0.21\% below analytic -- first-order upwind; the deviation HALVES with each mesh doubling: 1.40\% at N=25, 0.74\% at 50, 0.40\% at 100, 0.21\% at 200, 0.11\% at 400) INT(1 - c\_out/c\_in)dt = eps L / u = 4.0 s (telescopic conservation: exact identity once saturated) ledger closure \textasciitilde{} 1e-14 (k = 0)

A declared zero is not a silent default: the unit announces "k = 0 DECLARED: transport-only diagnostic, the solid is frozen (gate G1)". Note the He carrier fabrication of batch13 vanishes here (no uptake): carrier\_fabricated\_mol = 0.

\end{tutnarrative}
\begin{tutgolden}
\goldrow{kpi}{bed}{E\_adsorption\_seg0\_kJ}{0}
\goldrow{kpi}{bed}{E\_adsorption\_total\_kJ}{0}
\goldrow{kpi}{bed}{P\_final}{100000}
\goldrow{kpi}{bed}{Pe}{250}
\goldrow{kpi}{bed}{T\_final}{298}
\goldrow{kpi}{bed}{c\_out\_over\_cin\_final\_CO2}{0.999999999893}
\goldrow{kpi}{bed}{carrier\_fabricated\_mol}{0}
\goldrow{kpi}{bed}{integral\_anchor\_CO2}{3.99999999999}
\goldrow{kpi}{bed}{mass\_closure\_CO2}{4.09039795713e-14}
\goldrow{kpi}{bed}{n\_CO2\_final}{1.21079552058e-05}
\goldrow{kpi}{bed}{n\_He\_final}{6.86117461666e-05}
\goldrow{kpi}{bed}{qbar\_CO2\_final}{0}
\goldrow{kpi}{bed}{t\_50\_CO2}{3.97561121931}
\goldrow{kpi}{bed}{t\_breakthrough\_5pct\_CO2}{3.29598772214}
\goldmore{27}
\end{tutgolden}

\clearpage
\tutentry{batch15\_mesh\_study}{tutorials/batch/adsorber/batch15\_mesh\_study}
\begin{tutscheme}
\begin{tikzpicture}[>={Stealth[length=2mm]}, every node/.style={font=\scriptsize}]
  \node[draw, rounded corners=2pt, align=center, minimum width=1.9cm, minimum height=0.85cm, fill=black!4] (bed) at (0.0,-0.0) {\textbf{bed}\\[-1pt]{\tiny fixedBedAdsorber}};
\end{tikzpicture}
\end{tutscheme}
\tutwhat{Mesh-study member N = 25 of the batch13 breakthrough: t\_b(5\%) converges monotonically 2732 $\to$ 2839 $\to$ 2891 $\to$ 2916 s with observed order 1.03-1.05; full table in reports/meshStudy.md}
\begin{tutnarrative}
the declared-mesh lesson (A3 spec, gate G6).

The axial mesh is a DECLARED modelling choice (nCells; absent = FATAL), and this case is the receipt: the batch13 breakthrough run on the 25/50/100/200 family, everything else identical, with the scheme PINNED to `convectionScheme upwind;` (batch13 itself runs the default van Leer scheme since C375; its t\_b5 at 100 cells is 2934.07 s). First-order upwind advection dominates the error, so t\_b(5\%) must converge MONOTONICALLY with observed order p \textasciitilde{} 1:

N = 25 t\_b5 = 2732.027658 s (this case) N = 50 t\_b5 = 2838.542518 s N = 100 t\_b5 = 2890.669098 s (batch13's mesh, under upwind) N = 200 t\_b5 = 2915.764796 s p(25/50/100) = 1.031 p(50/100/200) = 1.055 Richardson (p = 1): t\_b5(inf) \textasciitilde{} 2940.86 s $\to$ N = 100 sits 1.7\% low

The full table, the four overlaid breakthrough curves and the verdict are PUBLISHED in reports/meshStudy.md (generated from the four runs; the report header cites mesh, tolerances and integrator). The ledger mass closure stays <= \textasciitilde{}1e-11 on EVERY mesh -- conservation is telescopic in the FV form, it does not improve with refinement because it was never discretisation-limited (gate G3).

This case pins the N = 25 member; its numbers are DELIBERATELY the worst of the family (t\_b5 5.8\% below batch13) -- coarse-mesh optimism made visible, the point of the lesson.

\end{tutnarrative}
\begin{tutgolden}
\goldrow{kpi}{bed}{E\_adsorption\_seg0\_kJ}{-413.616606535}
\goldrow{kpi}{bed}{E\_adsorption\_total\_kJ}{-413.616606535}
\goldrow{kpi}{bed}{P\_final}{100000}
\goldrow{kpi}{bed}{Pe}{250}
\goldrow{kpi}{bed}{T\_final}{298}
\goldrow{kpi}{bed}{c\_out\_over\_cin\_final\_CO2}{0.999973635802}
\goldrow{kpi}{bed}{carrier\_fabricated\_mol}{9.19148014522}
\goldrow{kpi}{bed}{integral\_anchor\_CO2}{3040.50945345}
\goldrow{kpi}{bed}{mass\_closure\_CO2}{7.40535661511e-15}
\goldrow{kpi}{bed}{n\_CO2\_final}{1.21079316034e-05}
\goldrow{kpi}{bed}{n\_He\_final}{6.86117697691e-05}
\goldrow{kpi}{bed}{qbar\_CO2\_final}{2.87233754538}
\goldrow{kpi}{bed}{retention\_factor\_CO2}{304.051250948}
\goldrow{kpi}{bed}{t\_50\_CO2}{3041.40641271}
\goldmore{41}
\end{tutgolden}

\clearpage
\tutentry{batch16\_ergun\_profile}{tutorials/batch/adsorber/batch16\_ergun\_profile}
\begin{tutscheme}
\begin{tikzpicture}[>={Stealth[length=2mm]}, every node/.style={font=\scriptsize}]
  \node[draw, rounded corners=2pt, align=center, minimum width=1.9cm, minimum height=0.85cm, fill=black!4] (bed) at (0.0,-0.0) {\textbf{bed}\\[-1pt]{\tiny fixedBedAdsorber}};
\end{tikzpicture}
\end{tutscheme}
\tutwhat{A4 F1: pure-He Ergun profile against the closed-form P-squared solution}
\begin{tutgolden}
\goldrow{kpi}{bed}{P\_first\_cell\_Pa}{99999.4670134}
\goldrow{kpi}{bed}{P\_mid\_cell\_Pa}{99946.1540012}
\goldrow{kpi}{bed}{P\_last\_cell\_Pa}{99893.8796441}
\goldrow{kpi}{bed}{P\_out\_boundary\_Pa}{99893.3459257}
\goldrow{kpi}{bed}{carrier\_fabricated\_mol}{0}
\goldrow{kpi}{bed}{carrier\_makeup\_rel}{0}
\goldrow{kpi}{campaign}{mass\_closure\_rel}{0}
\goldrow{kpi}{bed}{E\_adsorption\_seg0\_kJ}{0}
\goldrow{kpi}{bed}{E\_adsorption\_total\_kJ}{0}
\goldrow{kpi}{bed}{P\_final}{99946.6820773}
\goldrow{kpi}{bed}{T\_final}{298.15}
\goldrow{kpi}{bed}{deltaP\_first\_to\_out\_Pa}{106.12034902}
\goldrow{kpi}{bed}{gasInventory\_mol}{0.0806360746763}
\goldrow{kpi}{bed}{n\_He\_final}{8.06360746763e-05}
\goldmore{34}
\end{tutgolden}

\clearpage
\tutentry{batch17\_dilute\_ergun\_limit}{tutorials/batch/adsorber/batch17\_dilute\_ergun\_limit}
\begin{tutscheme}
\begin{tikzpicture}[>={Stealth[length=2mm]}, every node/.style={font=\scriptsize}]
  \node[draw, rounded corners=2pt, align=center, minimum width=1.9cm, minimum height=0.85cm, fill=black!4] (legacy) at (0.0,-0.0) {\textbf{legacy}\\[-1pt]{\tiny fixedBedAdsorber}};
  \node[draw, rounded corners=2pt, align=center, minimum width=1.9cm, minimum height=0.85cm, fill=black!4] (ergun) at (0.0,-1.7) {\textbf{ergun}\\[-1pt]{\tiny fixedBedAdsorber}};
\end{tikzpicture}
\end{tutscheme}
\tutwhat{A4 F2: 400 ppm uptake, Ergun must approach the A3 constant-velocity limit}
\begin{tutgolden}
\goldrow{kpi}{campaign}{element\_C\_closure\_rel}{5.33128627003e-13}
\goldrow{kpi}{campaign}{element\_He\_closure\_rel}{2.669253225e-13}
\goldrow{kpi}{campaign}{element\_O\_closure\_rel}{5.33128627003e-13}
\goldrow{kpi}{campaign}{element\_worst\_closure\_rel}{5.33128627003e-13}
\goldrow{kpi}{campaign}{energy\_balance\_available}{0}
\goldrow{kpi}{campaign}{energy\_records\_invalid}{0}
\goldrow{kpi}{campaign}{energy\_records\_valid}{3}
\goldrow{kpi}{campaign}{mass\_closure\_rel}{2.6805259481e-13}
\goldrow{kpi}{campaign}{mass\_kg\_external\_out}{0.161630811285}
\goldrow{kpi}{campaign}{mass\_kg\_final}{0.000565140820198}
\goldrow{kpi}{campaign}{mass\_kg\_initial}{0.162173202362}
\goldrow{kpi}{campaign}{mass\_residual\_kg}{4.34770458296e-14}
\goldrow{kpi}{campaign}{moles\_kmol\_external\_out}{0.0403338611118}
\goldrow{kpi}{campaign}{moles\_kmol\_final}{2.7503993825e-05}
\goldmore{80}
\end{tutgolden}

\clearpage
\tutentry{batch18\_ergun\_conservation}{tutorials/batch/adsorber/batch18\_ergun\_conservation}
\begin{tutscheme}
\begin{tikzpicture}[>={Stealth[length=2mm]}, every node/.style={font=\scriptsize}]
  \node[draw, rounded corners=2pt, align=center, minimum width=1.9cm, minimum height=0.85cm, fill=black!4] (bed) at (0.0,-0.0) {\textbf{bed}\\[-1pt]{\tiny fixedBedAdsorber}};
\end{tikzpicture}
\end{tutscheme}
\tutwhat{A4 F3: 15 percent CO2 with explicit carrier continuity and Ergun velocity}
\begin{tutgolden}
\goldrow{kpi}{bed}{E\_adsorption\_seg0\_kJ}{-413.69985154}
\goldrow{kpi}{bed}{E\_adsorption\_total\_kJ}{-413.69985154}
\goldrow{kpi}{bed}{P\_final}{100047.159985}
\goldrow{kpi}{bed}{P\_first\_cell\_Pa}{100092.419862}
\goldrow{kpi}{bed}{P\_last\_cell\_Pa}{100001.886963}
\goldrow{kpi}{bed}{P\_mid\_cell\_Pa}{100047.163684}
\goldrow{kpi}{bed}{P\_out\_boundary\_Pa}{100000}
\goldrow{kpi}{bed}{Pe}{250}
\goldrow{kpi}{bed}{T\_final}{298.15}
\goldrow{kpi}{bed}{c\_out\_over\_cin\_final\_CO2}{1.0000135596}
\goldrow{kpi}{bed}{carrier\_fabricated\_mol}{-7.95807864051e-13}
\goldrow{kpi}{bed}{deltaP\_first\_to\_out\_Pa}{92.4198624974}
\goldrow{kpi}{bed}{integral\_anchor\_CO2}{3032.65647404}
\goldrow{kpi}{bed}{mass\_closure\_CO2}{2.15282822286e-14}
\goldmore{51}
\end{tutgolden}

\clearpage
\tutentry{batch19\_feed\_switch\_purge}{tutorials/batch/adsorber/batch19\_feed\_switch\_purge}
\begin{tutscheme}
\begin{tikzpicture}[>={Stealth[length=2mm]}, every node/.style={font=\scriptsize}]
  \node[draw, rounded corners=2pt, align=center, minimum width=1.9cm, minimum height=0.85cm, fill=black!4] (bed) at (0.0,-0.0) {\textbf{bed}\\[-1pt]{\tiny fixedBedAdsorber}};
\end{tikzpicture}
\end{tutscheme}
\tutwhat{A5 feed switch: CO2/He load to 3600 s (batch13 anchors), then pure-He purge to 9000 s -- amended feed commitment ledgered, frozen breakthrough KPIs, desorption duty as the same state difference}
\begin{tutnarrative}
THE A5 WITNESS: load, then switch the feed.

batch13's bed (CO2 15\% in He on zeolite 13X, the A3 anchor), run through a TWO-STEP cycle in one campaign:

0 .. 3600 s LOAD -- identical to batch13's feed; the front arrives on batch13's front (t\_5 = 2934.1, t\_50 = 3035.1, t\_95 = 3165.6 s on the default van Leer scheme, C375) 3600 s SWITCH -- recipe setParameter bed.feed.CO2 = 0: the feed becomes pure He (concentration- swing regeneration) 3600 .. 9000 s PURGE -- q$\ast$(c $\to$ 0) = 0, the LDF drives the bed to desorb; the favourable (Langmuir) isotherm makes desorption the UNFAVOURABLE direction -- watch the long tail on c\_out

What the switch does, and where each consequence is VISIBLE:

The feed commitment is DECLARED matter (the campaign datum owns A$\ast$u$\ast$c\_in$\ast$(endTime - t) of future feed). Re-declaring the feed at t = 3600 s AMENDS that commitment: the un-fed CO2 leaves the datum and the extra He enters it, priced and recorded as `feedAmendment` ledger records against the external boundary -- the campaign mass/element balances read the records and still close (KPIs mass\_kg\_amend\_in/out, element\_worst\_closure\_vs\_declared\_rel \textasciitilde{} 1e-14).

The breakthrough sampler's reference basis (c\_out/c\_in) dies with the old feed: crossings and integral FREEZE at their pre-switch values (announced), so t\_5/t\_50/t\_95 here must MATCH batch13's anchors -- the frozen claim describes the first feed step. The desorption wave lives on in the trajectory's c\_out\_CO2 / qbar\_CO2 columns.

The adsorption duty stays an EXACT state difference on the adsorbed inventory: E\_adsorption = dH\_ads $\ast$ (n\_ads(end) - n\_ads(0)). batch13 (load only) ledgers -413.6 kJ; here (-20.3 kJ at 9000 s) the purge takes part of the loading back OUT, so |E| is SMALLER -- desorption is endothermic and the ledger shows the heat returning, with no new formula anywhere.

u, P, T stay DECLARED CONSTANTS: asking the bed for a thermal swing (feed.T), a flow transient (u) or a blowdown (P) is REFUSED by name (fired by bin/curate/check\_feed\_switch.py, batch12-style honesty).

The carrier-fabrication caveat of the constant-(u, P, T) closure applies unchanged: the carrier it implies is booked as the named stream `carrierMakeUp` (an inlet while the bed loads, an outlet on the purge's net desorption), so the He element closure holds and the stream is the explanation.

\end{tutnarrative}
\begin{tutgolden}
\goldrow{kpi}{bed}{E\_adsorption\_seg0\_kJ}{-20.3329948533}
\goldrow{kpi}{bed}{E\_adsorption\_total\_kJ}{-20.3329948533}
\goldrow{kpi}{bed}{P\_final}{100000}
\goldrow{kpi}{bed}{Pe}{250}
\goldrow{kpi}{bed}{T\_final}{298}
\goldrow{kpi}{bed}{carrier\_fabricated\_mol}{0.451844330074}
\goldrow{kpi}{bed}{integral\_anchor\_CO2}{3040.51207995}
\goldrow{kpi}{bed}{mass\_closure\_CO2}{8.63958271763e-14}
\goldrow{kpi}{bed}{n\_CO2\_final}{2.75111922805e-07}
\goldrow{kpi}{bed}{n\_He\_final}{8.04445894496e-05}
\goldrow{kpi}{bed}{qbar\_CO2\_final}{0.141201353148}
\goldrow{kpi}{bed}{t\_50\_CO2}{3035.10085907}
\goldrow{kpi}{bed}{t\_breakthrough\_5pct\_CO2}{2934.06908679}
\goldrow{kpi}{bed}{t\_breakthrough\_95pct\_CO2}{3165.56529487}
\goldmore{52}
\end{tutgolden}

\clearpage
\tutentry{batch20\_thermal\_breakthrough}{tutorials/batch/adsorber/batch20\_thermal\_breakthrough}
\begin{tutscheme}
\begin{tikzpicture}[>={Stealth[length=2mm]}, every node/.style={font=\scriptsize}]
  \node[draw, rounded corners=2pt, align=center, minimum width=1.9cm, minimum height=0.85cm, fill=black!4] (bed) at (0.0,-0.0) {\textbf{bed}\\[-1pt]{\tiny fixedBedAdsorber}};
\end{tikzpicture}
\end{tutscheme}
\tutwhat{A5-T1 adiabatic CO2/He breakthrough (ergun, one T per cell): thermal wave + earlier front vs batch18, anchors u\_th and DT\_ad announced pre-run, campaign energy on vessel enthalpy alone}
\begin{tutnarrative}
THE A5-T1 WITNESS: the adiabatic bed.

batch18's ergun-mode bed (15\% CO2/He, N = 25) with the energy balance DECLARED: `energyBalance adiabatic;` -- one temperature per cell, the isosteric heat staying IN the bed as a travelling thermal wave instead of being removed by an unnamed coolant. Design record: docs/design/fixed-bed-thermal-a5.md.

What the engine announces BEFORE integrating (equilibrium theory):

pure thermal-wave velocity u\_th = u c\_tot cpg / (eps c\_tot cpg + rho\_b cp\_s) and the time L/u\_th the heat needs to LEAVE the bed; the adiabatic LIMIT DT\_ad = q$\ast$(-dH\_ads)/(cp\_s + q$\ast$ cpg\_CO2 + eps c\_tot cpg/rho\_b) -- the all-heat-retained BOUND (the loaded solid carries the adsorbate's own Cp). Here u\_th < u\_sh: the heat LAGS the front, so the observed outlet-T rise sits well BELOW the bound (\textasciitilde{} +36 K vs the 126 K limit) while the WARM bed breaks through far earlier than batch18's isothermal front (t\_50 \textasciitilde{} 812 s vs 3042 s -- q$\ast$(T) falls via the van't Hoff b(T), the SAME datum the energy state carries; one dH\_ads, three readers).

The golden pins T\_out/T\_max, both anchors, the early crossings, and the campaign ENERGY closure -- claimed, adiabatically, on the vessel's internal energy against the boundary streams alone, and closed to round-off: the bed's energy rows are conservative finite volumes (each cell's internal energy U\_j, the faces carrying the enthalpy of the species fluxes; DEV.md 4c C258).

cp\_s enters as DECLARED case data (the kLDF pattern -- scope + source say where the number came from; curating a primary-cited cp onto the adsorbent asset record is a separate curation act). The adsorbed phase carries its gas's own Cp: a constant isosteric heat implies it (h\_ads = h\_gas(T) + dH\_ads). Feed enters at the declared T; each raffinate package carries the enthalpy the outlet face integrated over its window and leaves at its mixing-cup T, so the material ledger carries the thermal wave's enthalpy exactly.

Adiabatic means NO energy record: nothing is exchanged with the environment; the campaign energy balance closes on the vessel's internal energy (per cell: gas n(h - RT) and adsorbed n(h + dH\_ads) at T\_j, solid sensible cp\_s(T\_j-T0)) against the boundary streams alone.

\end{tutnarrative}
\begin{tutgolden}
\goldrow{kpi}{bed}{P\_final}{100050.550283}
\goldrow{kpi}{bed}{P\_first\_cell\_Pa}{100101.622717}
\goldrow{kpi}{bed}{P\_last\_cell\_Pa}{100002.081511}
\goldrow{kpi}{bed}{P\_mid\_cell\_Pa}{100051.159061}
\goldrow{kpi}{bed}{P\_out\_boundary\_Pa}{100000}
\goldrow{kpi}{bed}{Pe}{250}
\goldrow{kpi}{bed}{T\_final}{315.850484743}
\goldrow{kpi}{bed}{T\_max\_final\_K}{334.326452215}
\goldrow{kpi}{bed}{T\_out\_final\_K}{334.326452215}
\goldrow{kpi}{bed}{c\_out\_over\_cin\_final\_CO2}{0.823725183953}
\goldrow{kpi}{bed}{carrier\_fabricated\_mol}{2.4897417461e-11}
\goldrow{kpi}{bed}{dT\_adiabatic\_CO2\_K}{125.813871501}
\goldrow{kpi}{bed}{deltaP\_first\_to\_out\_Pa}{101.622716824}
\goldrow{kpi}{bed}{integral\_anchor\_CO2}{2839.09489826}
\goldmore{44}
\end{tutgolden}

\clearpage
\tutentry{batch21\_tsa\_hot\_purge}{tutorials/batch/adsorber/batch21\_tsa\_hot\_purge}
\begin{tutscheme}
\begin{tikzpicture}[>={Stealth[length=2mm]}, every node/.style={font=\scriptsize}]
  \node[draw, rounded corners=2pt, align=center, minimum width=1.9cm, minimum height=0.85cm, fill=black!4] (bed) at (0.0,-0.0) {\textbf{bed}\\[-1pt]{\tiny fixedBedAdsorber}};
\end{tikzpicture}
\end{tutscheme}
\tutwhat{T1.5 TSA half-cycle: adiabatic load to 3600 s, then 400 K clean He purge -- feed.T as the hot-purge control, the re-declared commitment ledgered OUT at 298 K / IN at 400 K, energy balance claimed across the swing}
\begin{tutnarrative}
THE T1.5 WITNESS: a real TSA half-cycle.

batch20's adiabatic bed run as LOAD $\to$ HOT CLEAN PURGE in one campaign:

0 .. 3600 s LOAD -- 15\% CO2/He at 298.15 K (batch20's feed; the warm front is through by \textasciitilde{}806 s) 3600 s SWITCH -- recipe, two setParameter events: feed.CO2 = 0 (clean He purge) feed.T = 400 K (the HOT purge) 3600 .. 12000 s PURGE -- the hot clean carrier sweeps the bed: q$\ast$(T) collapses via the van't Hoff b(T) and desorption runs FAR faster than batch19's cold concentration swing

This is the step the old refusal named: "thermal swing needs the non-isothermal bed energy balance". A5-T1 built that balance; feed.T is now the hot-purge CONTROL (in thermal mode ONLY -- the isothermal bed still refuses it, by name).

The LEDGER story is the point. A feed at a new temperature re-declares the WHOLE remaining commitment -- at the pinned feed pressure the total concentration drops to P/(R$\ast$400) (composition preserved) and even unchanged matter changes enthalpy -- so the switch retires the old commitment (one OUT feedAmendment package, all species, at 298.15 K) and declares the new one (one IN package at 400 K): the enthalpy repricing rides in the packages' own T, and the campaign energy balance stays CLAIMED across the swing, boundary streams alone.

Breakthrough crossings/integral freeze at the switch (they describe the load step); the desorption story lives in c\_out/qbar/T\_out. The frozen pre-run anchors (u\_th, DT\_ad bound) keep their initialise values -- a later switch never rewrites an announced claim.

\end{tutnarrative}
\begin{tutgolden}
\goldrow{kpi}{bed}{P\_final}{100030.237558}
\goldrow{kpi}{bed}{P\_first\_cell\_Pa}{100070.37249}
\goldrow{kpi}{bed}{P\_last\_cell\_Pa}{100001.142011}
\goldrow{kpi}{bed}{P\_mid\_cell\_Pa}{100028.478818}
\goldrow{kpi}{bed}{P\_out\_boundary\_Pa}{100000}
\goldrow{kpi}{bed}{Pe}{250}
\goldrow{kpi}{bed}{T\_final}{349.866000841}
\goldrow{kpi}{bed}{T\_max\_final\_K}{399.999997624}
\goldrow{kpi}{bed}{T\_out\_final\_K}{306.831902423}
\goldrow{kpi}{bed}{carrier\_fabricated\_mol}{-1.13686837722e-12}
\goldrow{kpi}{bed}{dT\_adiabatic\_CO2\_K}{125.813871501}
\goldrow{kpi}{bed}{deltaP\_first\_to\_out\_Pa}{70.3724897253}
\goldrow{kpi}{bed}{integral\_anchor\_CO2}{1315.23023598}
\goldrow{kpi}{bed}{mass\_closure\_CO2}{4.72969836841e-15}
\goldmore{53}
\end{tutgolden}

\clearpage
\tutentry{batch22\_wall\_cooled}{tutorials/batch/adsorber/batch22\_wall\_cooled}
\begin{tutscheme}
\begin{tikzpicture}[>={Stealth[length=2mm]}, every node/.style={font=\scriptsize}]
  \node[draw, rounded corners=2pt, align=center, minimum width=1.9cm, minimum height=0.85cm, fill=black!4] (bed) at (0.0,-0.0) {\textbf{bed}\\[-1pt]{\tiny fixedBedAdsorber}};
\end{tikzpicture}
\end{tutscheme}
\tutwhat{T2 wall-cooled CO2/He breakthrough (h=5 W/m2K, T\_wall=298.15 K): the bed BETWEEN its limits -- t\_50 bracketed by the adiabatic 806 s and isothermal 3036 s, Q\_wall as a state row of the same ODE, wallHeat record read from the state}
\begin{tutnarrative}
THE T2 WITNESS: the bed between its two limits.

batch20's adiabatic bed with `energyBalance wallCooled;` -- the energy row gains -h a\_w (T\_j - T\_wall) and the bed sits BETWEEN the limits it must contain:

h $\to$ 0 batch20 (adiabatic: t\_50 = 806 s, T\_max \textasciitilde{} 336 K) h $\to$ infinity batch18 (isothermal: t\_50 = 3036 s, T = T\_wall)

so the golden pins the BRACKETING:

t\_50(adiabatic) < t\_50(wallCooled) < t\_50(isothermal) T\_wall < T\_max(wallCooled) < T\_max(adiabatic)

at h = 5 W/(m2 K), T\_wall = 298.15 K, dBed 0.113 m DECLARED (never derived from the cross-section by a silent circle assumption); the engine announces a\_w = 4/dBed and the wall NTU before integrating.

The LEDGER is the design point (docs/design/fixed-bed-thermal-a5.md section 6): the removed heat integrates as ONE MORE STATE ROW of the same ODE -- the M\_in/M\_out ledger-row pattern -- so the `wallHeat` energy record READS a state (E = -Q\_wall, T\_service = T\_wall) and the campaign closure dH\_vessels = Q\_ledger + H\_external holds with every term a state function. Guards: `adiabatic` + a wallHeatTransfer block refuses (contradiction), `wallCooled` without one refuses, h <= 0 refuses (adiabatic is the way to SAY zero) -- gate check\_thermal\_bed.

\end{tutnarrative}
\begin{tutgolden}
\goldrow{kpi}{bed}{E\_wallHeat\_seg0\_kJ}{-185.062633189}
\goldrow{kpi}{bed}{E\_wallHeat\_total\_kJ}{-185.062633189}
\goldrow{kpi}{bed}{P\_final}{100048.490853}
\goldrow{kpi}{bed}{P\_first\_cell\_Pa}{100096.32317}
\goldrow{kpi}{bed}{P\_last\_cell\_Pa}{100001.949113}
\goldrow{kpi}{bed}{P\_mid\_cell\_Pa}{100048.72665}
\goldrow{kpi}{bed}{P\_out\_boundary\_Pa}{100000}
\goldrow{kpi}{bed}{Pe}{250}
\goldrow{kpi}{bed}{Q\_wall\_removed\_kJ}{185.062633189}
\goldrow{kpi}{bed}{T\_final}{306.216508093}
\goldrow{kpi}{bed}{T\_max\_final\_K}{312.429267214}
\goldrow{kpi}{bed}{T\_out\_final\_K}{312.429267214}
\goldrow{kpi}{bed}{c\_out\_over\_cin\_final\_CO2}{0.890377021134}
\goldrow{kpi}{bed}{carrier\_fabricated\_mol}{2.4698465495e-11}
\goldmore{52}
\end{tutgolden}

\clearpage
\tutentry{batch23\_tsa\_cycles}{tutorials/batch/adsorber/batch23\_tsa\_cycles}
\begin{tutscheme}
\begin{tikzpicture}[>={Stealth[length=2mm]}, every node/.style={font=\scriptsize}]
  \node[draw, rounded corners=2pt, align=center, minimum width=1.9cm, minimum height=0.85cm, fill=black!4] (bed) at (0.0,-0.0) {\textbf{bed}\\[-1pt]{\tiny fixedBedAdsorber}};
\end{tikzpicture}
\end{tutscheme}
\tutwhat{Toward A6: three hand-declared TSA cycles (cold load / 400 K purge) in one campaign -- eleven ledgered feed amendments, balances closed, final KPIs at the cycle boundary; the CSS verdict machinery stays a named proposal}
\begin{tutnarrative}
toward A6: three FULL TSA cycles, hand-declared.

batch21's half-cycle repeated three times in ONE campaign, with the recipe grammar that already exists -- nothing new in the engine:

cycle k (k = 0,1,2), each 7200 s: t = 7200k feed.T = 298.15 K (cold again) [k>0] feed.CO2 = 0.15 (feed returns) [k>0] t = 7200k + 3600 feed.CO2 = 0 (clean purge) feed.T = 400 K (hot purge)

ending EXACTLY at the third cycle boundary (21600 s), so the FINAL KPIs are end-of-cycle values. Every switch re-declares the feed commitment through the T1.5 ledger machinery (the T switches retire and redeclare the WHOLE commitment; the composition switches amend by delta) -- FIFTEEN feedAmendment records (three per switching instant: the composition delta + the T retire/declare pair), and the campaign mass/element/energy balances still close on records, never jumps.

The A6 QUESTION this case makes visible (and does not yet verdict): does the bed state at successive cycle ends approach a fixed point -- the cyclic steady state? MEASURED (this case's own run): qbar\_CO2 at the three cycle ends is 0.3151 $\to$ 0.2415 $\to$ 0.1748 mol/kg -- a slow geometric approach (ratio \textasciitilde{}0.9; each purge leaves residual heat, the next load starts warmer and holds less), NOT yet converged in three cycles. That unfinished convergence IS the pedagogical point: eyes on a trajectory can see it, and only the A6 verdict machinery could CLAIM it. The golden pins the third-cycle end. A first-class CSS verdict (declared cycles, per- cycle state snapshots, an invariance tolerance the case declares) is an ARCHITECTURE feature -- proposed in docs/design/fixed-bed-thermal-a5.md section 7, awaiting Vitor's decision; this case is deliberately its cheap, grammar-free witness.

TIME INTEGRATION: `solver \{ method rodas4; \}` selects the order-4 stiff integrator (Hairer \& Wanner's RODAS4) instead of the default 2(3) Rosenbrock pair. Same rtol, same Jacobian, same golden: 2.4x fewer steps and about half the wall clock on this case (it is the witness for that option; the other adsorber cases keep the default).

\end{tutnarrative}
\begin{tutgolden}
\goldrow{kpi}{bed}{P\_final}{100029.977681}
\goldrow{kpi}{bed}{P\_first\_cell\_Pa}{100063.941541}
\goldrow{kpi}{bed}{P\_last\_cell\_Pa}{100001.18892}
\goldrow{kpi}{bed}{P\_mid\_cell\_Pa}{100030.107889}
\goldrow{kpi}{bed}{P\_out\_boundary\_Pa}{100000}
\goldrow{kpi}{bed}{Pe}{250}
\goldrow{kpi}{bed}{T\_final}{328.511667754}
\goldrow{kpi}{bed}{T\_max\_final\_K}{399.868834476}
\goldrow{kpi}{bed}{T\_out\_final\_K}{314.928123967}
\goldrow{kpi}{bed}{carrier\_fabricated\_mol}{-5.09885467181e-11}
\goldrow{kpi}{bed}{dT\_adiabatic\_CO2\_K}{125.813871501}
\goldrow{kpi}{bed}{deltaP\_first\_to\_out\_Pa}{63.9415407171}
\goldrow{kpi}{bed}{integral\_anchor\_CO2}{1315.23082099}
\goldrow{kpi}{bed}{mass\_closure\_CO2}{6.95863667995e-14}
\goldmore{58}
\end{tutgolden}

\clearpage
\tutentry{batch24\_blowdown\_inert}{tutorials/batch/adsorber/batch24\_blowdown\_inert}
\begin{tutscheme}
\begin{tikzpicture}[>={Stealth[length=2mm]}, every node/.style={font=\scriptsize}]
  \node[draw, rounded corners=2pt, align=center, minimum width=1.9cm, minimum height=0.85cm, fill=black!4] (bed) at (0.0,-0.0) {\textbf{bed}\\[-1pt]{\tiny fixedBedAdsorber}};
\end{tikzpicture}
\end{tutscheme}
\tutwhat{P-swing witness A: inert isothermal blowdown, n/n0 = P/P0}
\begin{tutnarrative}
P-SWING WITNESS A: the inventory ratio.

The MATERIAL half of a blowdown, isolated from every thermal effect so it cannot pass by coincidence (design note fixed-bed-thermal-a5.md, section 8/9 amended battery, second-opinion review 2026-08-03):

ideal gas, INERT bed (kLDF = 0: the solid is frozen, pure transport), pure He; ISOTHERMAL ergun mode -- T is a DECLARED constant of this teaching slice, so the energy question is out of frame BY DECLARATION (the thermal half is witness B, batch25); at t = 200 s the recipe closes the feed valve (P2, `setParameter u = 0` -- the whole remaining feed commitment retires as a ledgered feedAmendment record) and drops the outlet boundary (P1, `setParameter P\_out` 1.0 $\to$ 0.5 bar); the bed drains through the SAME Ergun face machinery every run already uses; the vented gas rides the telescopic M\_out ledger rows.

THE ANCHORS (ideal-gas exactness, hand-computable):

n\_end / n\_0 = P\_end / P\_0 = 0.5 n(P) = eps A L P / (R T) = 0.0807 mol at 1 bar, 0.0403 at 0.5 bar vented = inventory drop (the mass-closure KPI is the proof)

The gate (check\_pswing) recomputes both analytically and pins the `gasInventory\_mol` KPI + the closure against them.

\end{tutnarrative}
\begin{tutgolden}
\goldrow{kpi}{bed}{E\_adsorption\_seg0\_kJ}{0}
\goldrow{kpi}{bed}{E\_adsorption\_total\_kJ}{0}
\goldrow{kpi}{bed}{P\_final}{49999.9998988}
\goldrow{kpi}{bed}{P\_first\_cell\_Pa}{49999.9998996}
\goldrow{kpi}{bed}{P\_last\_cell\_Pa}{49999.9998972}
\goldrow{kpi}{bed}{P\_mid\_cell\_Pa}{49999.999899}
\goldrow{kpi}{bed}{P\_out\_boundary\_Pa}{50000}
\goldrow{kpi}{bed}{Pe}{0}
\goldrow{kpi}{bed}{T\_final}{298.15}
\goldrow{kpi}{bed}{carrier\_fabricated\_mol}{-8.881784197e-16}
\goldrow{kpi}{bed}{deltaP\_first\_to\_out\_Pa}{1.45519152284e-11}
\goldrow{kpi}{bed}{gasInventory\_mol}{0.0403395454642}
\goldrow{kpi}{bed}{n\_CO2\_final}{0}
\goldrow{kpi}{bed}{n\_He\_final}{4.03395454642e-05}
\goldmore{48}
\end{tutgolden}

\clearpage
\tutentry{batch24\_tsa\_until\_css}{tutorials/batch/adsorber/batch24\_tsa\_until\_css}
\begin{tutscheme}
\begin{tikzpicture}[>={Stealth[length=2mm]}, every node/.style={font=\scriptsize}]
  \node[draw, rounded corners=2pt, align=center, minimum width=1.9cm, minimum height=0.85cm, fill=black!4] (bed) at (0.0,-0.0) {\textbf{bed}\\[-1pt]{\tiny fixedBedAdsorber}};
\end{tikzpicture}
\end{tutscheme}
\tutwhat{Toward A6: three hand-declared TSA cycles (cold load / 400 K purge) in one campaign -- eleven ledgered feed amendments, balances closed, final KPIs at the cycle boundary; the CSS verdict machinery stays a named proposal}
\begin{tutnarrative}
A6 item 4: the cyclic steady state as the STOPPING CONDITION, not a count.

batch23 beside this one declares `repeat 3;` and reports honestly that three cycles did not settle -- it makes the question visible and does not answer it. This case asks what an engineer actually wants to ask of a TSA bed: not "run three cycles" but "cycle until it repeats, and tell me how many that took".

cycle \{ period 7200; repeat untilCSS; cssTolerance 3e-2; maxCycles 6; ... \}

The declared cap is SIX and the run uses TWO, announcing why it stopped.

TWO KEYS ARE MANDATORY under `untilCSS` and the engine refuses without either: the tolerance, because whether a bed has converged is a modelling judgement the engine must never invent; and the cap, because a bed that never settles has to stop somewhere and the run must be able to say it stopped for want of cycles rather than because it arrived. Reaching the cap is NOT an error -- the trajectory is a physically valid campaign and every other KPI in it is real -- but it means the case's own question went unanswered, so it is published (`css\_stopped\_at\_cap`) and announced.

WHAT BUILDING IT FOUND, and it outlives the feature: the CSS norm and the loading are NOT the same quantity. See the measured numbers beside `cssTolerance` below. A tolerance chosen without looking can certify a bed whose loading is still visibly marching, because the norm over the whole packed state is dominated by its largest components. The verdict is only as meaningful as the number the case declares -- which is the argument for the engine refusing to supply one.

\end{tutnarrative}
\begin{tutgolden}
\goldrow{kpi}{bed}{P\_final}{100029.449382}
\goldrow{kpi}{bed}{P\_first\_cell\_Pa}{100062.639643}
\goldrow{kpi}{bed}{P\_last\_cell\_Pa}{100001.176121}
\goldrow{kpi}{bed}{P\_mid\_cell\_Pa}{100029.404542}
\goldrow{kpi}{bed}{P\_out\_boundary\_Pa}{100000}
\goldrow{kpi}{bed}{Pe}{250}
\goldrow{kpi}{bed}{T\_final}{324.097107435}
\goldrow{kpi}{bed}{T\_max\_final\_K}{399.868835232}
\goldrow{kpi}{bed}{T\_min\_final\_K}{297.290399767}
\goldrow{kpi}{bed}{T\_out\_final\_K}{310.560011256}
\goldrow{kpi}{bed}{carrier\_fabricated\_mol}{1.42108547152e-11}
\goldrow{kpi}{bed}{dT\_adiabatic\_CO2\_K}{125.813871501}
\goldrow{kpi}{bed}{deltaP\_first\_to\_out\_Pa}{62.6396434308}
\goldrow{kpi}{bed}{gasInventory\_mol}{0.0748325309122}
\goldmore{60}
\end{tutgolden}

\clearpage
\tutentry{batch25\_blowdown\_thermal}{tutorials/batch/adsorber/batch25\_blowdown\_thermal}
\begin{tutscheme}
\begin{tikzpicture}[>={Stealth[length=2mm]}, every node/.style={font=\scriptsize}]
  \node[draw, rounded corners=2pt, align=center, minimum width=1.9cm, minimum height=0.85cm, fill=black!4] (bed) at (0.0,-0.0) {\textbf{bed}\\[-1pt]{\tiny fixedBedAdsorber}};
\end{tikzpicture}
\end{tutscheme}
\tutwhat{P-swing witness B: adiabatic blowdown COOLS (expansion work)}
\begin{tutnarrative}
P-SWING WITNESS B: the cooling.

The THERMAL half of a blowdown, isolated from adsorption so the only heat effect in frame is the expansion work itself (design note fixed-bed-thermal-a5.md, section 8/9 amended battery):

same inert pure-He ergun bed as witness A (batch24), but `energyBalance adiabatic;` -- each cell's internal energy U\_j = eps c (h - RT) + rho\_b cp\_s (T - T0) is balanced in conservative finite-volume form (design note section 10); at t = 200 s the recipe closes the feed valve (P2) and halves the outlet boundary (P1) -- the bed drains, and the gas that stays behind did the work of pushing the vented gas out: read through the chain rule, the eps R T dc\_tot/dt term pulls T DOWN; the solid heat capacity is a declared SYNTHETIC LOW value (9.2 J/(kg K), scope says so): the real \textasciitilde{}920 of a zeolite pellet would thermally bury the gas-phase cooling under the solid's mass (dT \textasciitilde{} -0.03 K); at the teaching value the drop is visible (\textasciitilde{} -3 K class). Same honest-teaching-value discipline as the hxFluid of the pinch witness -- the physics on display is the GAS's, and the header says which knob made it visible.

THE ANCHOR: T\_min\_final\_K sits BELOW the declared 298.15 K by a visible margin -- a depressurising cell COOLS, which is exactly what the pre-widening cp-form could not show (the sabotage of the gate restores that form and this witness fails). The material anchors stay witness A's business; this case pins the sign and visibility of the thermal effect. AND ITS FIRST LAW: the campaign energy balance is claimed and closes to round-off (energy\_closure\_rel \textasciitilde{} 4e-10 of the 2.3 J the blowdown moves, against \textasciitilde{}200 J of gross hold-up energy), because the face enthalpy fluxes are shared by the two cells of each face and the vented gas leaves carrying exactly what crossed the outlet face. Until 2026-11-09 the bed carried a temperature form that missed this first law by 7.9 \% of the energy moved, at every step size, and the engine refused the claim by name (DEV.md 4c C258).

\end{tutnarrative}
\begin{tutgolden}
\goldrow{kpi}{bed}{P\_final}{49999.9998351}
\goldrow{kpi}{bed}{P\_first\_cell\_Pa}{49999.9997932}
\goldrow{kpi}{bed}{P\_last\_cell\_Pa}{50000}
\goldrow{kpi}{bed}{P\_mid\_cell\_Pa}{49999.9998123}
\goldrow{kpi}{bed}{P\_out\_boundary\_Pa}{50000}
\goldrow{kpi}{bed}{Pe}{0}
\goldrow{kpi}{bed}{T\_final}{294.92367086}
\goldrow{kpi}{bed}{T\_max\_final\_K}{295.066441477}
\goldrow{kpi}{bed}{T\_min\_final\_K}{294.887445096}
\goldrow{kpi}{bed}{T\_out\_final\_K}{295.066441477}
\goldrow{kpi}{bed}{carrier\_fabricated\_mol}{4.17443857259e-14}
\goldrow{kpi}{bed}{deltaP\_first\_to\_out\_Pa}{-0.000205334486964}
\goldrow{kpi}{bed}{gasInventory\_mol}{0.0407808424401}
\goldrow{kpi}{bed}{n\_CO2\_final}{0}
\goldmore{41}
\end{tutgolden}

\clearpage
\tutentry{batch29\_pulse\_two\_schemes}{tutorials/batch/adsorber/batch29\_pulse\_two\_schemes}
\begin{tutscheme}
\begin{tikzpicture}[>={Stealth[length=2mm]}, every node/.style={font=\scriptsize}]
  \node[draw, rounded corners=2pt, align=center, minimum width=1.9cm, minimum height=0.85cm, fill=black!4] (bedUpwind) at (0.0,-0.0) {\textbf{bedUpwind}\\[-1pt]{\tiny fixedBedAdsorber}};
  \node[draw, rounded corners=2pt, align=center, minimum width=1.9cm, minimum height=0.85cm, fill=black!4] (bedVanLeer) at (0.0,-1.7) {\textbf{bedVanLeer}\\[-1pt]{\tiny fixedBedAdsorber}};
\end{tikzpicture}
\end{tutscheme}
\tutwhat{A linear pulse through two identical fixed beds, upwind and van Leer, against the closed-form mean and variance: the grid widens the upwind peak by v dz/2, never moves its centre}
\begin{tutnarrative}
the same bed, two schemes, against the closed form.

NUMERICAL DISPERSION, for a student. The bed is cut into cells, and the convective flux through each face needs ONE concentration at the face. The first-order `upwind` scheme takes the upstream cell's value; a Taylor expansion shows that this differs from the true face value by -(dz/2) dc/dz, so the scheme quietly adds a Fickian flux: it solves the model with D\_L + v dz/2 instead of D\_L (v = u/eps). Here v dz/2 is 3.1e-4 m2/s at 100 cells against a declared D\_L = 2.5e-4: the grid more than doubles the dispersion. The `vanLeer` scheme reconstructs a limited linear profile in each cell, which removes that term (second order), and its limiter keeps the peak free of over- and undershoots.

THE EXPERIMENT. Two identical, independent beds of a synthetic LINEAR sorbent (Henry isotherm, case-local records under constant/), fed 1 \% CO2 in He for t\_inj = 2 s and then pure He (a rectangular pulse, the recipe below). Each bed publishes the moments of its outlet concentration (operation.schemeDiagnostics true;). The model is linear, so its moments have a closed form (derived in the Theory Guide, fixed-bed section; symbols: eps void fraction, u superficial velocity, v = u/eps, D\_L = Dax/eps, K = H R T the linear isotherm in m3/kg, rho\_b bulk density, k the LDF constant):

t0 = eps L/u = 4 s k' = rho\_b K/eps = 12.39479 t\_R = t0 (1 + k') = 53.57914 s Pe = v L/D\_L = u L/Dax = 250 mean = t\_R + t\_inj/2 = 54.57914 s variance = t\_R\textasciicircum{}2 [2/Pe - 2(1 - exp(-Pe))/Pe\textasciicircum{}2] (Danckwerts, closed) + 2 k' t0/k (LDF) + t\_inj\textasciicircum{}2/12 (the pulse) = 22.87393 + 9.91583 + 0.33333 = 33.12309 s2

bedVanLeer (100 cells) must reproduce THAT variance; bedUpwind (100 cells) reproduces it only with D\_L replaced by D\_L + v dz/2 (Pe\_eff = 111.1): 61.45714 s2. Both means are t\_R + t\_inj/2: the grid moves the width, never the centre. The `anchor` rows in `expected` are these closed-form values (DERIVED anchors: identities of the linear model, not measurements); the gate check\_fixed\_bed\_scheme repeats them at 200 cells and states the tolerance it measured.

\end{tutnarrative}
\begin{tutgolden}
\goldrow{kpi}{bedUpwind}{D\_axial\_interstitial}{0.00025}
\goldrow{kpi}{bedUpwind}{D\_numerical}{0.0003125}
\goldrow{kpi}{bedUpwind}{E\_adsorption\_seg0\_kJ}{-2.6519254011e-28}
\goldrow{kpi}{bedUpwind}{E\_adsorption\_total\_kJ}{-2.6519254011e-28}
\goldrow{kpi}{bedUpwind}{P\_final}{100000}
\goldrow{kpi}{bedUpwind}{Pe}{250}
\goldrow{kpi}{bedUpwind}{T\_final}{298.15}
\goldrow{kpi}{bedUpwind}{c\_max\_over\_feed\_CO2}{0.906065328263}
\goldrow{kpi}{bedUpwind}{c\_min\_CO2}{0}
\goldrow{kpi}{bedUpwind}{carrier\_fabricated\_mol}{1.32596270055e-29}
\goldrow{kpi}{bedUpwind}{carrier\_makeup\_rel}{3.28700455746e-30}
\goldrow{kpi}{bedUpwind}{integral\_anchor\_CO2}{2}
\goldrow{kpi}{bedUpwind}{mass\_closure\_CO2}{1.30352993952e-14}
\goldrow{kpi}{bedUpwind}{nCells\_numerical\_ratio\_0p1}{1251}
\goldmore{103}
\end{tutgolden}

\clearpage
\tutentry{batch30\_bed\_correlations}{tutorials/batch/adsorber/batch30\_bed\_correlations}
\begin{tutscheme}
\begin{tikzpicture}[>={Stealth[length=2mm]}, every node/.style={font=\scriptsize}]
  \node[draw, rounded corners=2pt, align=center, minimum width=1.9cm, minimum height=0.85cm, fill=black!4] (bed) at (0.0,-0.0) {\textbf{bed}\\[-1pt]{\tiny fixedBedAdsorber}};
\end{tikzpicture}
\end{tutscheme}
\tutwhat{Where do Dax, k\_f and k\_LDF come from? -- batch13's CO2/N2 bed with the declared Dax and k\_LDF beside the literature correlations (Wakao-Funazkri, Edwards-Richardson, Chung-Wen, Gunn; Wilson-Geankoplis; Ruthven series with Glueckauf's 15), each estimate with its window and page}
\begin{tutnarrative}
where do Dax, k\_f and k\_LDF come from?

batch13's bed (15 \% CO2 on 13X beads, 0.5 m, u = 0.05 m/s, 1 bar, 298 K) with the carrier changed to N2 (the catalogue's helium has no diffusion volume) and the three transport numbers put beside what the literature correlations say. The case DECLARES Dax and k\_LDF, as every adsorber case does, and a declared value always wins -- so the front below is the declared one. `transportCorrelations` makes the run ALSO print and publish (KPIs corr\_$\ast$) every correlation's estimate, each with its validity window and the page its coefficients were read at:

axial dispersion Wakao \& Funazkri 1978 eq (2); Edwards \& Richardson 1968 eq (17) (gas); Chung \& Wen 1968 eqs (15)-(16) (liquids); Gunn 1987 eq (1) film k\_f Wakao \& Funazkri 1978 eq (12); Wilson \& Geankoplis 1966 eqs (8)-(9) (liquids) k\_LDF (series) Ruthven 1984 eq (8.41): film + macropore + micropore, the internal terms with Glueckauf's 15 (1955 eqn 3a)

BASIS. Dax is stored on the BED (superficial) basis, Dax = eps D\_L; the correlations are written for the interstitial D\_L and are converted (the run prints both). A textbook D\_L typed into Dax overstates the dispersion by 1/eps.

WHAT THE RUN SHOWS (feed state: rho = 1.227 kg/m3, mu = 1.72e-5 Pa s, Re = 7.13, Sc = 0.854, Dm = 1.64e-5 m2/s from Fuller): Dax (bed basis, m2/s): Edwards \& Richardson 3.5e-5 and Gunn 3.2e-5 agree within 10 \% and sit BELOW the declared 1e-4 (batch13's Pe = 250 is a conservative teaching value, not a prediction); Wakao \& Funazkri 3.8e-4 is HIGHER because its stagnant term 20 Dm (a particle whose surface controls the transfer) dominates at this Re; Chung \& Wen 1.75e-4 is flagged OUTSIDE its phase (a liquid correlation) and the run says so -- the window is the lesson; the film coefficient is k\_f = 0.044 m/s (Wakao \& Funazkri, Sh = 5.39; Wilson \& Geankoplis answers 0.041 but is flagged: liquids, Sc 950 up); the series k\_LDF = 0.0113 1/s, against the declared 0.05: the shares of 1/k say WHERE the resistance is -- film 4.3 \%, macropore 20.0 \%, micropore 75.7 \% -- and the declared value, 4.4 times faster, is published beside the estimate, not used (the front below is the declared one: t\_b(5 \%) = 2934.07 s, as batch13). The estimates need particle data the catalogue cannot hold (porosity, tortuosity, micropore diffusivity of THIS sample), declared below as sample data with their scope.

Tolerances: the golden pins the estimates to 1e-4 relative; the anchors of every correlation are checked by check\_bed\_correlations.py.

\end{tutnarrative}
\begin{tutgolden}
\goldrow{kpi}{bed}{D\_axial\_interstitial}{0.00025}
\goldrow{kpi}{bed}{D\_numerical}{0}
\goldrow{kpi}{bed}{E\_adsorption\_seg0\_kJ}{-413.617021276}
\goldrow{kpi}{bed}{E\_adsorption\_total\_kJ}{-413.617021276}
\goldrow{kpi}{bed}{P\_final}{100000}
\goldrow{kpi}{bed}{Pe}{250}
\goldrow{kpi}{bed}{T\_final}{298}
\goldrow{kpi}{bed}{c\_out\_over\_cin\_final\_CO2}{0.999999999943}
\goldrow{kpi}{bed}{carrier\_fabricated\_mol}{9.19148936169}
\goldrow{kpi}{bed}{carrier\_makeup\_rel}{0.108446875338}
\goldrow{kpi}{bed}{corr\_Dax\_ChungWen}{0.000175250886463}
\goldrow{kpi}{bed}{corr\_Dax\_EdwardsRichardson}{3.53384489604e-05}
\goldrow{kpi}{bed}{corr\_Dax\_Gunn}{3.21840656335e-05}
\goldrow{kpi}{bed}{corr\_Dax\_WakaoFunazkri}{0.000378310977433}
\goldmore{78}
\end{tutgolden}

\subsection{combustion}
\grouptheory{combustion}
\clearpage
\tutentry{ignition01\_h2o2}{tutorials/batch/combustion/ignition01\_h2o2}
\begin{tutscheme}
\begin{tikzpicture}[>={Stealth[length=2mm]}, every node/.style={font=\scriptsize}]
  \node[draw, rounded corners=2pt, align=center, minimum width=1.9cm, minimum height=0.85cm, fill=black!4] (bomb) at (0.0,-0.0) {\textbf{bomb}\\[-1pt]{\tiny batchReactor}};
\end{tikzpicture}
\end{tutscheme}
\tutwhat{Adiabatic constant-V H2/O2 ignition at 1100 K (O Conaire 2004 H2/O2 kinetics, NASA-9 thermo); Rosenbrock23 stiff solver; reports ignition delay}
\begin{tutnarrative}
closed rigid vessel, stoichiometric H2/O2 in N2 at 1100 K.

H2 : O2 : N2 = 2 : 1 : 4 (stoichiometric H2 + 1/2 O2 $\to$ H2O, 57\% N2)

mode adiabatic + energy gasConstantV: the first law dU = 0 in a fixed volume, with the heat release coming from the species FORMATION enthalpies (curated thermo) -- no per-reaction $\Delta$H is supplied. Temperature, and the radicals H/O/OH/HO2, are integrated by the stiff solver until branching runs away.

The 21 reactions are read from the kinetics library (`reactionSet OConaire2004;` in constant/reactions): one record per reaction, each with its journal page.

solver \{ integrator Rosenbrock23; \} -- the L-stable stiff method. Swap to RK4 (see the ignition01\_h2o2\_rk4 sibling) and the same step blows up.

\end{tutnarrative}
\tutreadme{A closed rigid vessel holds a stoichiometric H$_2$/O$_2$ charge diluted in N$_2$ at 1100 K and 5 bar. A trace H-radical seed starts chain branching; after a short induction the radical pool explodes and the mixture ignites --- a sharp temperature jump from 1100 K to \textasciitilde{}3078 K.}
\tutreadme{and plot trajectory.csv (T and the radicals H/O/OH vs time). You will see:}
\tutreadme{The ignition delay $\tau$\_ign (the time of maximum dT/dt, printed in the run's epilogue on the 1 $\mu$s write grid) is $\approx$ 8 $\mu$s here. The trace H seed is what makes it that short: the same mechanism without a seed, at 2 atm, gives the 43 $\mu$s of the shock-tube pin in ignition02\_h2air\_slack.}
\tutreadme{The radical lifetimes are far shorter than the induction --- the solver prints a stiffness ratio of order 10?. That is stiffness. An explicit method (RK4) is stability-limited to a step of order the $\ast$fastest$\ast$ scale, so it cannot take the 1 $\mu$s steps this case uses. The solver here is Rosenbrock23, an L-stable linearly-implicit method that steps by $\ast$accuracy$\ast$, not stability; at verbosity 3 it prints its stiffness verdict and step counts (the no-silent- crutch credo).}
\tutreadme{See it fail on purpose --- two siblings run the same problem at the same step with other integrators:}
\tutreadme{(Both ship a .expect-nonconvergence marker --- their whole point is to fail.) $\ast$See, then decide$\ast$, applied to numerics.}
\begin{tutgolden}
\goldrow{kpi}{bomb}{P\_final}{1121202.82397}
\goldrow{kpi}{bomb}{T\_final}{3037.00819983}
\goldrow{kpi}{bomb}{n\_H2O2\_final}{8.00367416065e-11}
\goldrow{kpi}{bomb}{n\_H2\_final}{1.62963087713e-06}
\goldrow{kpi}{bomb}{n\_HO2\_final}{7.19758331359e-10}
\goldrow{kpi}{bomb}{n\_H\_final}{4.49541510152e-07}
\goldrow{kpi}{bomb}{n\_N2\_final}{2.857145e-05}
\goldrow{kpi}{bomb}{n\_O2\_final}{5.33369140811e-07}
\goldrow{kpi}{bomb}{n\_OH\_final}{1.20519239058e-06}
\goldrow{kpi}{bomb}{n\_O\_final}{1.86407481061e-07}
\goldrow{kpi}{bomb}{n\_water\_final}{1.18257622566e-05}
\goldrow{kpi}{bomb}{t\_end}{0.0002}
\goldrow{kpi}{bomb}{totalMoles\_final}{4.44021534514e-05}
\goldrow{kpi}{campaign}{element\_H\_closure\_rel}{1.51815023592e-14}
\goldmore{17}
\end{tutgolden}

\clearpage
\tutentry{ignition01\_h2o2\_eulersi}{tutorials/batch/combustion/ignition01\_h2o2\_eulersi}
\begin{tutscheme}
\begin{tikzpicture}[>={Stealth[length=2mm]}, every node/.style={font=\scriptsize}]
  \node[draw, rounded corners=2pt, align=center, minimum width=1.9cm, minimum height=0.85cm, fill=black!4] (bomb) at (0.0,-0.0) {\textbf{bomb}\\[-1pt]{\tiny batchReactor}};
\end{tikzpicture}
\end{tutscheme}
\tutwhat{WITNESS + EXPECTED FAILURE: semi-implicit Euler on the H2/O2 ignition -- A-stable but first-order, goes unphysical at this step and the engine refuses by name (see .expect-nonconvergence).}
\begin{tutnarrative}
WITNESS for the EulerSI integrator (semi-implicit Euler) -- an EXPECTED FAILURE: see .expect-nonconvergence for why first order at this step goes unphysical and how the engine now refuses instead of exiting 0. Registered model no case selected until 2026-08-22 (invariant I8).

\end{tutnarrative}
\tutreadme{The third integrator of the stiffness lesson. The same adiabatic constant-volume H$_2$/O$_2$ ignition as [ignition01\_h2o2](../ignition01\_h2o2/README.md) (Rosenbrock23, converges) and [ignition01\_h2o2\_rk4](../ignition01\_h2o2\_rk4/README.md) (explicit RK4, blows up), here with semi-implicit Euler (integrator EulerSI;) at the same deltaT = 1e-6 s.}
\tutreadme{Semi-implicit Euler is A-stable, so it does not oscillate to infinity the way RK4 does --- the first steps are physical (T 1100 $\to$ 1102 $\to$ 1108 $\to$ 1139 K) --- but it is FIRST ORDER, and once the ignition transient starts one linearised step is wrong enough to drive the state unphysical. The integrator reports that it gave up, and the reactor refuses by name instead of exiting 0:}
\tutreadme{.expect-nonconvergence records why that refusal exists (until 2026-08-22 this exact run reached a negative temperature and exited 0). L-stable second order converges, explicit fourth order blows up loudly, A-stable first order fails honestly: $\ast$see, then decide$\ast$, applied to numerics.}
\tutreadme{Same as the sibling: the ? Conaire et al. (2004) H$_2$/O$_2$ mechanism (constant/reactions, shared with ignition02\_h2air\_slack) and NASA-9 heat capacities from NASA/TP-2002-211556 for H$_2$, O$_2$, N$_2$ and water. Until 2026-10-25 it ran a GRI-Mech 3.0 subset, replaced because GRI-Mech grants no redistribution (DEV.md ?4c C164); the failure it teaches did not change.}

\clearpage
\tutentry{ignition01\_h2o2\_jac}{tutorials/batch/combustion/ignition01\_h2o2\_jac}
\begin{tutscheme}
\begin{tikzpicture}[>={Stealth[length=2mm]}, every node/.style={font=\scriptsize}]
  \node[draw, rounded corners=2pt, align=center, minimum width=1.9cm, minimum height=0.85cm, fill=black!4] (bomb) at (0.0,-0.0) {\textbf{bomb}\\[-1pt]{\tiny batchReactor}};
\end{tikzpicture}
\end{tutscheme}
\tutwhat{Adiabatic constant-V H2/O2 ignition at 1100 K on the ANALYTIC kinetics Jacobian, audited against the finite difference (O Conaire 2004 H2/O2 kinetics, NASA-9 thermo); Rosenbrock23 stiff solver; reports ignition delay}
\begin{tutnarrative}
ignition01\_h2o2 integrated on the ANALYTIC Jacobian.

The case is ignition01\_h2o2 unchanged (closed rigid vessel, stoichiometric H2/O2 in N2 at 1100 K, the O Conaire 2004 mechanism: modifiedArrhenius, thirdBody, falloff with Troe, every reaction reversible by detailed balance) with ONE added key in the vessel's solver block, `jacobian analytic;`. Rosenbrock23 then builds W = I - h d J from the rate laws' own derivatives (the species columns), the chain rule (the energy row) and one difference (the T column) instead of differencing all ten columns. The Jacobian is AUDITED against the finite difference at analytic Jacobians 1, 2, 4, 8, ... and the run refuses by name if a column disagrees: Rosenbrock23 is not a W-method, and a wrong J would only make it silently slower or less accurate. The answer is the same as ignition01\_h2o2's within its golden's tolerances.

H2 : O2 : N2 = 2 : 1 : 4 (stoichiometric H2 + 1/2 O2 $\to$ H2O, 57\% N2)

mode adiabatic + energy gasConstantV: the first law dU = 0 in a fixed volume, with the heat release coming from the species FORMATION enthalpies (curated thermo) -- no per-reaction $\Delta$H is supplied. Temperature, and the radicals H/O/OH/HO2, are integrated by the stiff solver until branching runs away.

solver \{ integrator Rosenbrock23; \} -- the L-stable stiff method. Swap to RK4 (see the ignition01\_h2o2\_rk4 sibling) and the same step blows up.

\end{tutnarrative}
\tutreadme{The same H$_2$/O$_2$ ignition as ignition01\_h2o2 (closed rigid vessel, 1100 K, 5 bar, O Conaire 2004 mechanism, Rosenbrock23), with one added key in the vessel's solver block:}
\tutreadme{A Rosenbrock method solves, at every step, a linear system with the matrix W = I - h$\cdot$d$\cdot$J, where J = $\partial$f/$\partial$y is the Jacobian of the right-hand side over the state y = (n\_1 ... n\_9, T). By default the integrator builds J by finite differences: it perturbs each of the ten state entries in turn and evaluates the whole right-hand side again --- eleven evaluations per Jacobian, every one of them recomputing 21 rate constants and 21 equilibrium constants.}
\tutreadme{With jacobian analytic; the reactor builds J from the rate laws themselves:}
\tutreadme{- species columns: $\partial$(dn\_i/dt)/$\partial$n\_k = $\Sigma$\_j $\nu$\_ij $\partial$r\_j/$\partial$c\_k, from the derivative of each mass-action leg (the product of the other factors --- never r\_j/c\_k, which would divide by a zero radical), the third-body term k0$\cdot$eff\_k of R5--R9 and R15, the fall-off blend of R9 and R15 (Troe, differentiated in [M] alone), and the reverse leg k/K\_c with K\_c computed once per reaction; - the energy row by the chain rule on those columns (dT/dt = -$\Sigma$ u\_i (dn\_i/dt) / $\Sigma$ n\_i C\_v,i); - the T column by one difference of the whole right-hand side.}
\tutreadme{Rosenbrock23 is not a W-method: its order conditions assume the $\ast$true$\ast$ Jacobian. A Jacobian with a missing term does not crash anything --- the steps still pass the error test, just smaller, or less accurate --- so nobody would see it. The run therefore measures its claim: at analytic Jacobians 1, 2, 4, 8, ... it also builds a central finite-difference one at the same state, and prints}
\tutreadme{A column that differs by more than 1e-4 of its own scale refuses the run by name (row, column, both numbers).}
\begin{tutgolden}
\goldrow{kpi}{bomb}{P\_final}{1121202.82397}
\goldrow{kpi}{bomb}{T\_final}{3037.00819983}
\goldrow{kpi}{bomb}{n\_H2O2\_final}{8.00367416066e-11}
\goldrow{kpi}{bomb}{n\_H2\_final}{1.62963087713e-06}
\goldrow{kpi}{bomb}{n\_HO2\_final}{7.1975833136e-10}
\goldrow{kpi}{bomb}{n\_H\_final}{4.49541510153e-07}
\goldrow{kpi}{bomb}{n\_N2\_final}{2.857145e-05}
\goldrow{kpi}{bomb}{n\_O2\_final}{5.33369140811e-07}
\goldrow{kpi}{bomb}{n\_OH\_final}{1.20519239058e-06}
\goldrow{kpi}{bomb}{n\_O\_final}{1.86407481061e-07}
\goldrow{kpi}{bomb}{n\_water\_final}{1.18257622566e-05}
\goldrow{kpi}{bomb}{t\_end}{0.0002}
\goldrow{kpi}{bomb}{totalMoles\_final}{4.44021534514e-05}
\goldrow{kpi}{campaign}{element\_H\_closure\_rel}{1.77908230772e-15}
\goldmore{17}
\end{tutgolden}

\clearpage
\tutentry{ignition01\_h2o2\_rk4}{tutorials/batch/combustion/ignition01\_h2o2\_rk4}
\begin{tutscheme}
\begin{tikzpicture}[>={Stealth[length=2mm]}, every node/.style={font=\scriptsize}]
  \node[draw, rounded corners=2pt, align=center, minimum width=1.9cm, minimum height=0.85cm, fill=black!4] (bomb) at (0.0,-0.0) {\textbf{bomb}\\[-1pt]{\tiny batchReactor}};
\end{tikzpicture}
\end{tutscheme}
\tutwhat{Adiabatic constant-V H2/O2 ignition at 1100 K (O Conaire 2004 H2/O2 kinetics, NASA-9 thermo); EXPLICIT RK4 -- fails on the stiff chemistry (the lesson)}
\begin{tutnarrative}
ignition01\_h2o2 -- closed rigid vessel, stoichiometric H2/O2 in N2 at 1100 K.

H2 : O2 : N2 = 2 : 1 : 4 (stoichiometric H2 + 1/2 O2 $\to$ H2O, 57\% N2)

mode adiabatic + energy gasConstantV: the first law dU = 0 in a fixed volume, with the heat release coming from the species FORMATION enthalpies (curated thermo) -- no per-reaction $\Delta$H is supplied. Temperature, and the radicals H/O/OH/HO2, are integrated by the stiff solver until branching runs away.

solver \{ integrator Rosenbrock23; \} -- the L-stable stiff method. Swap to RK4 (see the ignition01\_h2o2\_rk4 sibling) and the same step blows up.

\end{tutnarrative}
\tutreadme{This is the failure half of the stiffness lesson. It runs the exact same adiabatic constant-volume H$_2$/O$_2$ ignition as the sibling case [ignition01\_h2o2](../ignition01\_h2o2/README.md), but with the explicit RK4 integrator instead of the stiff Rosenbrock23 --- at the same fixed step (deltaT = 1e-6 s).}
\tutreadme{The first step is still plausible (T 1100 $\to$ 1110 K); the second already throws the state out of the physical region (T near 8000 K, H$_2$ at zero), and the third is nan. The engine refuses rather than exit 0:}
\tutreadme{That blow-up is the result. The stiff sibling prints a stiffness ratio of order 10? for this chemistry: an explicit method is stability-limited to a step of order the $\ast$fastest$\ast$ scale, so a step chosen for the $\ast$induction$\ast$ is murdered by stability long before accuracy matters.}
\tutreadme{Compare side-by-side with ignition01\_h2o2, whose L-stable adaptive sub-stepping integrates the $\ast$same$\ast$ chemistry to a clean burned state (\textasciitilde{}3078 K), and with ignition01\_h2o2\_eulersi, which is A-stable but first order and fails honestly by name. $\ast$See, then decide$\ast$, applied to numerics.}
\tutreadme{This case ships a .expect-nonconvergence marker, so bin/runTests counts it as an EXPECTED-FAIL --- its whole point is to fail.}
\tutreadme{Same as the sibling: the ? Conaire et al. (2004) H$_2$/O$_2$ mechanism (constant/reactions, shared with ignition02\_h2air\_slack) and NASA-9 heat capacities from NASA/TP-2002-211556 for H$_2$, O$_2$, N$_2$ and water. Until 2026-10-25 it ran a GRI-Mech 3.0 subset, replaced because GRI-Mech grants no redistribution (DEV.md ?4c C164); the failure it teaches did not change.}

\clearpage
\tutentry{ignition01\_h2o2\_rop}{tutorials/batch/combustion/ignition01\_h2o2\_rop}
\begin{tutscheme}
\begin{tikzpicture}[>={Stealth[length=2mm]}, every node/.style={font=\scriptsize}]
  \node[draw, rounded corners=2pt, align=center, minimum width=1.9cm, minimum height=0.85cm, fill=black!4] (bomb) at (0.0,-0.0) {\textbf{bomb}\\[-1pt]{\tiny batchReactor}};
\end{tikzpicture}
\end{tutscheme}
\tutwhat{Adiabatic constant-V H2/O2 ignition with rates of production per species and reaction at three instants at 1100 K (O Conaire 2004 H2/O2 kinetics, NASA-9 thermo); Rosenbrock23 stiff solver; reports ignition delay}
\begin{tutnarrative}
ignition01\_h2o2 plus the RATES OF PRODUCTION report.

The case is ignition01\_h2o2 unchanged (closed rigid vessel, stoichiometric H2/O2 in N2 at 1100 K) with ONE added block on the vessel, `rateOfProduction`: at three instants of the 1 us write grid it prints, per species, which reactions make and destroy it, and writes the whole matrix c\_ij = nu\_ij r\_j V [kmol/s] to rateOfProduction.csv. The instants were read from ignition01\_h2o2's own trajectory: 3 us (induction, T 1108 K, the radical pool still growing), 7 us (the ramp, T 1766 K, between the 6 and 8 us rows where dT/dt peaks) and 15 us (burned gas, T 3037 K, the chain reactions and their reverses nearly cancel). `pathFlux ( H O );` adds, at the same instants, where the atoms of H and O go: species-to-species paths (the atoms of each reaction shared between its reactants and products in proportion to the atoms taken), printed, written to pathFlux.csv and drawn as one Graphviz file per element and instant.

H2 : O2 : N2 = 2 : 1 : 4 (stoichiometric H2 + 1/2 O2 $\to$ H2O, 57\% N2)

mode adiabatic + energy gasConstantV: the first law dU = 0 in a fixed volume, with the heat release coming from the species FORMATION enthalpies (curated thermo) -- no per-reaction $\Delta$H is supplied. Temperature, and the radicals H/O/OH/HO2, are integrated by the stiff solver until branching runs away.

solver \{ integrator Rosenbrock23; \} -- the L-stable stiff method. Swap to RK4 (see the ignition01\_h2o2\_rk4 sibling) and the same step blows up.

\end{tutnarrative}
\tutreadme{The same H$_2$/O$_2$ ignition as ignition01\_h2o2 (closed rigid vessel, 1100 K, 5 bar, O Conaire 2004 mechanism, Rosenbrock23), with one added block on the vessel:}
\tutreadme{At the first committed step that reaches each time (here the 1 $\mu$s write grid hits them exactly) the run prints, for every species, its net rate of change, its gross production and consumption, and the top reactions on each side with their share of that gross. The whole matrix goes to rateOfProduction.csv (columns t\_s, species, reaction, nu, rate\_kmol\_m3\_s, contribution\_kmol\_s, plus a NET row and the integrator's own DNDT\_ODE row per species).}
\tutreadme{pathFlux ( H O ); adds the element path flux: at the same instants, for each listed element, the net flux of its atoms from each species to each other (pathFlux.csv: t\_s, element, from, to, reaction, flux\_kmol\_atoms\_s, one row per reaction and directed pair plus a NET row per pair), the top largest printed with their share and leading reaction, and one Graphviz file pathFlux\_<element>\_<t>.dot per element and instant (dot -Tpng renders it). For reaction $\ast$j$\ast$ and element $\ast$e$\ast$ the atoms leaving reactant A for product B are f\_AB = V ($\nu$?\_A a\_eA)($\nu$?\_B a\_eB) / N\_ej, with N\_ej the atoms of $\ast$e$\ast$ on one side; the identity "atoms in minus atoms out of species i = a\_ei $\cdot$ net\_i" is held by the same gate.}
\tutreadme{The three instants were read from ignition01\_h2o2's trajectory:}
\tutreadme{Definitions: r\_j is the net rate of reaction $\ast$j$\ast$ (kmol/m$^3$/s, the very number the integrator uses, third-body concentration and fall-off included), c\_ij = $\nu$\_ij r\_j V (kmol/s) is what reaction $\ast$j$\ast$ does to species $\ast$i$\ast$, and $\Sigma$\_j c\_ij equals the integrator's dn$_i$/dt exactly (held by bin/curate/check\_rate\_of\_production.py). Theory Guide: $\ast$Rates of production of a reaction mechanism$\ast$.}
\tutreadme{The solver is Rosenbrock23, as in ignition01\_h2o2, whose \_rk4 and \_eulersi siblings show the explicit integrators failing on this chemistry; this case adds nothing to that lesson and only reads the state the integrator reaches.}
\begin{tutgolden}
\goldrow{kpi}{bomb}{P\_final}{1121202.82397}
\goldrow{kpi}{bomb}{T\_final}{3037.00819983}
\goldrow{kpi}{bomb}{n\_H2O2\_final}{8.00367416065e-11}
\goldrow{kpi}{bomb}{n\_H2\_final}{1.62963087713e-06}
\goldrow{kpi}{bomb}{n\_HO2\_final}{7.19758331359e-10}
\goldrow{kpi}{bomb}{n\_H\_final}{4.49541510152e-07}
\goldrow{kpi}{bomb}{n\_N2\_final}{2.857145e-05}
\goldrow{kpi}{bomb}{n\_O2\_final}{5.33369140811e-07}
\goldrow{kpi}{bomb}{n\_OH\_final}{1.20519239058e-06}
\goldrow{kpi}{bomb}{n\_O\_final}{1.86407481061e-07}
\goldrow{kpi}{bomb}{n\_water\_final}{1.18257622566e-05}
\goldrow{kpi}{bomb}{t\_end}{0.0002}
\goldrow{kpi}{bomb}{totalMoles\_final}{4.44021534514e-05}
\goldrow{kpi}{campaign}{element\_H\_closure\_rel}{1.51815023592e-14}
\goldmore{17}
\end{tutgolden}

\clearpage
\tutentry{ignition01\_h2o2\_seulex}{tutorials/batch/combustion/ignition01\_h2o2\_seulex}
\begin{tutscheme}
\begin{tikzpicture}[>={Stealth[length=2mm]}, every node/.style={font=\scriptsize}]
  \node[draw, rounded corners=2pt, align=center, minimum width=1.9cm, minimum height=0.85cm, fill=black!4] (bomb) at (0.0,-0.0) {\textbf{bomb}\\[-1pt]{\tiny batchReactor}};
\end{tikzpicture}
\end{tutscheme}
\tutwhat{Adiabatic constant-V H2/O2 ignition at 1100 K (O Conaire 2004 H2/O2 kinetics, NASA-9 thermo); seulex stiff solver; reports ignition delay}
\begin{tutnarrative}
the ignition01\_h2o2 vessel with `integrator seulex`: closed rigid vessel, stoichiometric H2/O2 in N2 at 1100 K.

H2 : O2 : N2 = 2 : 1 : 4 (stoichiometric H2 + 1/2 O2 $\to$ H2O, 57\% N2)

mode adiabatic + energy gasConstantV: the first law dU = 0 in a fixed volume, with the heat release coming from the species FORMATION enthalpies (curated thermo) -- no per-reaction $\Delta$H is supplied. Temperature, and the radicals H/O/OH/HO2, are integrated by the stiff solver until branching runs away.

solver \{ integrator Rosenbrock23; \} -- the L-stable stiff method. Swap to RK4 (see the ignition01\_h2o2\_rk4 sibling) and the same step blows up.

\end{tutnarrative}
\tutreadme{The same closed vessel, mechanism and write grid as ignition01\_h2o2 (read its README below), integrated by seulex (extrapolated linearly-implicit Euler, Hairer \& Wanner 1996 IV.9). A small mechanism is where an extrapolation method could lose to a Rosenbrock one; here it does not (the same CPU as Rodas4, 0.106 s against 0.105 s, Rosenbrock23 0.117 s) and the ignition delay and every printed species agree with Rodas4. Theory Guide, section $\ast$Extrapolating the linearly-implicit Euler method$\ast$.}
\tutreadme{A closed rigid vessel holds a stoichiometric H$_2$/O$_2$ charge diluted in N$_2$ at 1100 K and 5 bar. A trace H-radical seed starts chain branching; after a short induction the radical pool explodes and the mixture ignites --- a sharp temperature jump from 1100 K to \textasciitilde{}3078 K.}
\tutreadme{and plot trajectory.csv (T and the radicals H/O/OH vs time). You will see:}
\tutreadme{The ignition delay $\tau$\_ign (the time of maximum dT/dt, printed in the run's epilogue on the 1 $\mu$s write grid) is $\approx$ 8 $\mu$s here. The trace H seed is what makes it that short: the same mechanism without a seed, at 2 atm, gives the 43 $\mu$s of the shock-tube pin in ignition02\_h2air\_slack.}
\tutreadme{The radical lifetimes are far shorter than the induction --- the solver prints a stiffness ratio of order 10?. That is stiffness. An explicit method (RK4) is stability-limited to a step of order the $\ast$fastest$\ast$ scale, so it cannot take the 1 $\mu$s steps this case uses. The solver here is Rosenbrock23, an L-stable linearly-implicit method that steps by $\ast$accuracy$\ast$, not stability; at verbosity 3 it prints its stiffness verdict and step counts (the no-silent- crutch credo).}
\tutreadme{See it fail on purpose --- two siblings run the same problem at the same step with other integrators:}
\begin{tutgolden}
\goldrow{kpi}{bomb}{P\_final}{1121202.82293}
\goldrow{kpi}{bomb}{T\_final}{3037.00819759}
\goldrow{kpi}{bomb}{n\_H2O2\_final}{8.00367412435e-11}
\goldrow{kpi}{bomb}{n\_H2\_final}{1.62963087041e-06}
\goldrow{kpi}{bomb}{n\_HO2\_final}{7.19758326731e-10}
\goldrow{kpi}{bomb}{n\_H\_final}{4.49541506383e-07}
\goldrow{kpi}{bomb}{n\_N2\_final}{2.857145e-05}
\goldrow{kpi}{bomb}{n\_O2\_final}{5.33369138823e-07}
\goldrow{kpi}{bomb}{n\_OH\_final}{1.20519238467e-06}
\goldrow{kpi}{bomb}{n\_O\_final}{1.8640747939e-07}
\goldrow{kpi}{bomb}{n\_water\_final}{1.18257622681e-05}
\goldrow{kpi}{bomb}{t\_end}{0.0002}
\goldrow{kpi}{bomb}{totalMoles\_final}{4.44021534429e-05}
\goldrow{kpi}{campaign}{element\_H\_closure\_rel}{1.23266682828e-12}
\goldmore{17}
\end{tutgolden}

\clearpage
\tutentry{ignition02\_h2air\_slack}{tutorials/batch/combustion/ignition02\_h2air\_slack}
\begin{tutscheme}
\begin{tikzpicture}[>={Stealth[length=2mm]}, every node/.style={font=\scriptsize}]
  \node[draw, rounded corners=2pt, align=center, minimum width=1.9cm, minimum height=0.85cm, fill=black!4] (bomb) at (0.0,-0.0) {\textbf{bomb}\\[-1pt]{\tiny batchReactor}};
\end{tikzpicture}
\end{tutscheme}
\tutwhat{Stoich H2/air ignition, 1100 K / 2 atm (Slack shock tube; O Conaire 2004 kinetics)}
\begin{tutnarrative}
/ flowsheetDict Stoichiometric H2/air autoignition at 1100 K / 2 atm against the Slack shock-tube point (O Conaire 2004 kinetics): the radical pool (H, O, OH, HO2) builds exponentially through the induction time, then runaway. A stiff-integrator showcase -- compare the ignition delay with the measured shock-tube value.

\end{tutnarrative}
\tutreadme{Stoichiometric H$_2$/air, 1100 K, 2 atm, constant volume --- the reflected-shock condition of Slack (1977) as compiled in ? Conaire et al. 2004 (Fig 8). Kinetics: the ? Conaire H$_2$/O$_2$ mechanism, Table I verbatim (19 reactions; two Troe falloffs; two summed pairs; A in cm$^3$-mol-s, Ea in kcal --- the engine converts and announces every one).}
\tutreadme{Result: $\tau$\_ign (max dT/dt) = 43 $\mu$s vs the measured \textasciitilde{}40--60 $\mu$s window.}
\tutreadme{Two lessons: 1. No seed. This mechanism self-initiates, and the 43 $\mu$s delay is the chain-BRANCHING chemistry (R1-R3), not the dissociation R5 (H$_2$ + M = 2H, Ea = 105 kcal/mol): MEASURED 2026-10-31, when R5-R8 first ran at their true rate (DEV.md 4c C185 item 9 -- they had run at ZERO), the delay stayed 43 $\mu$s, and removing R5 from the mechanism leaves it at 43 $\mu$s too. (This paragraph used to credit R5 with the induction; that was written while R5 ran at zero and was never true.) With the ignition01-style trace-H seed the delay collapses to 12 $\mu$s: initial conditions are physics, not knobs. 2. The engine now prints tau\_ign in the epilogue (time of max dT/dt on the write grid) whenever a vessel ignites --- no CSV spelunking.}
\tutreadme{Run ignition01\_h2o2 (the same kinetics, 5 bar, trace-H seeded) for the sibling case and the stiffness lesson (\_rk4).}
\begin{tutgolden}
\goldrow{kpi}{bomb}{P\_final}{555478.877115}
\goldrow{kpi}{bomb}{T\_final}{3002.22466196}
\goldrow{kpi}{bomb}{n\_H2O2\_final}{2.91490173643e-11}
\goldrow{kpi}{bomb}{n\_H2\_final}{9.52635145411e-07}
\goldrow{kpi}{bomb}{n\_HO2\_final}{3.22097300598e-10}
\goldrow{kpi}{bomb}{n\_H\_final}{3.11093339693e-07}
\goldrow{kpi}{bomb}{n\_N2\_final}{1.391515e-05}
\goldrow{kpi}{bomb}{n\_O2\_final}{3.15344857993e-07}
\goldrow{kpi}{bomb}{n\_OH\_final}{6.97676339692e-07}
\goldrow{kpi}{bomb}{n\_O\_final}{1.28141634456e-07}
\goldrow{kpi}{bomb}{n\_water\_final}{5.93268981723e-06}
\goldrow{kpi}{bomb}{t\_end}{0.0003}
\goldrow{kpi}{bomb}{totalMoles\_final}{2.22530823808e-05}
\goldrow{kpi}{campaign}{element\_H\_closure\_rel}{3.78247165177e-15}
\goldmore{17}
\end{tutgolden}

\clearpage
\tutentry{ignition03\_h2air\_constp}{tutorials/batch/combustion/ignition03\_h2air\_constp}
\begin{tutscheme}
\begin{tikzpicture}[>={Stealth[length=2mm]}, every node/.style={font=\scriptsize}]
  \node[draw, rounded corners=2pt, align=center, minimum width=1.9cm, minimum height=0.85cm, fill=black!4] (bomb) at (0.0,-0.0) {\textbf{bomb}\\[-1pt]{\tiny batchReactor}};
\end{tikzpicture}
\end{tutscheme}
\tutwhat{Adiabatic constant-P H2/O2 ignition at 1100 K (O Conaire 2004 H2/O2 kinetics, NASA-9 thermo); Rosenbrock23 stiff solver; reports ignition delay}
\begin{tutnarrative}
the SAME charge as ignition01\_h2o2, but the vessel is a free-floating adiabatic cell at CONSTANT PRESSURE (5 bar) instead of a rigid one. (This is the cell OpenFOAM's chemistry solver integrates.)

mode adiabatic + energy gasConstantP: the first law is dH = 0 at fixed P, dT/dt = - sum\_i h\_i(T) dn\_i/dt / sum\_i n\_i Cp\_i(T), V = n R T / P, with h\_i the ideal-gas molar enthalpy on the formation datum (the datum gasConstantV takes u\_i = h\_i - R T from). V is NOT an input: it follows from (n, T, P) at every state, so the holdup declares no V.

What it shows, against ignition01\_h2o2 (constant V, same T, P, composition): ignition is SLOWER and the burned gas COOLER. At constant V the heat of reaction raises the temperature of a gas that cannot expand, and the pressure climbs; at constant P the same heat also does the expansion work P dV and is spread over the larger Cp (= Cv + R per mole), so the temperature rises less. The cooler trajectory lowers the rate constants and lengthens the induction. The burned-gas volume is V = n R T / P of its own final state (hand-check in the README).

solver \{ integrator Rosenbrock23; \} -- the L-stable stiff method. `jacobian analytic` is refused for this basis by name (debt: c\_k = n\_k/V moves with every n\_j and with T); the finite-difference Jacobian is used.

\end{tutnarrative}
\tutreadme{ignition01\_h2o2 closes the charge in a rigid vessel. This case lets the same charge (1100 K, 5 bar, H$_2$ : O$_2$ : N$_2$ = 2 : 1 : 4, trace-H seed, the same ? Conaire mechanism and thermo) expand against a fixed pressure --- the adiabatic constant-pressure cell that OpenFOAM's chemistry solver integrates.}
\tutreadme{Closed, adiabatic, ideal gas, P fixed: the first law is dH = 0, so}
\tutreadme{with h$_i$ the ideal-gas molar enthalpy on the formation datum (the datum gasConstantV takes u$_i$ = h$_i$ - RT from) and Cp$_i$ the ideal-gas heat capacity. V is not an input: it is derived from (n, T, P) at every state, so the holdup declares no V (declaring one that disagrees with n R T / P is refused by name), and every reader of the vessel --- the time directories' internalStates/bomb, the rate-of-production report --- sees the V of its own state. The rates use c$_i$ = n$_i$ / V with that V. Theory Guide, sec:batch-gas-constp.}
\tutreadme{At constant V the heat release goes into a gas that cannot expand and the pressure climbs; at constant P the same heat also does the expansion work P dV and is spread over the larger Cp = Cv + R, so the temperature rises less and the burned gas is cooler. The cooler path lowers the rate constants, so the runaway is later and slower.}
\tutreadme{At 200 $\mu$s the run ends at n = 4.39569e-5 kmol, T = 2805.845 K, P = 5 bar: V = n R T / P = 4.39569e-5 $\times$ 1000 $\times$ 8.314463 $\times$ 2805.845 / 5e5 = 2.05095e-3 m$^3$, which is the V written in 0.0002/internalStates/bomb (2.050948e-3). The vessel has grown from 1.0 L to 2.05 L at fixed pressure. (Before 2026-11-21 the burned gas ended 22 K hotter: OH was a hand-fitted cubic Cp, and it is now the NASA row of the gas-species table, DEV.md 4c C458.)}
\tutreadme{The same balance on the real LLNL H$_2$ mechanism (O'Connaire 2004, local copy), stoichiometric H$_2$/air, 1100 K, 1 atm, against Cantera's IdealGasConstPressureReactor, agrees on $\tau$, T\_final (to the mK), V\_final and the final mole fractions --- docs/reviews/2026-11-19-chemkin-real-mechanism/REPORT.md section 8.}
\begin{tutgolden}
\goldrow{kpi}{bomb}{P\_final}{500000}
\goldrow{kpi}{bomb}{T\_final}{2805.84503362}
\goldrow{kpi}{bomb}{n\_H2O2\_final}{3.52297379621e-11}
\goldrow{kpi}{bomb}{n\_H2\_final}{1.25269330246e-06}
\goldrow{kpi}{bomb}{n\_HO2\_final}{3.31508723291e-10}
\goldrow{kpi}{bomb}{n\_H\_final}{2.776347193e-07}
\goldrow{kpi}{bomb}{n\_N2\_final}{2.857145e-05}
\goldrow{kpi}{bomb}{n\_O2\_final}{4.30243577405e-07}
\goldrow{kpi}{bomb}{n\_OH\_final}{8.50008647441e-07}
\goldrow{kpi}{bomb}{n\_O\_final}{1.07986690752e-07}
\goldrow{kpi}{bomb}{n\_water\_final}{1.24664840301e-05}
\goldrow{kpi}{bomb}{t\_end}{0.0002}
\goldrow{kpi}{bomb}{totalMoles\_final}{4.39568677059e-05}
\goldrow{kpi}{campaign}{element\_H\_closure\_rel}{9.65448665656e-14}
\goldmore{17}
\end{tutgolden}

\clearpage
\tutentry{ignition04\_h2air\_sensitivity}{tutorials/batch/combustion/ignition04\_h2air\_sensitivity}
\begin{tutscheme}
\begin{tikzpicture}[>={Stealth[length=2mm]}, every node/.style={font=\scriptsize}]
  \node[draw, rounded corners=2pt, align=center, minimum width=1.9cm, minimum height=0.85cm, fill=black!4] (bomb) at (0.0,-0.0) {\textbf{bomb}\\[-1pt]{\tiny batchReactor}};
\end{tikzpicture}
\end{tutscheme}
\tutwhat{Ignition-delay sensitivity of stoich H2/air, 1100 K / 2 atm (brute force, N+1 integrations)}
\begin{tutnarrative}
/ flowsheetDict WHICH REACTION GOVERNS THE IGNITION? The ignition02 vessel (stoichiometric H2/air, 1100 K, 2 atm, constant volume, O Conaire 2004) with ONE added block, `sensitivity \{ ... \}`. The engine integrates the vessel N+1 times (the base run, then once per reaction with that reaction's NET rate multiplied by `factor`) and prints, sorted by |S\_j|,

S\_j = d ln tau / d ln k\_j = ( ln tau\_j - ln tau\_0 ) / ln factor

with tau = the time T first exceeds T0 + 400 K, linearly interpolated between the two accepted states that bracket it. The multiplier scales forward and reverse together, so the equilibrium (the final state) is the same in every run. Read the sign first: NEGATIVE S speeds the ignition up when that reaction is faster (chain branching: R1, H + O2 = O + OH, is the dominant negative entry), POSITIVE S slows it down (chain termination: R9, H + O2 (+M) = HO2 (+M), competes with R1 for the same H atoms). The signs and magnitudes are the run's own, checked reaction by reaction against Cantera 3.2.0 (docs/reviews/2026-11-19-chemkin-real-mechanism/REPORT.md section 10). `sensitivity.csv` carries every reaction.

\end{tutnarrative}
\tutreadme{The stoichiometric H$_2$/air vessel of ignition02\_h2air\_slack (1100 K, 2 atm, constant volume, ? Conaire 2004, Table I verbatim) with one added block:}
\tutreadme{The engine integrates the vessel N + 1 times (the base run, then once per reaction with that reaction's $\ast$net$\ast$ rate multiplied by 1.1) and prints, sorted by , the normalised sensitivity of the ignition delay}
\tutreadme{with $\tau$ = the time T first exceeds T$_0$ + 400 K, linearly interpolated between the two accepted states that bracket it (not the write-grid maximum of dT/dt, whose resolution is the write interval). sensitivity.csv carries all 21 reactions; the golden pins $\tau$$_0$ and the six that are not noise.}
\tutreadme{The signs and sizes are the run's own. They were checked reaction by reaction against Cantera 3.2.0 (set\_multiplier, same definition of $\tau$): on the same mechanism and thermo the two agree to 2$\cdot$10?? in S for every reaction and to 6$\cdot$10?? in $\tau$$_0$ (docs/reviews/2026-11-19-chemkin-real-mechanism/REPORT.md, section 10).}
\tutreadme{$\ast$ The multiplier scales forward and reverse together (k\_rev = k/K\_c), so every multiplied run ends in the same equilibrium: gate check\_ignition\_sensitivity holds all 21 final temperatures to 10?? of the base run (measured 7$\cdot$10?$^1$$^2$), and S of R1 to a hand recomputation from two runs of the case with R1's A edited. $\ast$ This case is not ignition02's holdup. ignition02 states 2 atm but loads 2.5$\cdot$10?? kmol into 1 L, which at 1100 K is 2.26 atm; here the moles are p V / R T for 2 atm exactly and the air is 2 : 1 : 3.76, so a Cantera run at 2 atm is the same problem. $\tau$$_0$ is 45.07 $\mu$s here.}
\tutreadme{Cost: 22 integrations of an adaptive stiff solver, about 7 s. No adjoint, no direct method: the brute force is what a student can read. Theory Guide sec:ignition-sensitivity.}
\begin{tutgolden}
\goldrow{kpi}{bomb}{P\_final}{490941.728141}
\goldrow{kpi}{bomb}{T\_final}{2991.58757699}
\goldrow{kpi}{bomb}{n\_H2O2\_final}{2.4034952022e-11}
\goldrow{kpi}{bomb}{n\_H2\_final}{8.56640617167e-07}
\goldrow{kpi}{bomb}{n\_HO2\_final}{2.7469237546e-10}
\goldrow{kpi}{bomb}{n\_H\_final}{2.86010657238e-07}
\goldrow{kpi}{bomb}{n\_N2\_final}{1.23242606244e-05}
\goldrow{kpi}{bomb}{n\_O2\_final}{2.84485169167e-07}
\goldrow{kpi}{bomb}{n\_OH\_final}{6.25381238756e-07}
\goldrow{kpi}{bomb}{n\_O\_final}{1.1754891456e-07}
\goldrow{kpi}{bomb}{n\_water\_final}{5.24296030408e-06}
\goldrow{kpi}{bomb}{sens\_S\_R10\_HO2\_H\_H2\_O2}{-0.049033738442}
\goldrow{kpi}{bomb}{sens\_S\_R11\_HO2\_H\_2OH}{-0.0346645937846}
\goldrow{kpi}{bomb}{sens\_S\_R1\_H\_O2\_O\_OH}{-1.02673194138}
\goldmore{24}
\end{tutgolden}

\clearpage
\tutentry{kinetics01\_oh\_decay}{tutorials/batch/combustion/kinetics01\_oh\_decay}
\begin{tutscheme}
\begin{tikzpicture}[>={Stealth[length=2mm]}, every node/.style={font=\scriptsize}]
  \node[draw, rounded corners=2pt, align=center, minimum width=1.9cm, minimum height=0.85cm, fill=black!4] (bomb) at (0.0,-0.0) {\textbf{bomb}\\[-1pt]{\tiny batchReactor}};
\end{tikzpicture}
\end{tutscheme}
\tutwhat{Flash-photolysis cell: a trace of OH decaying in H2/N2 at 1000 K, 1 bar, isothermal (O Conaire 2004 H2/O2 kinetics); the signal the EduTool how-do-we-know-k fits for k}
\begin{tutnarrative}
the cell of a flash-photolysis experiment: a trace of OH made at t = 0 in a hot H2/N2 bath, followed in time.

One isothermal closed vessel at constant volume (the "cell"), the whole O Conaire 2004 H2/O2 mechanism (19 reactions, the same constant/reactions as ignition01\_h2o2) so that EVERYTHING the radical can do is in the model -- not only the reaction under study. `rateOfProduction` prints, at one instant, which reactions make and destroy each species: the secondary chemistry that contaminates a measured decay, listed by name.

The page that runs this case (EduTool `how-do-we-know-k`) rewrites T, the charge, the H2 and OH mole fractions, endTime and the instant of the rate-of-production report; this file is the experiment at its defaults.

THE JACOBIAN is the analytic one (the default; audited against a finite difference on a doubling schedule). Until 2026-11-22 this case declared the finite-difference Jacobian: at the first Jacobian of a vessel holding H2O2 alone, every radical at zero, the audit read d(n\_HO2)/d(n\_HO2) = -0.9618 against -0.9665 and refused. The analytic derivative was right; the audit's REFERENCE was biased -- a forward two-point difference of a rate quadratic in a species at zero (2 HO2, 2 OH) is off by k h, and with the audit's step larger than the trace parent itself the bias was O(1). The audit now takes the second-order one-sided difference there (C468, Theory Guide sec:batch-analytic-jacobian).

\end{tutnarrative}
\tutreadme{An isothermal closed vessel (1000 K, 1 bar, 1 litre) holds an H$_2$/N$_2$ bath (H$_2$ mole fraction 0.01) and a trace of OH (1e-6) made at t = 0: the radical a flash would leave. The whole ? Conaire et al. 2004 H$_2$/O$_2$ mechanism is switched on (the constant/reactions of ignition01\_h2o2), so the OH can do everything the mechanism lets it do.}
\tutreadme{and plot trajectory.csv (bomb.n\_OH against t) on a log axis. With H$_2$ in a large excess the OH decays by OH + H$_2$ = H$_2$O + H alone, as a first-order law: ln n\_OH is a straight line of slope -k?, with k? = k$\cdot$c\_H$_2$, the pseudo-first-order constant of the experimenter's method. Dividing k? by the H$_2$ concentration c = x P / (R T) gives k, which the run's own log announces ([kinetics] A: 2.16e+08 cm3/mol/s $\to$ 216000 (SI ...), the law the engine evaluates, k = A T? exp(-Ea/RT), b = 1.51).}
\tutreadme{rateOfProduction prints, at 5 $\mu$s, which reactions make and destroy each species. With the trace seed the table says OH is removed by R3 alone; raise the OH seed (0/internalStates/bomb) and the secondary chemistry that contaminates a measured decay appears in it by name.}
\tutreadme{The EduTool How do we know k? (EduTools, Reaction engineering) runs this case under other temperatures, H$_2$ concentrations and seeds, fits k?, k and the Arrhenius line from the signal, and compares each recovered k with the engine's. Its sibling case is kinetics02\_h2o2\_falloff.}
\tutreadme{The case declares jacobian finiteDifference; on purpose; the reason is in the header of system/flowsheetDict.}
\begin{tutgolden}
\goldrow{kpi}{bomb}{E\_reaction\_seg0\_kJ}{-7.4718269738e-07}
\goldrow{kpi}{bomb}{E\_reaction\_total\_kJ}{-7.4718269738e-07}
\goldrow{kpi}{bomb}{P\_final}{99998.0419037}
\goldrow{kpi}{bomb}{T\_final}{1000}
\goldrow{kpi}{bomb}{n\_H2O2\_final}{3.02390216516e-17}
\goldrow{kpi}{bomb}{n\_H2\_final}{1.20257973521e-07}
\goldrow{kpi}{bomb}{n\_HO2\_final}{4.43350182676e-20}
\goldrow{kpi}{bomb}{n\_H\_final}{1.20261498562e-11}
\goldrow{kpi}{bomb}{n\_N2\_final}{1.1906717973e-05}
\goldrow{kpi}{bomb}{n\_O2\_final}{2.77252517794e-19}
\goldrow{kpi}{bomb}{n\_OH\_final}{1.68804633753e-17}
\goldrow{kpi}{bomb}{n\_O\_final}{5.70235034959e-17}
\goldrow{kpi}{bomb}{n\_water\_final}{1.20268649695e-11}
\goldrow{kpi}{bomb}{t\_end}{0.0001}
\goldmore{31}
\end{tutgolden}

\clearpage
\tutentry{kinetics02\_h2o2\_falloff}{tutorials/batch/combustion/kinetics02\_h2o2\_falloff}
\begin{tutscheme}
\begin{tikzpicture}[>={Stealth[length=2mm]}, every node/.style={font=\scriptsize}]
  \node[draw, rounded corners=2pt, align=center, minimum width=1.9cm, minimum height=0.85cm, fill=black!4] (bomb) at (0.0,-0.0) {\textbf{bomb}\\[-1pt]{\tiny batchReactor}};
\end{tikzpicture}
\end{tutscheme}
\tutwhat{Fall-off cell: a trace of H2O2 decaying in N2 at 1000 K, 30 bar, isothermal (O Conaire 2004 H2/O2 kinetics, R15 Troe fall-off); the signal the EduTool how-do-we-know-k fits for k}
\begin{tutnarrative}
the fall-off cell: a trace of H2O2 in N2 at 1000 K and 30 bar, isothermal, followed as a first-order decay.

The whole O Conaire 2004 H2/O2 mechanism (19 reactions, the same constant/reactions as ignition01\_h2o2) is loaded, so everything the H2O2 and the OH it makes can do is in the model; the reaction under study is R15, H2O2 + M = 2 OH + M, a Troe fall-off. The page that runs this case (EduTool `how-do-we-know-k`) rewrites T, P, the charge, the H2O2 and N2 mole fractions and the window; this file is the experiment at its defaults.

THE JACOBIAN is the analytic one (the default; audited against a finite difference on a doubling schedule). Until 2026-11-22 this case declared the finite-difference Jacobian: at the first Jacobian of a vessel holding H2O2 alone, every radical at zero, the audit read d(n\_HO2)/d(n\_HO2) = -0.9618 against -0.9665 and refused. The analytic derivative was right; the audit's REFERENCE was biased -- a forward two-point difference of a rate quadratic in a species at zero (2 HO2, 2 OH) is off by k h, and with the audit's step larger than the trace parent itself the bias was O(1). The audit now takes the second-order one-sided difference there (C468, Theory Guide sec:batch-analytic-jacobian).

\end{tutnarrative}
\tutreadme{An isothermal closed vessel (1000 K, 30 bar, 1 litre) holds N$_2$ and a trace of H$_2$O$_2$ (1e-10 mole fraction). The reaction H$_2$O$_2$ + M = 2 OH + M (R15 of the ? Conaire et al. 2004 mechanism) is a Troe fall-off: its first-order constant lies between the low-pressure limit k$_0$[M] and the high-pressure limit k$\infty$, and the reduced pressure P\_r = k$_0$[M]/k$\infty$ is the coordinate of the curve. At these conditions P\_r $\approx$ 0.67, near the middle.}
\tutreadme{and plot bomb.n\_H2O2 in trajectory.csv on a log axis: a straight line of slope -k. The Lindemann blend k\_L = k$\infty$ P\_r/(1 + P\_r) is too sharp; the engine multiplies it by Troe's broadening factor F (Gilbert, Luther and Troe 1983, Eqs. 6.3 to 6.7), so k = k\_L F with F < 1. The seed is small so that the OH the decomposition makes does not attack the H$_2$O$_2$ (OH + H$_2$O$_2$ = H$_2$O + HO$_2$), which would add a chain contribution to the decay: at a seed of 1e-5 the decay reads about 1.8 times too fast (measured natively at 1000 K, 1 and 30 bar).}
\tutreadme{The EduTool How do we know k? reruns this cell at other pressures and temperatures and draws the measured k beside the record's Lindemann and Troe laws. Its sibling case is kinetics01\_oh\_decay, where the same page measures a bimolecular rate constant.}
\tutreadme{The case declares jacobian finiteDifference; on purpose; the reason is in the header of system/flowsheetDict.}
\begin{tutgolden}
\goldrow{kpi}{bomb}{E\_reaction\_seg0\_kJ}{7.73058062009e-09}
\goldrow{kpi}{bomb}{E\_reaction\_total\_kJ}{7.73058062009e-09}
\goldrow{kpi}{bomb}{P\_final}{3000024.40213}
\goldrow{kpi}{bomb}{T\_final}{1000}
\goldrow{kpi}{bomb}{n\_H2O2\_final}{1.13461079795e-17}
\goldrow{kpi}{bomb}{n\_H2\_final}{1.53561486095e-23}
\goldrow{kpi}{bomb}{n\_HO2\_final}{3.98850760902e-18}
\goldrow{kpi}{bomb}{n\_H\_final}{4.84086720312e-19}
\goldrow{kpi}{bomb}{n\_N2\_final}{0.000360819999964}
\goldrow{kpi}{bomb}{n\_O2\_final}{5.29311342566e-19}
\goldrow{kpi}{bomb}{n\_OH\_final}{7.2070981054e-14}
\goldrow{kpi}{bomb}{n\_O\_final}{2.83640396419e-17}
\goldrow{kpi}{bomb}{n\_water\_final}{3.29270525034e-17}
\goldrow{kpi}{bomb}{t\_end}{0.005}
\goldmore{25}
\end{tutgolden}

\clearpage
\tutentry{methane01\_h2o2\_core}{tutorials/batch/combustion/methane01\_h2o2\_core}
\begin{tutscheme}
\begin{tikzpicture}[>={Stealth[length=2mm]}, every node/.style={font=\scriptsize}]
  \node[draw, rounded corners=2pt, align=center, minimum width=1.9cm, minimum height=0.85cm, fill=black!4] (vessel) at (0.0,-0.0) {\textbf{vessel}\\[-1pt]{\tiny batchReactor}};
\end{tikzpicture}
\end{tutscheme}
\tutwhat{the H2/O2 core ALONE: methane is a spectator, nothing starts (Petersen, Davidson \& Hanson 1999, mixture 2, 1244 K, 114.4 atm)}
\begin{tutnarrative}
/ flowsheetDict -- ONE adiabatic constant-volume vessel, the sets of constant/reactions. the H2/O2 core ALONE: methane is a spectator, nothing starts. Sets read: OConaire2004.

\end{tutnarrative}
\tutreadme{The vessel holds methane, oxygen and nitrogen behind a reflected shock (Petersen, Davidson and Hanson 1999, Table 1 mixture 2, the point of Table 2 page 86: 1244 K, 114.4 atm, measured delay 254 us). The only reactions are the 21 records of the OConaire2004 set, none of which names a carbon species: the methane is a spectator, the temperature does not move, and the engine prints no ignition delay. A mechanism is not a mixture.}
\tutreadme{The ladder is methane01 to methane06 (the same vessel and start, only the reaction sets differ), and methane07 adds the sensitivity. The EduTool methane-assembly ("Assemble methane combustion") runs them in the browser. The constants of the carbon records are ONE mechanism's values (Miller and Bowman 1989, Appendix A), not evaluated constants; the shock-tube point is Petersen, Davidson and Hanson, J. Propulsion and Power 15 (1999) 82-91. Reverse rates come from equilibrium.}
\begin{tutgolden}
\goldrow{kpi}{campaign}{element\_C\_closure\_rel}{0}
\goldrow{kpi}{campaign}{element\_H\_closure\_rel}{0}
\goldrow{kpi}{campaign}{element\_N\_closure\_rel}{0}
\goldrow{kpi}{campaign}{element\_O\_closure\_rel}{1.6325521776e-14}
\goldrow{kpi}{campaign}{element\_worst\_closure\_rel}{1.6325521776e-14}
\goldrow{kpi}{campaign}{energy\_balance\_available}{0}
\goldrow{kpi}{campaign}{energy\_records\_invalid}{0}
\goldrow{kpi}{campaign}{energy\_records\_valid}{0}
\goldrow{kpi}{campaign}{mass\_closure\_rel}{2.72283043428e-15}
\goldrow{kpi}{campaign}{mass\_kg\_external\_out}{0}
\goldrow{kpi}{campaign}{mass\_kg\_final}{0.0293067386387}
\goldrow{kpi}{campaign}{mass\_kg\_initial}{0.0293067386387}
\goldrow{kpi}{campaign}{mass\_residual\_kg}{7.97972798949e-17}
\goldrow{kpi}{campaign}{moles\_kmol\_external\_out}{0}
\goldmore{29}
\end{tutgolden}

\clearpage
\tutentry{methane02\_ch4\_attack}{tutorials/batch/combustion/methane02\_ch4\_attack}
\begin{tutscheme}
\begin{tikzpicture}[>={Stealth[length=2mm]}, every node/.style={font=\scriptsize}]
  \node[draw, rounded corners=2pt, align=center, minimum width=1.9cm, minimum height=0.85cm, fill=black!4] (vessel) at (0.0,-0.0) {\textbf{vessel}\\[-1pt]{\tiny batchReactor}};
\end{tikzpicture}
\end{tutscheme}
\tutwhat{+ the attack on CH4: CH3 appears, and has nowhere to go (Petersen, Davidson \& Hanson 1999, mixture 2, 1244 K, 114.4 atm)}
\begin{tutnarrative}
/ flowsheetDict -- ONE adiabatic constant-volume vessel, the sets of constant/reactions. + the attack on CH4: CH3 appears, and has nowhere to go. Sets read: OConaire2004, MB1989\_attack.

\end{tutnarrative}
\tutreadme{The same vessel with five records of Miller and Bowman 1989 Appendix A (reactions 3, 4, 5, 5a, 6: CH4 attacked by O2, H, OH, O and HO2) added to the core. A trace of CH3 appears and nothing consumes it: the temperature barely moves and there is no ignition.}
\tutreadme{The ladder is methane01 to methane06 (the same vessel and start, only the reaction sets differ), and methane07 adds the sensitivity. The EduTool methane-assembly ("Assemble methane combustion") runs them in the browser. The constants of the carbon records are ONE mechanism's values (Miller and Bowman 1989, Appendix A), not evaluated constants; the shock-tube point is Petersen, Davidson and Hanson, J. Propulsion and Power 15 (1999) 82-91. Reverse rates come from equilibrium.}
\begin{tutgolden}
\goldrow{kpi}{campaign}{element\_C\_closure\_rel}{1.81394232913e-15}
\goldrow{kpi}{campaign}{element\_H\_closure\_rel}{1.81394232913e-15}
\goldrow{kpi}{campaign}{element\_N\_closure\_rel}{0}
\goldrow{kpi}{campaign}{element\_O\_closure\_rel}{7.25578745599e-16}
\goldrow{kpi}{campaign}{element\_worst\_closure\_rel}{1.81394232913e-15}
\goldrow{kpi}{campaign}{energy\_balance\_available}{0}
\goldrow{kpi}{campaign}{energy\_records\_invalid}{0}
\goldrow{kpi}{campaign}{energy\_records\_valid}{0}
\goldrow{kpi}{campaign}{mass\_closure\_rel}{7.01437534767e-09}
\goldrow{kpi}{campaign}{mass\_kg\_external\_out}{0}
\goldrow{kpi}{campaign}{mass\_kg\_final}{0.0293067384332}
\goldrow{kpi}{campaign}{mass\_kg\_initial}{0.0293067386387}
\goldrow{kpi}{campaign}{mass\_residual\_kg}{2.05568465028e-10}
\goldrow{kpi}{campaign}{moles\_kmol\_external\_out}{0}
\goldmore{29}
\end{tutgolden}

\clearpage
\tutentry{methane03\_methyl}{tutorials/batch/combustion/methane03\_methyl}
\begin{tutscheme}
\begin{tikzpicture}[>={Stealth[length=2mm]}, every node/.style={font=\scriptsize}]
  \node[draw, rounded corners=2pt, align=center, minimum width=1.9cm, minimum height=0.85cm, fill=black!4] (vessel) at (0.0,-0.0) {\textbf{vessel}\\[-1pt]{\tiny batchReactor}};
\end{tikzpicture}
\end{tutscheme}
\tutwhat{+ the methyl step: CH3 $\to$ CH3O $\to$ CH2O (Petersen, Davidson \& Hanson 1999, mixture 2, 1244 K, 114.4 atm)}
\begin{tutnarrative}
/ flowsheetDict -- ONE adiabatic constant-volume vessel, the sets of constant/reactions. + the methyl step: CH3 $\to$ CH3O $\to$ CH2O. Sets read: OConaire2004, MB1989\_attack, MB1989\_methyl.

\end{tutnarrative}
\tutreadme{Twelve more records: CH3 with O2 (a chain branching), O, HO2 and H, and CH3O to formaldehyde. The mixture ignites, at a small fraction of the measured delay, and the carbon stops at formaldehyde. A mechanism can ignite for the wrong reasons.}
\tutreadme{The ladder is methane01 to methane06 (the same vessel and start, only the reaction sets differ), and methane07 adds the sensitivity. The EduTool methane-assembly ("Assemble methane combustion") runs them in the browser. The constants of the carbon records are ONE mechanism's values (Miller and Bowman 1989, Appendix A), not evaluated constants; the shock-tube point is Petersen, Davidson and Hanson, J. Propulsion and Power 15 (1999) 82-91. Reverse rates come from equilibrium.}
\begin{tutgolden}
\goldrow{kpi}{campaign}{element\_C\_closure\_rel}{1.57208335191e-15}
\goldrow{kpi}{campaign}{element\_H\_closure\_rel}{3.62788465826e-16}
\goldrow{kpi}{campaign}{element\_N\_closure\_rel}{0}
\goldrow{kpi}{campaign}{element\_O\_closure\_rel}{3.627893728e-16}
\goldrow{kpi}{campaign}{element\_worst\_closure\_rel}{1.57208335191e-15}
\goldrow{kpi}{campaign}{energy\_balance\_available}{0}
\goldrow{kpi}{campaign}{energy\_records\_invalid}{0}
\goldrow{kpi}{campaign}{energy\_records\_valid}{0}
\goldrow{kpi}{campaign}{mass\_closure\_rel}{5.11457227272e-06}
\goldrow{kpi}{campaign}{mass\_kg\_external\_out}{0}
\goldrow{kpi}{campaign}{mass\_kg\_final}{0.0293065887473}
\goldrow{kpi}{campaign}{mass\_kg\_initial}{0.0293067386387}
\goldrow{kpi}{campaign}{mass\_residual\_kg}{1.49891432846e-07}
\goldrow{kpi}{campaign}{moles\_kmol\_external\_out}{0}
\goldmore{29}
\end{tutgolden}

\clearpage
\tutentry{methane04\_formyl}{tutorials/batch/combustion/methane04\_formyl}
\begin{tutscheme}
\begin{tikzpicture}[>={Stealth[length=2mm]}, every node/.style={font=\scriptsize}]
  \node[draw, rounded corners=2pt, align=center, minimum width=1.9cm, minimum height=0.85cm, fill=black!4] (vessel) at (0.0,-0.0) {\textbf{vessel}\\[-1pt]{\tiny batchReactor}};
\end{tikzpicture}
\end{tutscheme}
\tutwhat{+ CH2O $\to$ HCO $\to$ CO (and CH2): the carbon reaches CO, not yet CO2 (Petersen, Davidson \& Hanson 1999, mixture 2, 1244 K, 114.4 atm)}
\begin{tutnarrative}
/ flowsheetDict -- ONE adiabatic constant-volume vessel, the sets of constant/reactions. + CH2O $\to$ HCO $\to$ CO (and CH2): the carbon reaches CO, not yet CO2. Sets read: OConaire2004, MB1989\_attack, MB1989\_methyl, MB1989\_formyl.

\end{tutnarrative}
\tutreadme{Eighteen more records take formaldehyde through HCO to carbon monoxide and oxidise the CH2 of the methyl step. The ignition delay hardly moves (it was decided by the branching) and the final temperature rises.}
\tutreadme{The ladder is methane01 to methane06 (the same vessel and start, only the reaction sets differ), and methane07 adds the sensitivity. The EduTool methane-assembly ("Assemble methane combustion") runs them in the browser. The constants of the carbon records are ONE mechanism's values (Miller and Bowman 1989, Appendix A), not evaluated constants; the shock-tube point is Petersen, Davidson and Hanson, J. Propulsion and Power 15 (1999) 82-91. Reverse rates come from equilibrium.}
\begin{tutgolden}
\goldrow{kpi}{campaign}{element\_C\_closure\_rel}{2.29766028357e-15}
\goldrow{kpi}{campaign}{element\_H\_closure\_rel}{1.81394232913e-15}
\goldrow{kpi}{campaign}{element\_N\_closure\_rel}{0}
\goldrow{kpi}{campaign}{element\_O\_closure\_rel}{5.26044590559e-15}
\goldrow{kpi}{campaign}{element\_worst\_closure\_rel}{5.26044590559e-15}
\goldrow{kpi}{campaign}{energy\_balance\_available}{0}
\goldrow{kpi}{campaign}{energy\_records\_invalid}{0}
\goldrow{kpi}{campaign}{energy\_records\_valid}{0}
\goldrow{kpi}{campaign}{mass\_closure\_rel}{5.08816962847e-06}
\goldrow{kpi}{campaign}{mass\_kg\_external\_out}{0}
\goldrow{kpi}{campaign}{mass\_kg\_final}{0.0293065895211}
\goldrow{kpi}{campaign}{mass\_kg\_initial}{0.0293067386387}
\goldrow{kpi}{campaign}{mass\_residual\_kg}{1.49117657451e-07}
\goldrow{kpi}{campaign}{moles\_kmol\_external\_out}{0}
\goldmore{29}
\end{tutgolden}

\clearpage
\tutentry{methane05\_c1\_chain}{tutorials/batch/combustion/methane05\_c1\_chain}
\begin{tutscheme}
\begin{tikzpicture}[>={Stealth[length=2mm]}, every node/.style={font=\scriptsize}]
  \node[draw, rounded corners=2pt, align=center, minimum width=1.9cm, minimum height=0.85cm, fill=black!4] (vessel) at (0.0,-0.0) {\textbf{vessel}\\[-1pt]{\tiny batchReactor}};
\end{tikzpicture}
\end{tutscheme}
\tutwhat{+ CO $\to$ CO2: the whole C1 chain, from CH4 to CO2 (Petersen, Davidson \& Hanson 1999, mixture 2, 1244 K, 114.4 atm)}
\begin{tutnarrative}
/ flowsheetDict -- ONE adiabatic constant-volume vessel, the sets of constant/reactions. + CO $\to$ CO2: the whole C1 chain, from CH4 to CO2. Sets read: OConaire2004, MB1989\_attack, MB1989\_methyl, MB1989\_formyl, MB1989\_carbonMonoxide.

\end{tutnarrative}
\tutreadme{Four more records turn CO into CO2, where most of the heat is released. The C1 chain is complete, from CH4 to CO2, and the delay is still far too short.}
\tutreadme{The ladder is methane01 to methane06 (the same vessel and start, only the reaction sets differ), and methane07 adds the sensitivity. The EduTool methane-assembly ("Assemble methane combustion") runs them in the browser. The constants of the carbon records are ONE mechanism's values (Miller and Bowman 1989, Appendix A), not evaluated constants; the shock-tube point is Petersen, Davidson and Hanson, J. Propulsion and Power 15 (1999) 82-91. Reverse rates come from equilibrium.}
\begin{tutgolden}
\goldrow{kpi}{campaign}{element\_C\_closure\_rel}{1.81394232913e-15}
\goldrow{kpi}{campaign}{element\_H\_closure\_rel}{8.46506420261e-16}
\goldrow{kpi}{campaign}{element\_N\_closure\_rel}{0}
\goldrow{kpi}{campaign}{element\_O\_closure\_rel}{1.813946864e-16}
\goldrow{kpi}{campaign}{element\_worst\_closure\_rel}{1.81394232913e-15}
\goldrow{kpi}{campaign}{energy\_balance\_available}{0}
\goldrow{kpi}{campaign}{energy\_records\_invalid}{0}
\goldrow{kpi}{campaign}{energy\_records\_valid}{0}
\goldrow{kpi}{campaign}{mass\_closure\_rel}{4.63455769928e-06}
\goldrow{kpi}{campaign}{mass\_kg\_external\_out}{0}
\goldrow{kpi}{campaign}{mass\_kg\_final}{0.029306602815}
\goldrow{kpi}{campaign}{mass\_kg\_initial}{0.0293067386387}
\goldrow{kpi}{campaign}{mass\_residual\_kg}{1.35823771199e-07}
\goldrow{kpi}{campaign}{moles\_kmol\_external\_out}{0}
\goldmore{29}
\end{tutgolden}

\clearpage
\tutentry{methane06\_assembled}{tutorials/batch/combustion/methane06\_assembled}
\begin{tutscheme}
\begin{tikzpicture}[>={Stealth[length=2mm]}, every node/.style={font=\scriptsize}]
  \node[draw, rounded corners=2pt, align=center, minimum width=1.9cm, minimum height=0.85cm, fill=black!4] (vessel) at (0.0,-0.0) {\textbf{vessel}\\[-1pt]{\tiny batchReactor}};
\end{tikzpicture}
\end{tutscheme}
\tutwhat{+ the ethane route (two CH3 recombine): the assembled mechanism (Petersen, Davidson \& Hanson 1999, mixture 2, 1244 K, 114.4 atm)}
\begin{tutnarrative}
/ flowsheetDict -- ONE adiabatic constant-volume vessel, the sets of constant/reactions. + the ethane route (two CH3 recombine): the assembled mechanism. Sets read: OConaire2004, MillerBowman1989.

\end{tutnarrative}
\tutreadme{Seven records of the ethane route (2CH3(+M) = C2H6(+M) with its Troe fall-off, and the attack on ethane) join the chain. At this pressure the recombination of two methyl radicals dominates, the ignition delay grows by about an order of magnitude, and it lands near the measurement. This is the set MillerBowman1989 (46 of the 235 reactions of Appendix A) beside the 21 of the OConaire2004 core.}
\tutreadme{The ladder is methane01 to methane06 (the same vessel and start, only the reaction sets differ), and methane07 adds the sensitivity. The EduTool methane-assembly ("Assemble methane combustion") runs them in the browser. The constants of the carbon records are ONE mechanism's values (Miller and Bowman 1989, Appendix A), not evaluated constants; the shock-tube point is Petersen, Davidson and Hanson, J. Propulsion and Power 15 (1999) 82-91. Reverse rates come from equilibrium.}
\begin{tutgolden}
\goldrow{kpi}{campaign}{element\_C\_closure\_rel}{1.20929488609e-16}
\goldrow{kpi}{campaign}{element\_H\_closure\_rel}{1.4511538633e-15}
\goldrow{kpi}{campaign}{element\_N\_closure\_rel}{0}
\goldrow{kpi}{campaign}{element\_O\_closure\_rel}{2.3581309232e-15}
\goldrow{kpi}{campaign}{element\_worst\_closure\_rel}{2.3581309232e-15}
\goldrow{kpi}{campaign}{energy\_balance\_available}{0}
\goldrow{kpi}{campaign}{energy\_records\_invalid}{0}
\goldrow{kpi}{campaign}{energy\_records\_valid}{0}
\goldrow{kpi}{campaign}{mass\_closure\_rel}{4.51752559371e-06}
\goldrow{kpi}{campaign}{mass\_kg\_external\_out}{0}
\goldrow{kpi}{campaign}{mass\_kg\_final}{0.0293066062448}
\goldrow{kpi}{campaign}{mass\_kg\_initial}{0.0293067386387}
\goldrow{kpi}{campaign}{mass\_residual\_kg}{1.32393941869e-07}
\goldrow{kpi}{campaign}{moles\_kmol\_external\_out}{0}
\goldmore{29}
\end{tutgolden}

\clearpage
\tutentry{methane07\_sensitivity}{tutorials/batch/combustion/methane07\_sensitivity}
\begin{tutscheme}
\begin{tikzpicture}[>={Stealth[length=2mm]}, every node/.style={font=\scriptsize}]
  \node[draw, rounded corners=2pt, align=center, minimum width=1.9cm, minimum height=0.85cm, fill=black!4] (vessel) at (0.0,-0.0) {\textbf{vessel}\\[-1pt]{\tiny batchReactor}};
\end{tikzpicture}
\end{tutscheme}
\tutwhat{the assembled mechanism with the sensitivity of the ignition delay to every reaction (Petersen, Davidson \& Hanson 1999, mixture 2, 1244 K, 114.4 atm)}
\begin{tutnarrative}
/ flowsheetDict -- ONE adiabatic constant-volume vessel, the sets of constant/reactions. the assembled mechanism with the sensitivity of the ignition delay to every reaction. Sets read: OConaire2004, MillerBowman1989.

\end{tutnarrative}
\tutreadme{methane06 with sensitivity \{ response tauIgnition; factor 1.1; reactions all; top 8; \}: the engine integrates the vessel once per reaction with that reaction's net rate multiplied and prints S\_j = d ln tau / d ln k\_j. Two reactions of the methyl radical head the table with opposite signs: the branching CH3 + O2 = CH3O + O (A8, negative) and the recombination of two CH3 (A1, positive).}
\tutreadme{The ladder is methane01 to methane06 (the same vessel and start, only the reaction sets differ), and methane07 adds the sensitivity. The EduTool methane-assembly ("Assemble methane combustion") runs them in the browser. The constants of the carbon records are ONE mechanism's values (Miller and Bowman 1989, Appendix A), not evaluated constants; the shock-tube point is Petersen, Davidson and Hanson, J. Propulsion and Power 15 (1999) 82-91. Reverse rates come from equilibrium.}
\begin{tutgolden}
\goldrow{kpi}{campaign}{element\_C\_closure\_rel}{1.20929488609e-16}
\goldrow{kpi}{campaign}{element\_H\_closure\_rel}{1.4511538633e-15}
\goldrow{kpi}{campaign}{element\_N\_closure\_rel}{0}
\goldrow{kpi}{campaign}{element\_O\_closure\_rel}{2.3581309232e-15}
\goldrow{kpi}{campaign}{element\_worst\_closure\_rel}{2.3581309232e-15}
\goldrow{kpi}{campaign}{energy\_balance\_available}{0}
\goldrow{kpi}{campaign}{energy\_records\_invalid}{0}
\goldrow{kpi}{campaign}{energy\_records\_valid}{0}
\goldrow{kpi}{campaign}{mass\_closure\_rel}{4.51752559371e-06}
\goldrow{kpi}{campaign}{mass\_kg\_external\_out}{0}
\goldrow{kpi}{campaign}{mass\_kg\_final}{0.0293066062448}
\goldrow{kpi}{campaign}{mass\_kg\_initial}{0.0293067386387}
\goldrow{kpi}{campaign}{mass\_residual\_kg}{1.32393941869e-07}
\goldrow{kpi}{campaign}{moles\_kmol\_external\_out}{0}
\goldmore{38}
\end{tutgolden}

\clearpage
\tutentry{nox01\_zeldovich\_postflame}{tutorials/batch/combustion/nox01\_zeldovich\_postflame}
\begin{tutscheme}
\begin{tikzpicture}[>={Stealth[length=2mm]}, every node/.style={font=\scriptsize}]
  \node[draw, rounded corners=2pt, align=center, minimum width=1.9cm, minimum height=0.85cm, fill=black!4] (postflame) at (0.0,-0.0) {\textbf{postflame}\\[-1pt]{\tiny batchReactor}};
\end{tikzpicture}
\end{tutscheme}
\tutwhat{Thermal NOx in post-flame H2/air gas, 2125 K / 1 atm (Glarborg 2018 N chemistry)}
\begin{tutnarrative}
/ flowsheetDict Thermal NOx in post-flame H2/air gas at 2125 K / 1 atm: the Zeldovich N-chemistry (Glarborg 2018 rate set) integrated in a batch reactor -- NO grows on the SECONDS scale while the O/H radical pool stays equilibrated, the classic separation of scales that makes thermal NOx a post-flame problem.

\end{tutnarrative}
\tutreadme{Hot burnt gas (lean H$_2$/air, 2125 K, 1 atm --- Homer \& Sutton Flame 1 as compiled in Glarborg 2018) held isothermal, starting from the equilibrated H/O radical pool with all nitrogen still as N$_2$. The Glarborg 2018 nitrogen chemistry (extended Zeldovich + N$_2$O routes, Tables 2/5 verbatim) then forms NO.}
\tutreadme{What is not carried. Glarborg's NNH route (R20--R24) is absent since 2026-10-25: NNH's thermo had come from GRI-Mech 3.0's thermo file, which grants no redistribution, and no public-domain primary for it was found (NASA/TP-2002-211556 has no NNH entry). The heat capacities of the flame species are NASA-9 polynomials from that report. On 2026-07-02, with the NNH route and the earlier thermo, the same table read 5227 / 5167 ppm (98.9 \%), 0.50 s and 1.56 s.}
\tutreadme{1. The cross-engine pin. The kinetic plateau lands on the Gibbs equilibrium computed by a $\ast$different engine$\ast$ (steady Gibbs minimisation vs stiff batch integration). Kinetics choose the path and the clock; the destination belongs to the thermodynamics. Reverse rates come from the same equilibrium constants --- so this agreement is a genuine consistency check of the mechanism transcription + thermo, not a tautology. 2. NOx forms in engineering time. Ignition takes microseconds (ignition02: 43 $\mu$s); thermal NO takes seconds at 2125 K. That timescale gap is the whole basis of NOx control by quench: cool the gas before the Zeldovich clock runs.}
\tutreadme{The burnt-gas initial pool was computed by a gibbsReactor at 2125 K/1 atm (14 species on 2026-07-02; it still agrees with the current 13-species thermo to the printed digits) (25 \% H$_2$ / 17.5 \% O$_2$ / 57.5 \% N$_2$ feed); the equilibrium NO share was returned to N$_2$/O$_2$ elementally.}
\begin{tutgolden}
\goldrow{kpi}{postflame}{P\_final}{100741.759342}
\goldrow{kpi}{postflame}{T\_final}{2125}
\goldrow{kpi}{postflame}{n\_H2O2\_final}{1.03158641521e-12}
\goldrow{kpi}{postflame}{n\_H2\_final}{4.85008397878e-09}
\goldrow{kpi}{postflame}{n\_HO2\_final}{1.62257148216e-11}
\goldrow{kpi}{postflame}{n\_H\_final}{6.02690833353e-10}
\goldrow{kpi}{postflame}{n\_N2O\_final}{1.37519144211e-12}
\goldrow{kpi}{postflame}{n\_N2\_final}{3.72506317809e-06}
\goldrow{kpi}{postflame}{n\_NH\_final}{2.49395194269e-15}
\goldrow{kpi}{postflame}{n\_NO\_final}{2.9429747874e-08}
\goldrow{kpi}{postflame}{n\_N\_final}{2.26203990434e-14}
\goldrow{kpi}{postflame}{n\_O2\_final}{3.0502331009e-07}
\goldrow{kpi}{postflame}{n\_OH\_final}{2.77512944532e-08}
\goldrow{kpi}{postflame}{n\_O\_final}{2.16840040834e-09}
\goldmore{29}
\end{tutgolden}

\clearpage
\tutentry{nox02\_millerbowman\_denox}{tutorials/batch/combustion/nox02\_millerbowman\_denox}
\begin{tutscheme}
\begin{tikzpicture}[>={Stealth[length=2mm]}, every node/.style={font=\scriptsize}]
  \node[draw, rounded corners=2pt, align=center, minimum width=1.9cm, minimum height=0.85cm, fill=black!4] (T900) at (0.0,-0.0) {\textbf{T900}\\[-1pt]{\tiny batchReactor}};
  \node[draw, rounded corners=2pt, align=center, minimum width=1.9cm, minimum height=0.85cm, fill=black!4] (T1000) at (0.0,-1.7) {\textbf{T1000}\\[-1pt]{\tiny batchReactor}};
  \node[draw, rounded corners=2pt, align=center, minimum width=1.9cm, minimum height=0.85cm, fill=black!4] (T1050) at (0.0,-3.4) {\textbf{T1050}\\[-1pt]{\tiny batchReactor}};
  \node[draw, rounded corners=2pt, align=center, minimum width=1.9cm, minimum height=0.85cm, fill=black!4] (T1100) at (0.0,-5.1) {\textbf{T1100}\\[-1pt]{\tiny batchReactor}};
  \node[draw, rounded corners=2pt, align=center, minimum width=1.9cm, minimum height=0.85cm, fill=black!4] (T1125) at (0.0,-6.8) {\textbf{T1125}\\[-1pt]{\tiny batchReactor}};
  \node[draw, rounded corners=2pt, align=center, minimum width=1.9cm, minimum height=0.85cm, fill=black!4] (T1150) at (0.0,-8.5) {\textbf{T1150}\\[-1pt]{\tiny batchReactor}};
  \node[draw, rounded corners=2pt, align=center, minimum width=1.9cm, minimum height=0.85cm, fill=black!4] (T1175) at (0.0,-10.2) {\textbf{T1175}\\[-1pt]{\tiny batchReactor}};
  \node[draw, rounded corners=2pt, align=center, minimum width=1.9cm, minimum height=0.85cm, fill=black!4] (T1200) at (0.0,-11.9) {\textbf{T1200}\\[-1pt]{\tiny batchReactor}};
  \node[draw, rounded corners=2pt, align=center, minimum width=1.9cm, minimum height=0.85cm, fill=black!4] (T1225) at (0.0,-13.6) {\textbf{T1225}\\[-1pt]{\tiny batchReactor}};
  \node[draw, rounded corners=2pt, align=center, minimum width=1.9cm, minimum height=0.85cm, fill=black!4] (T1250) at (0.0,-15.3) {\textbf{T1250}\\[-1pt]{\tiny batchReactor}};
  \node[draw, rounded corners=2pt, align=center, minimum width=1.9cm, minimum height=0.85cm, fill=black!4] (T1275) at (0.0,-17.0) {\textbf{T1275}\\[-1pt]{\tiny batchReactor}};
  \node[draw, rounded corners=2pt, align=center, minimum width=1.9cm, minimum height=0.85cm, fill=black!4] (T1300) at (0.0,-18.7) {\textbf{T1300}\\[-1pt]{\tiny batchReactor}};
  \node[draw, rounded corners=2pt, align=center, minimum width=1.9cm, minimum height=0.85cm, fill=black!4] (T1325) at (0.0,-20.4) {\textbf{T1325}\\[-1pt]{\tiny batchReactor}};
  \node[draw, rounded corners=2pt, align=center, minimum width=1.9cm, minimum height=0.85cm, fill=black!4] (T1350) at (0.0,-22.1) {\textbf{T1350}\\[-1pt]{\tiny batchReactor}};
  \node[draw, rounded corners=2pt, align=center, minimum width=1.9cm, minimum height=0.85cm, fill=black!4] (T1400) at (0.0,-23.8) {\textbf{T1400}\\[-1pt]{\tiny batchReactor}};
  \node[draw, rounded corners=2pt, align=center, minimum width=1.9cm, minimum height=0.85cm, fill=black!4] (T1450) at (0.0,-25.5) {\textbf{T1450}\\[-1pt]{\tiny batchReactor}};
  \node[draw, rounded corners=2pt, align=center, minimum width=1.9cm, minimum height=0.85cm, fill=black!4] (T1500) at (0.0,-27.2) {\textbf{T1500}\\[-1pt]{\tiny batchReactor}};
\end{tikzpicture}
\end{tutscheme}
\tutwhat{Miller \& Bowman 1989 Fig. 57 (Lyon \& Hardy flow reactor, data set 1): thermal DeNOx window, Appendix B, 17 isothermal vessels, 0.1 s}
\begin{tutnarrative}
/$\ast$---------------------------------------------------------------------------$\ast$ nox02\_millerbowman\_denox / flowsheetDict -- one isothermal vessel per temperature of the Thermal DeNOx window (900-1500 K, 25 K steps through the window), each a plug-flow element held 0.1 s.

\end{tutnarrative}
\tutreadme{The question. NH3 is injected into hot gas that carries NO and some O2. Between about 1100 and 1400 K it removes NO; below the window nothing reacts, above it the ammonia burns to NO. Where is the window, and how deep? Miller \& Bowman (Prog. Energy Combust. Sci. 15 (1989) 287-338, doi:10.1016/0360-1285(89)90017-8) answer in Fig. 57 (p. 313) against the flow-reactor data of Lyon \& Hardy. This case runs that figure.}
\tutreadme{Setup (the paper's data set 1). NO 225 ppm, NH3 450 ppm, O2 1.23 \%, balance He, 1.1 atm, isothermal, reaction time 0.1 s, temperature 900-1500 K. A flow reactor held at a wall temperature is a plug-flow element, and a plug-flow element is a batch vessel in its own travel time, so the case is 17 isothermal batchReactor units, one per temperature (energy gasConstantP, mode isothermal), integrated to 0.1 s in one run. The temperature sweep is therefore a list of vessels, not an outerDict: choupoBatch has no outer driver.}
\tutreadme{Mechanism. Appendix B of the paper, 73 reactions, 19 species plus He, A in cm3-mol-s and E in cal/mol as printed (constant/reactions; the paper's numbers name the reactions, B1...B73, each with its page). The paper computed this figure with Appendix B. Reverse rates come from equilibrium, as the paper does. Thermochemistry: one set per run, the NASA/TP-2002-211556 NASA-9 data (stable species as case-local overrides, radicals from the NASA table); N2H2 and NNH are read from the NASA table and the RRHO-computed table respectively (Theory Guide sec:rrho-thermo; NNH has no NASA entry).}
\tutreadme{The window opens at 1100 K, reaches 17 ppm (92 \% removal) near 1225 K, and closes by 1400 K, where the ammonia has all burned and NOx is back at 204 ppm. Below 1050 K nothing happens in 0.1 s. The curve is the paper's: root-mean-square difference 2.9 ppm over the model curve's 14 read points, largest 7.8 ppm on the steep flank at 1200 K. The readings are digitised from the page image (thresholded, the axis frame found by pixel, +-3 ppm, more on the steep flank): information, not a pin. Against the Lyon \& Hardy data themselves the model (the paper's and this case's alike) runs about 25-30 K cold on the flank and 25 ppm high above the minimum -- that is the paper's own remaining disagreement with the experiment, not a difference between the two models (docs/reviews/2026-11-22-miller-bowman/compare\_paper.py).}
\tutreadme{N2O is the intermediate to watch: it peaks at 13.0 ppm near 1325 K and is gone (0.2 ppm) by 1500 K; the ammonia is exhausted from 1275 K up, and from there on the NO that remains is what the ammonia burned to.}
\tutreadme{Same transcription, same thermodynamics (the NASA-9 coefficients Choupo reads, written into the Cantera input), Cantera 3.2.0, CVODES at rtol 1e-10: the end values of NOx agree to 5 significant digits in all 17 vessels, every species at every write time to 5.7e-4 relative at worst (rtol 1e-6, atol 1e-19 kmol as shipped), temperatures identical. docs/reviews/2026-11-22-miller-bowman/compare\_cantera.py denox.}
\begin{tutgolden}
\goldrow{kpi}{T1000}{E\_reaction\_seg0\_kJ}{-2.97108044833e-08}
\goldrow{kpi}{T1000}{E\_reaction\_total\_kJ}{-2.97108044833e-08}
\goldrow{kpi}{T1000}{P\_final}{111457.5}
\goldrow{kpi}{T1000}{T\_final}{1000}
\goldrow{kpi}{T1000}{n\_H2O2\_final}{5.53382336698e-27}
\goldrow{kpi}{T1000}{n\_H2\_final}{2.55152725357e-16}
\goldrow{kpi}{T1000}{n\_HNO\_final}{4.74843402439e-14}
\goldrow{kpi}{T1000}{n\_HO2\_final}{1.10785665093e-17}
\goldrow{kpi}{T1000}{n\_H\_final}{1.62823152097e-17}
\goldrow{kpi}{T1000}{n\_He\_final}{0.0132313228152}
\goldrow{kpi}{T1000}{n\_N2H2\_final}{1.37980428254e-26}
\goldrow{kpi}{T1000}{n\_N2O\_final}{3.79636333037e-17}
\goldrow{kpi}{T1000}{n\_N2\_final}{8.77642970437e-14}
\goldrow{kpi}{T1000}{n\_NH2\_final}{2.29341613127e-16}
\goldmore{475}
\end{tutgolden}

\clearpage
\tutentry{nox03\_millerbowman\_large}{tutorials/batch/combustion/nox03\_millerbowman\_large}
\begin{tutscheme}
\begin{tikzpicture}[>={Stealth[length=2mm]}, every node/.style={font=\scriptsize}]
  \node[draw, rounded corners=2pt, align=center, minimum width=1.9cm, minimum height=0.85cm, fill=black!4] (vessel) at (0.0,-0.0) {\textbf{vessel}\\[-1pt]{\tiny batchReactor}};
\end{tikzpicture}
\end{tutscheme}
\tutwhat{Miller \& Bowman 1989 Appendix A (235 reactions, 51 species): lean CH4/air + 100 ppm NO, adiabatic constant volume, 1300 K, 1 atm}
\begin{tutnarrative}
/$\ast$---------------------------------------------------------------------------$\ast$ nox03\_millerbowman\_large / flowsheetDict -- one adiabatic constant-volume vessel with the 235-reaction Appendix A mechanism. Two declared figures (system/postDict): when it ignites and what NO it leaves, and where the NO comes from (routeIntegral below: five declared reaction classes, integrated on the dense output).

\end{tutnarrative}
\tutreadme{What it is. The whole of Appendix A of Miller \& Bowman (Prog. Energy Combust. Sci. 15 (1989) 287-338, doi:10.1016/0360-1285(89)90017-8): 235 elementary reactions among 51 species (C, H, O, N), 5 of them pressure fall-offs (one Troe, one SRI, three Lindemann) and 21 with a bare third body, every rate coefficient read from the page images and every reaction named by the number the paper gives it (A1 ... A234, A5a), with its page, in constant/reactions. It is the case that shows what the engine does with a mechanism of a size nobody types by hand.}
\tutreadme{The vessel. Lean methane/air (phi = 0.6) with 100 ppm NO, 1300 K, 1 atm, 1 L, closed, adiabatic at constant volume (mode adiabatic; energy gasConstantV, the shock-tube / engine cell). Constant volume is chosen because the analytic Jacobian, audited at run time, is available there (constant pressure integrates on the finite difference, ignition03) and because the burned gas then carries both the heat release and the pressure rise. The mechanism was built for exactly this chemistry: the CH + N2 prompt route (A151), the HCN/NCO/NH chain, the NNH route (A216, A221-A227), N2O (A203-A207), NO2 (A188-A191).}
\tutreadme{Read at commit cec127f1562d (2026-11-24); the times below are not all write points of the shipped log grid (91 points, 20 per decade from 1 us), so they are taken from a scratch copy of the case on a uniform 1 ms write grid (the same dense output; at 1, 10 and 30 ms, which the shipped grid also writes, the values are identical).}
\tutreadme{$\ast$ Induction (0-2.4 ms): the methane starts to break down and the radical pool builds, and part of the initial NO becomes NO2 (17 ppm at 1 ms; the mechanism's route to it is HO2 + NO = NO2 + OH, A188 -- not checked here by a rate-of-production report). $\ast$ Ignition at 2.4 ms (tauIgnition = 2.380 ms, maximum dT/dt on the integrator's dense output; Cantera 3.2.0, maximum dT/dt of its own steps: 2.3802 ms; a 1 ms write grid alone would say "3 ms"): T jumps to 2612 K, P from 1.0 to 2.0 atm. $\ast$ Thermal NO then grows in the burned gas, 90 \% of its final value within 9 ms of the ignition, to 15008 ppm of NO (NO2 10 ppm) at 30 ms: 99.997 \% of the constant-volume adiabatic equilibrium Cantera computes from the same charge (NOx 15019.0 ppm, 2593.63 K, 2.0457 bar; here 2593.63 K and 2.0457 bar).}
\tutreadme{Same CHEMKIN transcription, the same NASA-9 coefficients Choupo reads (written into the Cantera input by cantera\_mech.py), Cantera 3.2.0 IdealGasReactor, adiabatic, constant volume, CVODES rtol 1e-10, atol 1e-20, compared by docs/reviews/2026-11-22-miller-bowman/compare\_cantera.py large at the write times of the Choupo run (deviation max / max(x\_Cantera, 1e-9) over all times and species).}
\tutreadme{Current (2026-11-24, commit cec127f1562d, Cantera 3.2.0, atol 1e-15 kmol). The agreement depends on the write grid, because since C471 the trajectory between the integrator's steps is its dense output (a stage polynomial, not error-controlled):}
\begin{tutgolden}
\goldrow{kpi}{campaign}{element\_C\_closure\_rel}{7.04929450598e-15}
\goldrow{kpi}{campaign}{element\_H\_closure\_rel}{2.66730062388e-15}
\goldrow{kpi}{campaign}{element\_N\_closure\_rel}{3.40483552427e-15}
\goldrow{kpi}{campaign}{element\_O\_closure\_rel}{3.42851910167e-16}
\goldrow{kpi}{campaign}{element\_worst\_closure\_rel}{7.04929450598e-15}
\goldrow{kpi}{campaign}{energy\_balance\_available}{0}
\goldrow{kpi}{campaign}{energy\_records\_invalid}{0}
\goldrow{kpi}{campaign}{energy\_records\_valid}{0}
\goldrow{kpi}{campaign}{mass\_closure\_rel}{2.16692977705e-06}
\goldrow{kpi}{campaign}{mass\_kg\_external\_out}{0}
\goldrow{kpi}{campaign}{mass\_kg\_final}{0.000263335416197}
\goldrow{kpi}{campaign}{mass\_kg\_initial}{0.000263335986828}
\goldrow{kpi}{campaign}{mass\_residual\_kg}{5.70630591227e-10}
\goldrow{kpi}{campaign}{moles\_kmol\_external\_out}{0}
\goldmore{69}
\end{tutgolden}

\clearpage
\tutentry{nox03\_millerbowman\_large\_seulex}{tutorials/batch/combustion/nox03\_millerbowman\_large\_seulex}
\begin{tutscheme}
\begin{tikzpicture}[>={Stealth[length=2mm]}, every node/.style={font=\scriptsize}]
  \node[draw, rounded corners=2pt, align=center, minimum width=1.9cm, minimum height=0.85cm, fill=black!4] (vessel) at (0.0,-0.0) {\textbf{vessel}\\[-1pt]{\tiny batchReactor}};
\end{tikzpicture}
\end{tutscheme}
\tutwhat{Miller \& Bowman 1989 Appendix A (235 reactions, 51 species): lean CH4/air + 100 ppm NO, adiabatic constant volume, 1300 K, 1 atm}
\begin{tutnarrative}
/$\ast$---------------------------------------------------------------------------$\ast$ nox03\_millerbowman\_large / flowsheetDict -- one adiabatic constant-volume vessel with the 235-reaction Appendix A mechanism. The same two declared figures as nox03 (system/postDict), on the seulex integrator's own dense output.

\end{tutnarrative}
\tutreadme{What it shows. The case nox03\_millerbowman\_large (read its README for the mechanism, the vessel and the comparison with Cantera) integrated with seulex, the extrapolated linearly-implicit Euler method of Hairer \& Wanner (1996, IV.9) that OpenFOAM's chemistry uses by default, instead of Rosenbrock23. Only the integrator word differs. Theory Guide, section $\ast$Extrapolating the linearly-implicit Euler method$\ast$.}
\tutreadme{What to look at. Run it with CHOUPO\_ODE\_STATS=1 to print the work counters: about 95 accepted steps (Rodas4 takes \textasciitilde{}730), a Jacobian built on \textasciitilde{}90 attempts (the convergence monitor decides, threshold min(1e-4, rtol); Rodas4 builds one per step) and \textasciitilde{}640 factorisations of W (one per extrapolation column, because W = I - h\_j J changes with the sub-step h\_j = H/n\_j). The final state (the post-ignition equilibrium) agrees with the Rosenbrock23 case.}
\tutreadme{Between the steps. The rows written between two steps come from seulex's OWN dense output, the polynomial through the sub-steps of the step (Hairer \& Wanner IV.9, Theory Guide $\ast$Dense output: the extrapolated polynomial$\ast$), not from the driver's cubic Hermite as before 2026-11-22 (which put the rows inside the induction period 4e-2 of a column maximum away from a tight reference). Grid error against a tight reference, 1 ms grid: 2.4e-4 at rtol 1e-8 and 2.2e-4 at 1e-10 with linearSolver dense, 6.8e-5 and 4.1e-5 with sparse (the step ends themselves: 9e-5, the floor atol sets on the trace columns); the dense and sparse runs, which take the identical steps, differ by 2.4e-4 on the grid (they differed by 1.8e-2). The polynomial's own error is not fed back into the step size (that would take 695 steps instead of 95).}
\tutreadme{What it is. The whole of Appendix A of Miller \& Bowman (Prog. Energy Combust. Sci. 15 (1989) 287-338, doi:10.1016/0360-1285(89)90017-8): 235 elementary reactions among 51 species (C, H, O, N), 5 of them pressure fall-offs (one Troe, one SRI, three Lindemann) and 21 with a bare third body, every rate coefficient read from the page images and every reaction named by the number the paper gives it (A1 ... A234, A5a), with its page, in constant/reactions. It is the case that shows what the engine does with a mechanism of a size nobody types by hand.}
\tutreadme{The vessel. Lean methane/air (phi = 0.6) with 100 ppm NO, 1300 K, 1 atm, 1 L, closed, adiabatic at constant volume (mode adiabatic; energy gasConstantV, the shock-tube / engine cell). Constant volume is chosen because the analytic Jacobian, audited at run time, is available there (constant pressure integrates on the finite difference, ignition03) and because the burned gas then carries both the heat release and the pressure rise. The mechanism was built for exactly this chemistry: the CH + N2 prompt route (A151), the HCN/NCO/NH chain, the NNH route (A216, A221-A227), N2O (A203-A207), NO2 (A188-A191).}
\tutreadme{$\ast$ Induction (0-2.4 ms): the methane starts to break down and the radical pool builds, and part of the initial NO becomes NO2 (17 ppm at 1 ms; the mechanism's route to it is HO2 + NO = NO2 + OH, A188 -- not checked here by a rate-of-production report). $\ast$ Ignition at 2.4 ms (Cantera: maximum dT/dt; the write grid of 1 ms says "3 ms"): T jumps to 2612 K, P from 1.0 to 2.0 atm. $\ast$ Thermal NO then grows in the burned gas, 90 \% of its final value within 9 ms of the ignition, to 15008 ppm of NO (NO2 10 ppm) at 30 ms: 99.997 \% of the constant-volume adiabatic equilibrium Cantera computes from the same charge (NOx 15019.0 ppm, 2593.63 K, 2.0457 bar; here 2593.63 K and 2.0457 bar).}
\begin{tutgolden}
\goldrow{kpi}{campaign}{element\_C\_closure\_rel}{1.0821619674e-13}
\goldrow{kpi}{campaign}{element\_H\_closure\_rel}{1.78518620327e-13}
\goldrow{kpi}{campaign}{element\_N\_closure\_rel}{1.96994055333e-14}
\goldrow{kpi}{campaign}{element\_O\_closure\_rel}{5.73705529679e-14}
\goldrow{kpi}{campaign}{element\_worst\_closure\_rel}{1.78518620327e-13}
\goldrow{kpi}{campaign}{energy\_balance\_available}{0}
\goldrow{kpi}{campaign}{energy\_records\_invalid}{0}
\goldrow{kpi}{campaign}{energy\_records\_valid}{0}
\goldrow{kpi}{campaign}{mass\_closure\_rel}{2.16692974453e-06}
\goldrow{kpi}{campaign}{mass\_kg\_external\_out}{0}
\goldrow{kpi}{campaign}{mass\_kg\_final}{0.000263335416197}
\goldrow{kpi}{campaign}{mass\_kg\_initial}{0.000263335986828}
\goldrow{kpi}{campaign}{mass\_residual\_kg}{5.70630582662e-10}
\goldrow{kpi}{campaign}{moles\_kmol\_external\_out}{0}
\goldmore{69}
\end{tutgolden}

\subsection{crystallisation}
\grouptheory{crystallisation}
\clearpage
\tutentry{batch13\_nacl\_pitzer}{tutorials/batch/crystallisation/batch13\_nacl\_pitzer}
\begin{tutscheme}
\begin{tikzpicture}[>={Stealth[length=2mm]}, every node/.style={font=\scriptsize}]
  \node[draw, rounded corners=2pt, align=center, minimum width=1.9cm, minimum height=0.85cm, fill=black!4] (cryst) at (0.0,-0.0) {\textbf{cryst}\\[-1pt]{\tiny batchCrystalliser}};
\end{tikzpicture}
\end{tutscheme}
\tutwhat{Batch NaCl crystalliser: saturation from the Pitzer ion activity, not a solubility curve}
\begin{tutnarrative}
One closed vessel. The unit declares no saturation of its own: it asks the thermodynamics the case declares, and the case declares Pitzer.

\end{tutnarrative}
\begin{tutgolden}
\goldrow{kpi}{campaign}{element\_Cl\_closure\_rel}{9.08953975799e-16}
\goldrow{kpi}{campaign}{element\_H\_closure\_rel}{0}
\goldrow{kpi}{campaign}{element\_Na\_closure\_rel}{9.08953975799e-16}
\goldrow{kpi}{campaign}{element\_O\_closure\_rel}{0}
\goldrow{kpi}{campaign}{element\_worst\_closure\_rel}{9.08953975799e-16}
\goldrow{kpi}{campaign}{energy\_balance\_available}{0}
\goldrow{kpi}{campaign}{energy\_records\_invalid}{0}
\goldrow{kpi}{campaign}{energy\_records\_valid}{1}
\goldrow{kpi}{campaign}{mass\_closure\_rel}{3.49878229369e-16}
\goldrow{kpi}{campaign}{mass\_kg\_external\_out}{0}
\goldrow{kpi}{campaign}{mass\_kg\_final}{1299.73034249}
\goldrow{kpi}{campaign}{mass\_kg\_initial}{1299.73034249}
\goldrow{kpi}{campaign}{mass\_residual\_kg}{4.54747350886e-13}
\goldrow{kpi}{campaign}{moles\_kmol\_external\_out}{0}
\goldmore{18}
\end{tutgolden}

\clearpage
\tutentry{batch14\_nacl\_pitzer\_antisolvent}{tutorials/batch/crystallisation/batch14\_nacl\_pitzer\_antisolvent}
\begin{tutscheme}
\begin{tikzpicture}[>={Stealth[length=2mm]}, every node/.style={font=\scriptsize}]
  \node[draw, rounded corners=2pt, align=center, minimum width=1.9cm, minimum height=0.85cm, fill=black!4] (cryst) at (0.0,-0.0) {\textbf{cryst}\\[-1pt]{\tiny batchCrystalliser}};
\end{tikzpicture}
\end{tutscheme}
\tutwhat{Batch NaCl crystalliser with ethanol present: a pairwise-aqueous model announces the mixed-solvent term it does not carry}
\begin{tutnarrative}
One closed vessel, identical to batch13 except for the ethanol in the charge. The unit declares no saturation of its own: it asks the thermodynamics the case declares, and the case declares Pitzer -- which has nothing to say about a second solvent, and says so.

\end{tutnarrative}
\begin{tutgolden}
\goldrow{kpi}{campaign}{element\_C\_closure\_rel}{0}
\goldrow{kpi}{campaign}{element\_Cl\_closure\_rel}{6.05235197043e-16}
\goldrow{kpi}{campaign}{element\_H\_closure\_rel}{0}
\goldrow{kpi}{campaign}{element\_Na\_closure\_rel}{6.05235197043e-16}
\goldrow{kpi}{campaign}{element\_O\_closure\_rel}{0}
\goldrow{kpi}{campaign}{element\_worst\_closure\_rel}{6.05235197043e-16}
\goldrow{kpi}{campaign}{energy\_balance\_available}{0}
\goldrow{kpi}{campaign}{energy\_records\_invalid}{0}
\goldrow{kpi}{campaign}{energy\_records\_valid}{1}
\goldrow{kpi}{campaign}{mass\_closure\_rel}{1.3612844745e-16}
\goldrow{kpi}{campaign}{mass\_kg\_external\_out}{0}
\goldrow{kpi}{campaign}{mass\_kg\_final}{1670.28772973}
\goldrow{kpi}{campaign}{mass\_kg\_initial}{1670.28772973}
\goldrow{kpi}{campaign}{mass\_residual\_kg}{2.27373675443e-13}
\goldmore{20}
\end{tutgolden}

\subsection{drying}
\grouptheory{drying}
\clearpage
\tutentry{dryer01\_sucrose\_tray}{tutorials/batch/drying/dryer01\_sucrose\_tray}
\begin{tutscheme}
\begin{tikzpicture}[>={Stealth[length=2mm]}, every node/.style={font=\scriptsize}]
  \node[draw, rounded corners=2pt, align=center, minimum width=1.9cm, minimum height=0.85cm, fill=black!4] (dryer) at (0.0,-0.0) {\textbf{dryer}\\[-1pt]{\tiny batchDryer}};
\end{tikzpicture}
\end{tutscheme}
\tutwhat{Tray drying of wet sucrose into declared 60 C air: the classical constant-rate / falling-rate drying curve, with the wet bulb and the GAB equilibrium moisture as results}
\begin{tutnarrative}
the DRYING CURVE, both periods, in one batch.

A tray holding 2.0 kg of dry sucrose wet to X\_0 = 0.30 kg water / kg dry solid, exposed over 0.5 m2 to air of DECLARED, CONSTANT condition: 60 C (333.15 K) and Y = 0.010 kg water / kg dry N2. Nothing about that air is integrated -- it is an ENVIRONMENT, the FixedBedAdsorber's constant-carrier posture, and the water that leaves the solid leaves the campaign across that boundary (reported as the unit's DECLARED material residual, never as a silent leak).

THE HAND ESTIMATE, computed from this case's own curated records (water Antoine + Watson Hvap + the two ideal-gas Cp polynomials, the case-local GAB isotherm on sucrose) -- the run must reproduce it:

air water activity p\_v = P Y/(Y + Mv/Mc) = 1551 Pa a\_w = p\_v / Psat(333.15 K) = 0.0774 GAB equilibrium K a\_w = 0.0743 X\_eq = Xm C Ka /((1-Ka)(1-Ka+C Ka)) = 0.01580 kg/kg wet bulb (Lewis = 1) T\_wb = 300.53 K (27.4 C, a 32.6 K depression) driving force Y\_sat(T\_wb) - Y = 0.014011 kg/kg constant-rate flux R\_c = k\_Y dY = 0.05 $\ast$ 0.014011 = 7.0057e-4 kg/(m2 s) R\_c A = 3.5028e-4 kg/s latent load R\_c A lambda(T\_wb) = 0.8645 kW (the AIR's, not the campaign's -- KPI only, never ledgered) break time t\_c = m\_s (X\_0 - X\_c)/(R\_c A) = 1027.7 s falling-rate constant tau = m\_s (X\_c - X\_eq)/(R\_c A) = 595.0 s end of run X(3000 s) = X\_eq + (X\_c - X\_eq) exp(-(3000-t\_c)/tau) = 0.019584 kg/kg

What to look at in the output: trajectory.csv columns dryer.X\_kg\_kg and dryer.R\_kg\_m2\_s: the drying curve and the rate curve, engine columns, not a post-processing. R is FLAT at 7.0057e-4 until X falls through X\_c = 0.12, then decays linearly in (X - X\_eq) -- the textbook break, visible as a corner. The falling rate is DECLARED as the characteristic drying curve f = Phi\textasciicircum{}n with n = 1, which IS the linear law; the drying EduTool varies n (n < 1 reaches X\_eq in finite time, n > 1 an algebraic tail slower than the exponential), and dryer02\_characteristic\_curves runs four shapes side by side. the KPI t\_critical is OBSERVED (the interpolated crossing between two accepted states), not the analytic 1027.7 s above; they agree because dX/dt is exactly constant on that leg. THE CAMPAIGN MASS BALANCE CLOSES BECAUSE THE DRIED WATER IS AN OUTLET. The solid is charged with 2.600000 kg and ends with 2.039168 kg; the 0.560832 kg of difference is the DRIED WATER. It left the solid and went into the air, and the campaign books it as an external OUTLET STREAM of the dryer (the ledger record `dryer $\to$ (external boundary)`, "evaporated water leaving the solid into the drying air"): initial = final + external out, mass and element closure at round-off. The air stays a DECLARED ENVIRONMENT (a temperature and a humidity, not a vessel this campaign owns), but matter crossing the campaign boundary is a stream whatever receives it. Until 2026-11-09 this unit left the water off the books and the balance read a 21.6 \% shortfall that only the log explained; a violated balance on screen is a defect, so it is a stream now (DEV.md 4c C258).

the campaign energy balance is UNAVAILABLE, and says why in one sentence: the drying heat comes from air this campaign never owned. The number is not hidden -- it is the latentDuty\_kW KPI.

Method hypotheses, each announced at run time: Lewis = 1 in the adiabatic-saturation equation; the humid heat evaluated at the air's own temperature; the surface held at T\_wb while X > X\_c; the LINEAR falling-rate curve (n = 1; a modelling CHOICE, not physics, and its declared source says so); the solid's own temperature NOT integrated.

\end{tutnarrative}
\begin{tutgolden}
\goldrow{kpi}{campaign}{element\_C\_closure\_rel}{0}
\goldrow{kpi}{campaign}{element\_H\_closure\_rel}{0}
\goldrow{kpi}{campaign}{element\_O\_closure\_rel}{0}
\goldrow{kpi}{campaign}{element\_worst\_closure\_rel}{0}
\goldrow{kpi}{campaign}{energy\_balance\_available}{0}
\goldrow{kpi}{campaign}{energy\_records\_invalid}{0}
\goldrow{kpi}{campaign}{energy\_records\_valid}{0}
\goldrow{kpi}{campaign}{mass\_closure\_rel}{0}
\goldrow{kpi}{campaign}{mass\_kg\_external\_out}{0.560738660881}
\goldrow{kpi}{campaign}{mass\_kg\_final}{2.03916808853}
\goldrow{kpi}{campaign}{mass\_kg\_initial}{2.60000001375}
\goldrow{kpi}{campaign}{mass\_residual\_kg}{0}
\goldrow{kpi}{campaign}{moles\_kmol\_external\_out}{0.031126209319}
\goldrow{kpi}{campaign}{moles\_kmol\_final}{0.00802224868102}
\goldmore{23}
\end{tutgolden}

\clearpage
\tutentry{dryer02\_characteristic\_curves}{tutorials/batch/drying/dryer02\_characteristic\_curves}
\begin{tutscheme}
\begin{tikzpicture}[>={Stealth[length=2mm]}, every node/.style={font=\scriptsize}]
  \node[draw, rounded corners=2pt, align=center, minimum width=1.9cm, minimum height=0.85cm, fill=black!4] (linear) at (0.0,-0.0) {\textbf{linear}\\[-1pt]{\tiny batchDryer}};
  \node[draw, rounded corners=2pt, align=center, minimum width=1.9cm, minimum height=0.85cm, fill=black!4] (concave) at (0.0,-1.7) {\textbf{concave}\\[-1pt]{\tiny batchDryer}};
  \node[draw, rounded corners=2pt, align=center, minimum width=1.9cm, minimum height=0.85cm, fill=black!4] (convex) at (0.0,-3.4) {\textbf{convex}\\[-1pt]{\tiny batchDryer}};
  \node[draw, rounded corners=2pt, align=center, minimum width=1.9cm, minimum height=0.85cm, fill=black!4] (measured) at (0.0,-5.1) {\textbf{measured}\\[-1pt]{\tiny batchDryer}};
\end{tikzpicture}
\end{tutscheme}
\tutwhat{The characteristic drying curve: one wet sucrose tray, one declared air, four falling-rate shapes -- linear, Phi\textasciicircum{}0.5, Phi\textasciicircum{}2 and a table}
\begin{tutnarrative}
THE CHARACTERISTIC DRYING CURVE.

Four identical trays of dryer01 (2.0 kg dry sucrose at X\_0 = 0.30, 0.5 m2, air at 60 C and Y = 0.010, X\_c = 0.12) that differ in ONE thing: the shape of the falling-rate period. van Meel (Chem. Eng. Sci. 9 (1958) 36-44) observed that for many materials the drying rate normalised by the constant rate is a single function of the normalised free moisture,

Phi = (X - X\_eq)/(X\_c - X\_eq), f = R/R\_c = f(Phi), f(1) = 1, f(0) = 0,

so ONE measured curve describes the material at any air condition -- the air sets R\_c (the wet bulb) and X\_eq (the isotherm), the curve sets the rest. The shape is the material's, MEASURED; each tray declares one (src/unitOperations/heatTransfer/DryingCurve.H):

linear f = Phi the textbook tail, dryer01's law, declared concave f = Phi\textasciicircum{}0.5 rate held high: reaches X\_eq in FINITE time, t' = tau/(1 - n) = 2 tau convex f = Phi\textasciicircum{}2 rate dropping early: an algebraic tail, Phi = 1/(1 + t'/tau), slower than linear measured a (Phi, f) table S-shaped, interpolated linearly

tau = m\_s (X\_c - X\_eq)/(R\_c A) = 595.0 s (dryer01's hand estimate) is the SAME for all four -- the constant-rate period, t\_c = 1027.7 s, too. What differs is only the falling period, and the end-of-run moisture, at t' = 3000 - 1027.7 = 1972.3 s into it:

linear X = X\_eq + (X\_c - X\_eq) exp(-t'/tau) = 0.019584 (= dryer01) convex X = X\_eq + (X\_c - X\_eq)/(1 + t'/tau) = 0.03995 concave X = X\_eq = 0.015798, reached at t' = 2 tau = 1190 s -- the engine lands 3.8e-8 kg/kg BELOW it (2.4e-6 relative): RK4 stepping onto a point where f = Phi\textasciicircum{}0.5 has an infinite slope, the size (dt/tau)\textasciicircum{}2 (X\_c - X\_eq)/4 predicts. Not clamped (no clamp in this unit, by rule); it shrinks with dt. measured 0.03502 (no closed form), between linear and convex, as its points lie below the linear curve bin/curate/check\_drying\_curve.py recomputes all of this from the run.

EVERY CURVE HERE IS HYPOTHETICAL. No measured characteristic curve of sucrose is in this repository; the shapes are chosen to span what materials do, and each tray says so in its `source`. A real design reads its curve from a thin-layer drying test of ITS material.

\end{tutnarrative}
\begin{tutgolden}
\goldrow{kpi}{campaign}{element\_C\_closure\_rel}{0}
\goldrow{kpi}{campaign}{element\_H\_closure\_rel}{8.24896884633e-15}
\goldrow{kpi}{campaign}{element\_O\_closure\_rel}{8.24896884633e-15}
\goldrow{kpi}{campaign}{element\_worst\_closure\_rel}{8.24896884633e-15}
\goldrow{kpi}{campaign}{energy\_balance\_available}{0}
\goldrow{kpi}{campaign}{energy\_records\_invalid}{0}
\goldrow{kpi}{campaign}{energy\_records\_valid}{0}
\goldrow{kpi}{campaign}{mass\_closure\_rel}{5.4657133231e-15}
\goldrow{kpi}{campaign}{mass\_kg\_external\_out}{2.17896631527}
\goldrow{kpi}{campaign}{mass\_kg\_final}{8.22069230365}
\goldrow{kpi}{campaign}{mass\_kg\_initial}{10.400000055}
\goldrow{kpi}{campaign}{mass\_residual\_kg}{5.68434188608e-14}
\goldrow{kpi}{campaign}{moles\_kmol\_external\_out}{0.120952890106}
\goldrow{kpi}{campaign}{moles\_kmol\_final}{0.0356409418936}
\goldmore{80}
\end{tutgolden}

\clearpage
\tutentry{dryer03\_solid\_temperature}{tutorials/batch/drying/dryer03\_solid\_temperature}
\begin{tutscheme}
\begin{tikzpicture}[>={Stealth[length=2mm]}, every node/.style={font=\scriptsize}]
  \node[draw, rounded corners=2pt, align=center, minimum width=1.9cm, minimum height=0.85cm, fill=black!4] (dryer) at (0.0,-0.0) {\textbf{dryer}\\[-1pt]{\tiny batchDryer}};
\end{tikzpicture}
\end{tutscheme}
\tutwhat{Tray drying of wet sucrose into declared 60 C air with the SOLID TEMPERATURE integrated: the surface holds the wet bulb through the constant-rate period, then the lumped solid heats toward the air (opt-in solidTemperature integrated)}
\begin{tutnarrative}
dryer01\_sucrose\_tray with the SOLID'S TEMPERATURE integrated (C512). Everything below about the air, the wet bulb, the break and the moisture is dryer01's, unchanged; what is new is `solidTemperature integrated;` and its KPIs (T\_solid\_final, solidEnergy\_$\ast$), the trajectory column dryer.T\_solid\_K.

WHAT THE SOLID DOES (C512; Theory Guide sec:batch-dryer-solid-T): X > X\_c: T\_s IS the wet bulb, 300.52 K (dryer01's), held (hypothesis 2). h = k\_Y c\_H = 0.05 $\ast$ 1060.2 = 53.01 W/(m2 K) (Lewis = 1: the relation the wet-bulb equation is written in), so at the wet bulb h A (T\_air - T\_wb) = R\_c A lambda(T\_wb) = 0.8648 kW: the heat in equals the latent load, and the temperature leaves the wet bulb continuously at X\_c. X <= X\_c: C(X) dT\_s/dt = h A (T\_air - T\_s) - R A lambda(T\_s), C = m\_s (cp\_s + X cp\_w), cp\_s = 1243 J/(kg K) (the record's 425.5 J/(mol K)), cp\_w = 4191. The time scale of the warm-up is C/(h A) = 100-130 s, short against the falling-rate time tau = 595 s, so T\_s follows the air with a lag and T\_air - T\_s decays with the drying curve (T\_air - T\_s \textasciitilde{} 1.4 K at 3000 s). LIMIT, by hand: X $\to$ X\_eq makes R $\to$ 0 and T\_s $\to$ T\_air (333.15 K); at 3000 s X = 0.01963 against X\_eq = 0.01586, 3.6 \% of the free moisture left, and T\_s = 331.78 K, 1.37 K below the air. INDEPENDENT CHECK: the same two equations integrated outside the engine (a ten-line RK4 with the printed constants and the Watson latent heat) end at T\_s = 331.7796 K against the engine's 331.7787 K, X = 0.019628 against 0.019631. THE ENERGY LEDGER closes: E\_conv = 1470.48 kJ in from the air, E\_out = 1387.63 kJ out with the vapour (latent + its sensible), and the solid stores (m\_s cp\_s + m\_w cp\_w)(T\_s - T\_wb) = 2650.7 J/K $\ast$ 31.26 K = 82.85 kJ: 1470.48 - 1387.63 - 82.85 = 0 to 1.1e-9 of E\_conv (KPI solidEnergy\_closure\_rel). The campaign's own energy balance stays unavailable (the air is outside it), as in dryer01. Announced assumptions: R is the characteristic drying curve and does not read T\_s; one temperature for the lumped solid (small Biot); cp at the mean of T\_wb and T\_air; the charge starts at T\_wb.

A tray holding 2.0 kg of dry sucrose wet to X\_0 = 0.30 kg water / kg dry solid, exposed over 0.5 m2 to air of DECLARED, CONSTANT condition: 60 C (333.15 K) and Y = 0.010 kg water / kg dry N2. Nothing about that air is integrated -- it is an ENVIRONMENT, the FixedBedAdsorber's constant-carrier posture, and the water that leaves the solid leaves the campaign across that boundary (reported as the unit's DECLARED material residual, never as a silent leak).

THE HAND ESTIMATE, computed from this case's own curated records (water Antoine + Watson Hvap + the two ideal-gas Cp polynomials, the case-local GAB isotherm on sucrose) -- the run must reproduce it:

air water activity p\_v = P Y/(Y + Mv/Mc) = 1551 Pa a\_w = p\_v / Psat(333.15 K) = 0.0774 GAB equilibrium K a\_w = 0.0743 X\_eq = Xm C Ka /((1-Ka)(1-Ka+C Ka)) = 0.01580 kg/kg wet bulb (Lewis = 1) T\_wb = 300.53 K (27.4 C, a 32.6 K depression) driving force Y\_sat(T\_wb) - Y = 0.014011 kg/kg constant-rate flux R\_c = k\_Y dY = 0.05 $\ast$ 0.014011 = 7.0057e-4 kg/(m2 s) R\_c A = 3.5028e-4 kg/s latent load R\_c A lambda(T\_wb) = 0.8645 kW (the AIR's, not the campaign's -- KPI only, never ledgered) break time t\_c = m\_s (X\_0 - X\_c)/(R\_c A) = 1027.7 s falling-rate constant tau = m\_s (X\_c - X\_eq)/(R\_c A) = 595.0 s end of run X(3000 s) = X\_eq + (X\_c - X\_eq) exp(-(3000-t\_c)/tau) = 0.019584 kg/kg

What to look at in the output: trajectory.csv columns dryer.X\_kg\_kg and dryer.R\_kg\_m2\_s: the drying curve and the rate curve, engine columns, not a post-processing. R is FLAT at 7.0057e-4 until X falls through X\_c = 0.12, then decays linearly in (X - X\_eq) -- the textbook break, visible as a corner. The falling rate is DECLARED as the characteristic drying curve f = Phi\textasciicircum{}n with n = 1, which IS the linear law; the drying EduTool varies n (n < 1 reaches X\_eq in finite time, n > 1 an algebraic tail slower than the exponential), and dryer02\_characteristic\_curves runs four shapes side by side. the KPI t\_critical is OBSERVED (the interpolated crossing between two accepted states), not the analytic 1027.7 s above; they agree because dX/dt is exactly constant on that leg. THE CAMPAIGN MASS BALANCE CLOSES BECAUSE THE DRIED WATER IS AN OUTLET. The solid is charged with 2.600000 kg and ends with 2.039168 kg; the 0.560832 kg of difference is the DRIED WATER. It left the solid and went into the air, and the campaign books it as an external OUTLET STREAM of the dryer (the ledger record `dryer $\to$ (external boundary)`, "evaporated water leaving the solid into the drying air"): initial = final + external out, mass and element closure at round-off. The air stays a DECLARED ENVIRONMENT (a temperature and a humidity, not a vessel this campaign owns), but matter crossing the campaign boundary is a stream whatever receives it. Until 2026-11-09 this unit left the water off the books and the balance read a 21.6 \% shortfall that only the log explained; a violated balance on screen is a defect, so it is a stream now (DEV.md 4c C258).

the campaign energy balance is UNAVAILABLE, and says why in one sentence: the drying heat comes from air this campaign never owned. The number is not hidden -- it is the latentDuty\_kW KPI.

Method hypotheses, each announced at run time: Lewis = 1 in the adiabatic-saturation equation; the humid heat evaluated at the air's own temperature; the surface held at T\_wb while X > X\_c; the LINEAR falling-rate curve (n = 1; a modelling CHOICE, not physics, and its declared source says so); the solid's own temperature NOT integrated.

\end{tutnarrative}
\begin{tutgolden}
\goldrow{kpi}{campaign}{element\_C\_closure\_rel}{0}
\goldrow{kpi}{campaign}{element\_H\_closure\_rel}{0}
\goldrow{kpi}{campaign}{element\_O\_closure\_rel}{0}
\goldrow{kpi}{campaign}{element\_worst\_closure\_rel}{0}
\goldrow{kpi}{campaign}{energy\_balance\_available}{0}
\goldrow{kpi}{campaign}{energy\_records\_invalid}{0}
\goldrow{kpi}{campaign}{energy\_records\_valid}{0}
\goldrow{kpi}{campaign}{mass\_closure\_rel}{0}
\goldrow{kpi}{campaign}{mass\_kg\_external\_out}{0.560738660881}
\goldrow{kpi}{campaign}{mass\_kg\_final}{2.03926135287}
\goldrow{kpi}{campaign}{mass\_kg\_initial}{2.60000001375}
\goldrow{kpi}{campaign}{mass\_residual\_kg}{0}
\goldrow{kpi}{campaign}{moles\_kmol\_external\_out}{0.031126209319}
\goldrow{kpi}{campaign}{moles\_kmol\_final}{0.00802224868102}
\goldmore{29}
\end{tutgolden}

\subsection{electrodialysis}
\grouptheory{electrodialysis}
\clearpage
\tutentry{edbatch01\_constant\_current}{tutorials/batch/electrodialysis/edbatch01\_constant\_current}
\begin{tutscheme}
\begin{tikzpicture}[>={Stealth[length=2mm]}, every node/.style={font=\scriptsize}]
  \node[draw, rounded corners=2pt, align=center, minimum width=1.9cm, minimum height=0.85cm, fill=black!4] (rig) at (0.0,-0.0) {\textbf{rig}\\[-1pt]{\tiny batchElectrodialysis}};
\end{tikzpicture}
\end{tutscheme}
\tutwhat{Batch recirculating electrodialysis at constant current: the limiting current falls with the depleting diluate until i crosses it}
\begin{tutnarrative}
a batch electrodialysis rig run at CONSTANT CURRENT, and the limiting current falling under it until they cross.

THE OPERATION. Five litres of 0.1 mol/kg brine in the diluate tank, five litres of 0.5 mol/kg in the concentrate tank, both recirculated through the bench stack of `ed03\_stack\_record` (Eurodia EUR2C-7P18, seven cell pairs, 0.020 m2 each) at the paper's own 175 L/h. The rectifier is held at 3.5 A, which is 175 A/m2 on one cell pair.

WHY IT IS A DIFFERENT EXPERIMENT FROM ed01-ed05. Those are steady: one pass, a fixed diluate composition, so the limiting current is evaluated ONCE. Here the diluate tank DEPLETES, and i\_lim -- which the GeraldesAfonso2010 model predicts from the ion diffusivities and the diluate composition -- falls with it, while i does not move at all. They must cross.

WHAT THE RUN MEASURES (every number below is printed by this run; nothing here is a target):

i 175.00 A/m2, flat, all 2000 s i\_lim 643.04 A/m2 at t = 0 88.26 A/m2 at t = 2000 s THE CROSSING t = 1593 s (i/i\_lim reaches 1) i/i\_lim at the end 1.983 xi\_eff 0.900 $\to$ 0.4539 (the plateau, below) demineralisation 86.27 \% ACTUAL 91.41 \% what the hand calculation predicts the hand calculation's own answer: the tank is EMPTY at t = 2188 s concentrate 0.5 $\to$ 0.5863 mol/kg (x1.173) specific energy 1.946 kWh per m3 of diluate treated material closure 1.4e-15 on mass, 8.5e-14 on the worst element

THE LESSON. Faraday's law at constant current is a STRAIGHT LINE:

n(t) = n(0) - xi I N t / (z F)

and it is what a student computes by hand. The run publishes it as `demin\_ideal`, evaluated from THIS run's own initial current, own initial inventory and own declared xi -- computed, never declared, so the gap cannot have been arranged -- beside `demin\_actual`.

The two agree exactly until t = 1593 s and part after it, and the reason is printed at the instant it happens: i has caught i\_lim. At the limiting current the counter-ion concentration vanishes at the diluate/membrane interface and the film can deliver no more of it, so a larger current is carried by water-splitting products instead of by salt. This case DECLARES that model -- `overLimiting \{ model saltFluxPlateau; \}`, xi\_eff = xi i\_lim/i, which adds no parameter -- so the depletion stops being linear and becomes asymptotic: the straight line hits zero at 2188 s and the tank never does.

WITHOUT the declaration (the default) the current efficiency would stay at xi and the two curves would coincide to the integrator's own precision. That is not a better answer; it is a model with nothing in it that could part from the hand calculation, and the run SAYS so when i crosses i\_lim. `edbatch02\_constant\_voltage` is the same rig with the rectifier held at a voltage instead: there the current DECAYS, the crossing never happens, and the same hand calculation fails for an entirely different reason.

WHAT IS PINNED AND WHAT IS NOT. The golden pins the answer THIS model gives. Nothing here is validated against a measured batch electrodialysis: the stack record reproduces the paper's bench unit and its limiting-current correlation, but no run of that unit in recirculating batch is transcribed anywhere in this tree. `t\_overLimiting\_s` is resolved to the TIME STEP -- it is the first accepted step at which i exceeds i\_lim, so it moves with deltaT (1593 s at deltaT 1, 1592.5 s at deltaT 0.25) while everything else is converged to \textasciitilde{}1e-8.

NOT MODELLED, and it matters to how far this run can be trusted: no salt diffuses back from the concentrate, no water is transported with the ions, the current efficiency is a declared constant apart from the plateau above, and the rig is isothermal with the Joule heat removed by no modelled duty (which is why the campaign energy balance reports UNAVAILABLE and says why). The full list is in the unit's own header.

THE OPERATOR PANEL. `system/panelDict` declares what a person at the rectifier would read: six elements (the diluate conductivity, the actual demineralisation against Faraday's straight line, the stack voltage, the over-limiting lamp, the current density against the stack rating, the rig temperature), each with the reason it is there, two limits taken from the stack record and one from the unit's own film model, three alarms and two messages. The engine reads the file AFTER the time loop and reports FACTS about it -- in the log (`[panel]` lines) and in the result's `panel` block, pinned by the golden's `panel` rows: the lamp's limit is crossed at 9 of 41 written instants from t = 1600 s on the write grid (the unit's own step-resolved crossing is 1593 s); the current density and the temperature never move at constant current. It reports NO GRADE: whether this panel tells a person what the rig is doing is the exercise. It cannot change a number of the run above, which is held by the unchanged golden rows.

\end{tutnarrative}
\begin{tutgolden}
\goldrow{kpi}{campaign}{element\_Cl\_closure\_rel}{8.44232091619e-14}
\goldrow{kpi}{campaign}{element\_H\_closure\_rel}{0}
\goldrow{kpi}{campaign}{element\_Na\_closure\_rel}{8.44232091619e-14}
\goldrow{kpi}{campaign}{element\_O\_closure\_rel}{0}
\goldrow{kpi}{campaign}{element\_worst\_closure\_rel}{8.44232091619e-14}
\goldrow{kpi}{campaign}{energy\_balance\_available}{0}
\goldrow{kpi}{campaign}{energy\_records\_invalid}{0}
\goldrow{kpi}{campaign}{energy\_records\_valid}{0}
\goldrow{kpi}{campaign}{mass\_closure\_rel}{1.39662189668e-15}
\goldrow{kpi}{campaign}{mass\_kg\_external\_out}{0}
\goldrow{kpi}{campaign}{mass\_kg\_final}{10.1751624752}
\goldrow{kpi}{campaign}{mass\_kg\_initial}{10.1751624752}
\goldrow{kpi}{campaign}{mass\_residual\_kg}{1.42108547152e-14}
\goldrow{kpi}{campaign}{moles\_kmol\_external\_out}{0}
\goldmore{119}
\end{tutgolden}

\clearpage
\tutentry{edbatch02\_constant\_voltage}{tutorials/batch/electrodialysis/edbatch02\_constant\_voltage}
\begin{tutscheme}
\begin{tikzpicture}[>={Stealth[length=2mm]}, every node/.style={font=\scriptsize}]
  \node[draw, rounded corners=2pt, align=center, minimum width=1.9cm, minimum height=0.85cm, fill=black!4] (rig) at (0.0,-0.0) {\textbf{rig}\\[-1pt]{\tiny batchElectrodialysis}};
\end{tikzpicture}
\end{tutscheme}
\tutwhat{Batch recirculating electrodialysis at constant voltage: the current decays as the diluate loses conductivity, and never reaches the limiting current}
\begin{tutnarrative}
THE SAME RIG, one knob different, and the hand calculation failing for the opposite reason.

Everything is identical to `edbatch01\_constant\_current` -- the same stack, the same two tanks, the same recirculation flow, the same xi, the same electrode term -- except that the rectifier is held at 4 V instead of 3.5 A. The difference between the two runs IS the difference between the two ways of driving a batch electrodialysis.

WHAT THE RUN MEASURES (printed by this run; nothing here is a target):

I 4.365 A at t = 0 $\to$ 2.018 A at t = 1500 s i/i\_lim 0.3394 $\to$ 0.4086 NEVER reaches 1 demineralisation 61.59 \% ACTUAL 85.51 \% what the hand calculation predicts the hand calculation's own answer: the tank is EMPTY at t = 1754 s diluate 0.1 $\to$ 0.03841 mol/kg kappa (diluate) 1.262 $\to$ 0.4846 S/m E\_mem (per pair) 0.0794 $\to$ 0.1313 V specific energy 1.048 kWh per m3 of diluate treated material closure machine level, as in edbatch01

WHY THE CURRENT FALLS. The stack equation is

U = N ( E\_mem + I R\_pair ) + E\_electrodes

and BOTH terms inside the bracket move against the current as the run proceeds. The diluate loses conductivity, so R\_pair rises (0.0636 $\to$ 0.1119 Ohm per cell pair); and the concentration RATIO between the two tanks grows, so the Nernst back-EMF rises with it. At fixed U the current has nowhere to go but down, and it more than halves.

WHY IT NEVER CROSSES i\_lim. i\_lim falls as the diluate depletes -- exactly as in edbatch01 -- but so does i, and here they fall together: the ratio ends at 0.409 having started at 0.339. That is the reason a real bench rig is run this way when the membranes matter: constant voltage is self-limiting, and the price is time. This case declares NO `overLimiting \{\}` model at all, and the run announces that default; it is irrelevant here precisely because the crossing never happens.

THE LESSON, and it is the same lesson as edbatch01 arriving by a different road. `demin\_ideal` is Faraday's law at the INITIAL current -- a straight line, computed from this run's own I(0), own inventory and own xi. It is 24 percentage points optimistic at 1500 s, and not one of those points has anything to do with the limiting current: the current simply is not what it was. Read `I\_A` in trajectory.csv beside `deminActual` and `deminIdeal` and the whole of the gap is in that first column.

WHAT IS PINNED AND WHAT IS NOT: as edbatch01. The golden pins the answer this model gives; no measured batch electrodialysis is reproduced anywhere in this tree.

\end{tutnarrative}
\begin{tutgolden}
\goldrow{kpi}{campaign}{element\_Cl\_closure\_rel}{1.30104260695e-15}
\goldrow{kpi}{campaign}{element\_H\_closure\_rel}{0}
\goldrow{kpi}{campaign}{element\_Na\_closure\_rel}{1.30104260695e-15}
\goldrow{kpi}{campaign}{element\_O\_closure\_rel}{0}
\goldrow{kpi}{campaign}{element\_worst\_closure\_rel}{1.30104260695e-15}
\goldrow{kpi}{campaign}{energy\_balance\_available}{0}
\goldrow{kpi}{campaign}{energy\_records\_invalid}{0}
\goldrow{kpi}{campaign}{energy\_records\_valid}{0}
\goldrow{kpi}{campaign}{mass\_closure\_rel}{0}
\goldrow{kpi}{campaign}{mass\_kg\_external\_out}{0}
\goldrow{kpi}{campaign}{mass\_kg\_final}{10.1751624752}
\goldrow{kpi}{campaign}{mass\_kg\_initial}{10.1751624752}
\goldrow{kpi}{campaign}{mass\_residual\_kg}{0}
\goldrow{kpi}{campaign}{moles\_kmol\_external\_out}{0}
\goldmore{46}
\end{tutgolden}

\clearpage
\tutentry{edbatch03\_multiionic\_split}{tutorials/batch/electrodialysis/edbatch03\_multiionic\_split}
\begin{tutscheme}
\begin{tikzpicture}[>={Stealth[length=2mm]}, every node/.style={font=\scriptsize}]
  \node[draw, rounded corners=2pt, align=center, minimum width=1.9cm, minimum height=0.85cm, fill=black!4] (rig) at (0.0,-0.0) {\textbf{rig}\\[-1pt]{\tiny batchElectrodialysis}};
\end{tikzpicture}
\end{tutscheme}
\tutwhat{Batch electrodialysis of a mixed NaCl + K2SO4 brine: two counter-ions on each membrane share its current by z D c}
\begin{tutnarrative}
a batch electrodialysis rig on a MIXED brine, and two counter-ions on each membrane sharing its current.

THE OPERATION. The rig of edbatch01 (Eurodia EUR2C-7P18, seven cell pairs, 175 L/h recirculation), but the diluate tank holds 0.05 mol/kg NaCl plus 0.025 mol/kg K2SO4 -- two cations (Na+, K+) on the cation-exchange membrane and two anions (Cl-, SO4 2-) on the anion-exchange one. The rectifier is held at 2.0 A (100 A/m2 on one cell pair).

THE RULE (Vitor's ruling C35, 2026-10-05). Until that day this unit REFUSED a diluate with more than one cation or one anion, for want of a rule dividing a membrane's current between its counter-ions. The ruling shares it by

t\_k = |z\_k| D\_k c\_k / SUM\_j |z\_j| D\_j c\_j over that membrane's counter-ions

-- the transport numbers of Geraldes \& Afonso, J. Membr. Sci. 360 (2010) 499-508, Eqs. (12)/(13), applied at every current and not only at the limit -- with D the ions' curated D0 and NO new parameter. Each ion then leaves at t\_k xi I N / (|z\_k| F) mol/s, so each membrane moves exactly xi I N / F equivalents and both tanks stay electroneutral.

WHAT A STUDENT SHOULD SEE. Divide that rate by the ion's own inventory:

d ln n\_k / dt = - D\_k xi I N / ( F SUM\_j |z\_j| D\_j n\_j )

-- the FRACTIONAL depletion rate of a counter-ion is proportional to its diffusivity and to nothing else about it, so on a log scale the counter-ions of one membrane fall along straight lines whose slopes stand in the ratio of their D0. K+ (1.96e-9 m2/s) is depleted 1.47 times faster than Na+ (1.33e-9) in log terms; Cl- (2.03e-9) 1.92 times faster than SO4 2- (1.06e-9), the charge of the sulfate making no difference to its FRACTIONAL rate (it carries twice the current per mole and so moves half the moles per coulomb, exactly compensating its doubled share). The run publishes `demin\_<ion>` for every ion and `t\_transfer\_final\_<ion>`, the shares at the end, and its trajectory carries every molality.

WHAT THE RUN MEASURES (every number is printed by this run; the golden pins them): i 100.00 A/m2, flat, all 2400 s i\_lim 696.83 A/m2 at t = 0, 244.39 at t = 2400 s (i/i\_lim 0.144 $\to$ 0.409: never crossed) removed at t = 2400 s K+ 69.77 \% Na+ 55.60 \% (CEM) Cl- 74.43 \% SO4 2- 50.94 \% (AEM) share of the current Na+ 0.404 $\to$ 0.499 of the CEM's Cl- 0.657 $\to$ 0.500 of the AEM's material closure 0 on mass, 1.4e-15 on the worst element

READ THE SHARES BACKWARDS AND THE LESSON IS THERE. At t = 0 the two cations are equally concentrated, so K+, the faster, takes 0.596 of the cation current; it therefore depletes faster, its c falls, and the share drifts toward Na+. The hand calculation `demin\_ideal` (Faraday's law at the INITIAL current AND the reference cation's INITIAL share, computed from the run) predicts 50.68 \% of the sodium removed; the run removes 55.60 \%, because sodium's share of the current GROWS as potassium is stripped -- the opposite sign from edbatch01's gap, and for a different reason. And ln(n/n0)/D0 comes out as one number for K+ and Na+, and one number for Cl- and SO4 2-, to the integrator's precision.

WHAT IS PINNED AND WHAT IS NOT. The golden pins the answer THIS model gives. The share z D c is a RULING, not a measurement: real ion-exchange membranes add a preference between counter-ions of their own (a Donnan preference for divalent ions, or a monovalent-selective layer) which no `kind IEM` record in this tree carries and this unit does not model. The brine is an author-set assumption; nothing here is validated against a measured rig. The membrane potential of a membrane with two counter-ions is taken as the MEAN of their own Nernst potentials (announced). Back- diffusion, water transport and every thermal effect are not modelled -- the unit's own header lists them.

\end{tutnarrative}
\begin{tutgolden}
\goldrow{kpi}{campaign}{element\_Cl\_closure\_rel}{2.45752492414e-15}
\goldrow{kpi}{campaign}{element\_H\_closure\_rel}{0}
\goldrow{kpi}{campaign}{element\_K\_closure\_rel}{1.44560289655e-15}
\goldrow{kpi}{campaign}{element\_Na\_closure\_rel}{1.15648231724e-15}
\goldrow{kpi}{campaign}{element\_O\_closure\_rel}{0}
\goldrow{kpi}{campaign}{element\_S\_closure\_rel}{5.7824115862e-16}
\goldrow{kpi}{campaign}{element\_worst\_closure\_rel}{2.45752492414e-15}
\goldrow{kpi}{campaign}{energy\_balance\_available}{0}
\goldrow{kpi}{campaign}{energy\_records\_invalid}{0}
\goldrow{kpi}{campaign}{energy\_records\_valid}{0}
\goldrow{kpi}{campaign}{mass\_closure\_rel}{0}
\goldrow{kpi}{campaign}{mass\_kg\_external\_out}{0}
\goldrow{kpi}{campaign}{mass\_kg\_final}{10.2181899752}
\goldrow{kpi}{campaign}{mass\_kg\_initial}{10.2181899752}
\goldmore{58}
\end{tutgolden}

\subsection{membrane}
\grouptheory{membrane}
\clearpage
\tutentry{diafilter01\_nf\_desalting}{tutorials/batch/membrane/diafilter01\_nf\_desalting}
\begin{tutscheme}
\begin{tikzpicture}[>={Stealth[length=2mm]}, every node/.style={font=\scriptsize}]
  \node[draw, rounded corners=2pt, align=center, minimum width=1.9cm, minimum height=0.85cm, fill=black!4] (retentate) at (0.0,-0.0) {\textbf{retentate}\\[-1pt]{\tiny batchDiafilter}};
\end{tikzpicture}
\end{tutscheme}
\tutwhat{Constant-volume NF diafiltration: NaCl washed out of an MgSO4 retentate, with the OBSERVED rejection published beside the constant-R washout law it violates}
\begin{tutnarrative}
constant-volume diafiltration on a loose NF membrane, and the exponential washout law failing in a way you can watch.

THE OPERATION. 10 litres of a brine carrying TWO salts sit behind 5 m2 of NF270 at 10 bar. Solvent is made up as fast as it permeates, so the volume is held and the run is measured in DIAVOLUMES -- how many vessel volumes of wash water have passed through. NaCl is poorly rejected by this membrane and leaves; MgSO4 is well rejected and stays. That is the whole point of a nanofiltration desalting step, and it is the same arithmetic whether the retained species is a divalent salt, a sugar or a protein.

THE LESSON. Derive the washout by hand. Over a constant volume V, with the permeate carrying c\_p = (1 - R) c,

V dc/dt = -Q\_p (1 - R) c , N = Q\_p t / V => c/c\_0 = exp( -(1 - R) N )

which is a straight line on a log axis against N. The derivation needs R to be CONSTANT. It is not:

as NaCl leaves, the osmotic pressure of the retentate falls, so the driving force rises and so does the water flux; a higher flux sweeps more solute to the wall, so the polarisation factor rises and the WALL concentration is no longer the bulk one; the observed rejection is measured against the bulk, so it moves even when the membrane's own permeability has not changed.

The run therefore publishes R\_obs for each solute at every instant, beside the diavolumes, and -- because this is constant-volume mode -- two KPIs per solute: `washoutActual\_<name>` (what the vessel did) and `washoutIdeal\_<name>` (what exp(-(1-R)N) predicts using THIS run's own initial R). The idealisation is computed from the run, not declared, so the gap between them cannot have been arranged.

Read them together. Where the two agree, the hand derivation was safe. Where they part, the R\_obs column says which assumption failed and when.

WHAT IS PINNED AND WHAT IS NOT. The golden pins the answer this model gives; NOTHING here is validated against a measured diafiltration. The NF270 permeabilities are the case-fitted pedagogical set the record declares (`provenance fittedToCase`, fitted against membrane02\_NF\_sugar), so the rejections are representative of a loose NF membrane and are not a claim about a particular element. Fouling is NOT modelled -- the flux in a real rig would decline with time and this one does not.

\end{tutnarrative}
\begin{tutgolden}
\goldrow{kpi}{campaign}{element\_Cl\_closure\_rel}{0}
\goldrow{kpi}{campaign}{element\_H\_closure\_rel}{4.9535921512e-16}
\goldrow{kpi}{campaign}{element\_Mg\_closure\_rel}{0}
\goldrow{kpi}{campaign}{element\_Na\_closure\_rel}{0}
\goldrow{kpi}{campaign}{element\_O\_closure\_rel}{4.95248741147e-16}
\goldrow{kpi}{campaign}{element\_S\_closure\_rel}{0}
\goldrow{kpi}{campaign}{element\_worst\_closure\_rel}{4.9535921512e-16}
\goldrow{kpi}{campaign}{energy\_balance\_available}{0}
\goldrow{kpi}{campaign}{energy\_records\_invalid}{0}
\goldrow{kpi}{campaign}{energy\_records\_valid}{0}
\goldrow{kpi}{campaign}{mass\_closure\_rel}{4.39709286757e-16}
\goldrow{kpi}{campaign}{mass\_kg\_amend\_in}{54.6016152827}
\goldrow{kpi}{campaign}{mass\_kg\_amend\_out}{0}
\goldrow{kpi}{campaign}{mass\_kg\_external\_out}{54.6016152827}
\goldmore{52}
\end{tutgolden}

\clearpage
\tutentry{diafilter02\_fouling\_decline}{tutorials/batch/membrane/diafilter02\_fouling\_decline}
\begin{tutscheme}
\begin{tikzpicture}[>={Stealth[length=2mm]}, every node/.style={font=\scriptsize}]
  \node[draw, rounded corners=2pt, align=center, minimum width=1.9cm, minimum height=0.85cm, fill=black!4] (retentate) at (0.0,-0.0) {\textbf{retentate}\\[-1pt]{\tiny batchDiafilter}};
\end{tikzpicture}
\end{tutscheme}
\tutwhat{The same diafiltration with a DECLARED Hermia cake-law decline: what fouling costs a washout schedule}
\begin{tutnarrative}
the same diafiltration, with the membrane losing permeance as it works.

WHY THIS CASE EXISTS. diafilter01's flux RISES during the run (233 $\to$ 256 LMH): the salt leaves, the osmotic pressure falls, the driving force grows. That is the physics the model has. It is not what a rig does. On real equipment the flux falls, because the membrane fouls, and every schedule a student will ever write is sized on that decline.

THE MODEL. Resistance in series on the permeance the unit hands the transport law:

1 / A\_eff = 1 / A\_w + r\_f(v), v = V\_permeated / A

with r\_f growing by one of HERMIA'S BLOCKING LAWS (Hermia, Trans. IChemE 60:183-187, 1982). Here `cake`: r\_f = k v, the deposited-layer form. Applied where the transport CONTEXT is assembled, so solution-diffusion and DSPM-DE are served alike and neither is modified.

THE CONSTANT IN THIS CASE IS HYPOTHETICAL AND IS NOT A FITTED VALUE FOR THIS SYSTEM. $\ast$$\ast$ A brine of NaCl and MgSO4 is fully dissolved; there is no suspended matter to deposit, so `cake` is not a mechanism this feed can honestly claim. What this case demonstrates is the CONSEQUENCE of a declared decline on a diafiltration schedule -- how many more minutes the same washout costs -- on a system whose other numbers are already pinned by diafilter01, so the difference between the two runs is the fouling and nothing else. It is a STRUCTURAL witness, in the same sense as edwards01\_sour\_water\_activity: no number in it is validated against a measured fouling experiment, and the tree carries no fouling data at all.

A fouling constant belongs to a RIG AND A MATERIAL, not to a membrane, so it is declared here in the case and no catalogue record is touched.

WHAT TO READ. `A\_eff\_over\_A\_w` in the trajectory is the permeance the membrane still has, as a fraction of clean. Put diafilter01 and diafilter02 side by side: same charge, same area, same pressure, and the diavolume count at t = 160 s is the whole lesson.

\end{tutnarrative}
\begin{tutgolden}
\goldrow{kpi}{campaign}{element\_Cl\_closure\_rel}{0}
\goldrow{kpi}{campaign}{element\_H\_closure\_rel}{5.55943179339e-16}
\goldrow{kpi}{campaign}{element\_Mg\_closure\_rel}{0}
\goldrow{kpi}{campaign}{element\_Na\_closure\_rel}{0}
\goldrow{kpi}{campaign}{element\_O\_closure\_rel}{5.55804034071e-16}
\goldrow{kpi}{campaign}{element\_S\_closure\_rel}{0}
\goldrow{kpi}{campaign}{element\_worst\_closure\_rel}{5.55943179339e-16}
\goldrow{kpi}{campaign}{energy\_balance\_available}{0}
\goldrow{kpi}{campaign}{energy\_records\_invalid}{0}
\goldrow{kpi}{campaign}{energy\_records\_valid}{0}
\goldrow{kpi}{campaign}{mass\_closure\_rel}{4.93453705585e-16}
\goldrow{kpi}{campaign}{mass\_kg\_amend\_in}{47.5616338569}
\goldrow{kpi}{campaign}{mass\_kg\_amend\_out}{0}
\goldrow{kpi}{campaign}{mass\_kg\_external\_out}{47.5616338569}
\goldmore{53}
\end{tutgolden}

\clearpage
\tutentry{diafilter03\_dspmde\_batch}{tutorials/batch/membrane/diafilter03\_dspmde\_batch}
\begin{tutscheme}
\begin{tikzpicture}[>={Stealth[length=2mm]}, every node/.style={font=\scriptsize}]
  \node[draw, rounded corners=2pt, align=center, minimum width=1.9cm, minimum height=0.85cm, fill=black!4] (retentate) at (0.0,-0.0) {\textbf{retentate}\\[-1pt]{\tiny batchDiafilter}};
\end{tikzpicture}
\end{tutscheme}
\tutwhat{Batch NF concentration on the DSPM-DE pore model: the observed rejection per instant and the campaign's time-averaged rejection}
\begin{tutnarrative}
batch CONCENTRATION of a two-salt brine on the DSPM-DE pore model, and the rejection of a whole campaign stated as a TIME average.

THE OPERATION. The charge of diafilter01 (10 litres, 0.02 kmol/m3 each of NaCl and MgSO4) sits behind 1 m2 of NF270 at 10 bar with NO make-up solvent, so the vessel shrinks and both salts concentrate. The membrane is the pore tier of the curated NF270 record (`NF270\_dspmde`: pore radius, A\_k/dx, a constant volumetric charge, a pore dielectric constant), and the transport law is DSPM-DE -- steric, Donnan and dielectric exclusion with the extended Nernst-Planck equations SOLVED across the pore -- asked at every instant of the integration, exactly as the spiral module asks it at every station of its channel. The film is the declared k\_film with the default film theory (announced).

WHAT THE RUN SHOWS. The observed rejection is not a property of the membrane alone; it moves with the solution the membrane sees. As the retentate concentrates the NaCl rejection FALLS (0.691 at the start, 0.554 at the end of this run) while the MgSO4 rejection RISES slightly (0.973 $\to$ 0.979). In DSPM-DE two things push the monovalent salt the same way -- a bulk concentration growing past the fixed charge weakens the Donnan exclusion, and a growing share of retained divalent ions changes what must cross to keep the permeate neutral -- and this run does not separate them. The flux falls with the osmotic pressure (244 $\to$ 116 LMH) and the volume reduction reaches 3.9. Read `R\_obs\_<salt>` in the trajectory against `V\_m3`.

THE REJECTION OF THE CAMPAIGN. A batch rig is usually judged by the permeate it COLLECTED, not by one instant, and two time averages of the rejection are published at the end (DEFINED in BatchDiafilter.cpp, kpis(), and in docs/ai/membrane-research.md).

R\_avg = 1 - c\_permPooled / c\_bMeanPerm, where c\_permPooled = (n\_0 - n) / V\_perm is the pooled permeate (every drop collected, mixed) and c\_bMeanPerm = INTEGRAL c\_b dV\_perm / V\_perm is the retentate's bulk concentration averaged over what permeated. It is a weighted mean of the instantaneous R\_obs (weights c\_b Q\_p), so it lies between the smallest and the largest of them (NaCl 0.635).

R\_avgVsInitialFeed = 1 - c\_permPooled / c\_bInitial compares the pooled permeate with the CHARGE and, because the retentate only concentrates, lies below R\_avg (NaCl 0.467) -- the number a laboratory gets by analysing its collected permeate against the feed it started from, and a reminder that it is not the membrane's rejection. The steady spiral module states the SPATIAL counterpart: its R\_obs is the mixed permeate of the whole channel against the inlet feed.

WHAT IS VERIFIED AND WHAT IS NOT. Verified (check\_diafiltration): the DSPM-DE pore model runs inside the batch vessel at every instant; matter is conserved to machine precision; the two averages reproduce from the published integrals and from the trajectory; R\_avg lies inside the range of R\_obs; the pooled-permeate rejection lies below it, as it must when the retentate concentrates. NOT validated: no number here is compared with a measured batch filtration. The pore parameters are the record's case-fitted pedagogical set (`provenance fittedToCase`, fitted against membrane10\_dspmde\_divalent), so the rejections are those of a representative loose NF membrane, not a claim about a particular element. In a MIXED feed the salt basis carries each salt's cation leg through the pore and the run says so; fouling is not declared; k\_film is declared, not correlated with a stirrer or a channel.

\end{tutnarrative}
\begin{tutgolden}
\goldrow{kpi}{campaign}{element\_Cl\_closure\_rel}{0}
\goldrow{kpi}{campaign}{element\_H\_closure\_rel}{1.51983140987e-12}
\goldrow{kpi}{campaign}{element\_Mg\_closure\_rel}{0}
\goldrow{kpi}{campaign}{element\_Na\_closure\_rel}{0}
\goldrow{kpi}{campaign}{element\_O\_closure\_rel}{1.51764442727e-12}
\goldrow{kpi}{campaign}{element\_S\_closure\_rel}{0}
\goldrow{kpi}{campaign}{element\_worst\_closure\_rel}{1.51983140987e-12}
\goldrow{kpi}{campaign}{energy\_balance\_available}{0}
\goldrow{kpi}{campaign}{energy\_records\_invalid}{0}
\goldrow{kpi}{campaign}{energy\_records\_valid}{0}
\goldrow{kpi}{campaign}{mass\_closure\_rel}{1.51406233018e-12}
\goldrow{kpi}{campaign}{mass\_kg\_external\_out}{7.4717123351}
\goldrow{kpi}{campaign}{mass\_kg\_final}{2.56417330931}
\goldrow{kpi}{campaign}{mass\_kg\_initial}{10.0358856444}
\goldmore{41}
\end{tutgolden}

\subsection{plant}
\grouptheory{plant}
\clearpage
\tutentry{apalutamide01\_tablets}{tutorials/batch/plant/apalutamide01\_tablets}
\begin{tutscheme}
\begin{tikzpicture}[>={Stealth[length=2mm]}, every node/.style={font=\scriptsize}]
  \node[draw, rounded corners=2pt, align=center, minimum width=1.9cm, minimum height=0.85cm, fill=black!4] (acetoneTank) at (0.0,-0.0) {\textbf{acetoneTank}\\[-1pt]{\tiny storageTank}};
  \node[draw, rounded corners=2pt, align=center, minimum width=1.9cm, minimum height=0.85cm, fill=black!4] (solutionPrep) at (2.9,-0.0) {\textbf{solutionPrep}\\[-1pt]{\tiny mixer}};
  \node[draw, rounded corners=2pt, align=center, minimum width=1.9cm, minimum height=0.85cm, fill=black!4] (n2Mix) at (513.3,-0.0) {\textbf{n2Mix}\\[-1pt]{\tiny mixer}};
  \node[draw, rounded corners=2pt, align=center, minimum width=1.9cm, minimum height=0.85cm, fill=black!4] (blower) at (516.2,-0.0) {\textbf{blower}\\[-1pt]{\tiny compressor}};
  \node[draw, rounded corners=2pt, align=center, minimum width=1.9cm, minimum height=0.85cm, fill=black!4] (gasHeater) at (519.1,-0.0) {\textbf{gasHeater}\\[-1pt]{\tiny phaseChanger}};
  \node[draw, rounded corners=2pt, align=center, minimum width=1.9cm, minimum height=0.85cm, fill=black!4] (SD) at (522.0,-0.0) {\textbf{SD}\\[-1pt]{\tiny sprayDryer}};
  \node[draw, rounded corners=2pt, align=center, minimum width=1.9cm, minimum height=0.85cm, fill=black!4] (exhaustCooler) at (524.9,-0.0) {\textbf{exhaustCooler}\\[-1pt]{\tiny phaseChanger}};
  \node[draw, rounded corners=2pt, align=center, minimum width=1.9cm, minimum height=0.85cm, fill=black!4] (solventCondenser) at (527.8,-0.0) {\textbf{solventCondenser}\\[-1pt]{\tiny phaseChanger}};
  \node[draw, rounded corners=2pt, align=center, minimum width=1.9cm, minimum height=0.85cm, fill=black!4] (condensatePot) at (530.7,-0.0) {\textbf{condensatePot}\\[-1pt]{\tiny isothermalFlash}};
  \node[draw, rounded corners=2pt, align=center, minimum width=1.9cm, minimum height=0.85cm, fill=black!4] (purgeSplit) at (533.6,-0.0) {\textbf{purgeSplit}\\[-1pt]{\tiny splitter}};
  \node[draw, rounded corners=2pt, align=center, minimum width=1.9cm, minimum height=0.85cm, fill=black!4] (recoveredTank) at (533.6,-1.7) {\textbf{recoveredTank}\\[-1pt]{\tiny storageTank}};
  \node[draw, rounded corners=2pt, align=center, minimum width=1.9cm, minimum height=0.85cm, fill=black!4] (blender) at (524.9,-1.7) {\textbf{blender}\\[-1pt]{\tiny mixer}};
  \node[draw, rounded corners=2pt, align=center, minimum width=1.9cm, minimum height=0.85cm, fill=black!4] (rollerCompactor) at (527.8,-1.7) {\textbf{rollerCompactor}\\[-1pt]{\tiny mixer}};
  \node[draw, rounded corners=2pt, align=center, minimum width=1.9cm, minimum height=0.85cm, fill=black!4] (mill) at (530.7,-1.7) {\textbf{mill}\\[-1pt]{\tiny splitter}};
  \node[draw, rounded corners=2pt, align=center, minimum width=1.9cm, minimum height=0.85cm, fill=black!4] (screen) at (533.6,-3.4) {\textbf{screen}\\[-1pt]{\tiny splitter}};
  \node[draw, rounded corners=2pt, align=center, minimum width=1.9cm, minimum height=0.85cm, fill=black!4] (finalBlender) at (536.5,-0.0) {\textbf{finalBlender}\\[-1pt]{\tiny mixer}};
  \node[draw, rounded corners=2pt, align=center, minimum width=1.9cm, minimum height=0.85cm, fill=black!4] (tabletPress) at (539.4,-0.0) {\textbf{tabletPress}\\[-1pt]{\tiny splitter}};
  \node[draw, rounded corners=2pt, align=center, minimum width=1.9cm, minimum height=0.85cm, fill=black!4] (coatingSuspensionPrep) at (0.0,-1.7) {\textbf{coatingSuspensionPrep}\\[-1pt]{\tiny mixer}};
  \node[draw, rounded corners=2pt, align=center, minimum width=1.9cm, minimum height=0.85cm, fill=black!4] (airHeater) at (0.0,-3.4) {\textbf{airHeater}\\[-1pt]{\tiny phaseChanger}};
  \node[draw, rounded corners=2pt, align=center, minimum width=1.9cm, minimum height=0.85cm, fill=black!4] (coatingPan) at (542.3,-0.0) {\textbf{coatingPan}\\[-1pt]{\tiny mixer}};
  \node[draw, rounded corners=2pt, align=center, minimum width=1.9cm, minimum height=0.85cm, fill=black!4] (coatingDryer) at (545.2,-0.0) {\textbf{coatingDryer}\\[-1pt]{\tiny solidDryer}};
  \node[draw, rounded corners=2pt, align=center, minimum width=1.9cm, minimum height=0.85cm, fill=black!4] (packagingLine) at (548.1,-0.0) {\textbf{packagingLine}\\[-1pt]{\tiny splitter}};
  \draw[->] ($(acetoneTank.west)+(-1.05,0.00)$) -- ($(acetoneTank.west)+(0,0.00)$) node[at start, above, font=\tiny, anchor=south west] {acetoneFeed};
  \draw[->] ($(solutionPrep.west)+(-1.05,-0.17)$) -- ($(solutionPrep.west)+(0,-0.17)$) node[at start, below, font=\tiny, anchor=north west] {apiFeed};
  \draw[->] ($(solutionPrep.west)+(-1.05,0.17)$) -- ($(solutionPrep.west)+(0,0.17)$) node[at start, above, font=\tiny, anchor=south west] {hpmcasFeed};
  \draw[->] (acetoneTank) -- (solutionPrep) node[midway, above, font=\tiny] {acetoneToPrep};
  \draw[->] ($(n2Mix.west)+(-1.05,0.00)$) -- ($(n2Mix.west)+(0,0.00)$) node[at start, above, font=\tiny, anchor=south west] {n2Makeup};
  \draw[->, dashed] (purgeSplit.south) .. controls +(0,-1.0) and +(0,-1.0) .. (n2Mix.south) node[midway, below, font=\tiny] {gasRecycle};
  \draw[->] (n2Mix) -- (blower) node[midway, above, font=\tiny] {loopGas};
  \draw[->] (blower) -- (gasHeater) node[midway, above, font=\tiny] {blowerOut};
  \draw[->] (solutionPrep) -- (SD) node[midway, above, font=\tiny] {sprayFeed};
  \draw[->] (gasHeater) -- (SD) node[midway, above, font=\tiny] {dryingGas};
  \draw[->] (SD) -- (exhaustCooler) node[midway, above, font=\tiny] {exhaustGas};
  \draw[->] (exhaustCooler) -- (solventCondenser) node[midway, above, font=\tiny] {cooledExhaust};
  \draw[->] (solventCondenser) -- (condensatePot) node[midway, above, font=\tiny] {chilledExhaust};
  \draw[->] (condensatePot) -- (purgeSplit) node[midway, above, font=\tiny] {condenserGas};
  \draw[->] (condensatePot) -- (recoveredTank) node[midway, above, font=\tiny] {recoveredSolvent};
  \draw[->] (SD) -- (blender) node[midway, above, font=\tiny] {asdPowder};
  \draw[->] ($(blender.west)+(-1.05,-0.34)$) -- ($(blender.west)+(0,-0.34)$) node[at start, below, font=\tiny, anchor=north west] {mccFeed};
  \draw[->] ($(blender.west)+(-1.05,0.00)$) -- ($(blender.west)+(0,0.00)$) node[at start, above, font=\tiny, anchor=south west] {ccsIntraFeed};
  \draw[->] ($(blender.west)+(-1.05,0.34)$) -- ($(blender.west)+(0,0.34)$) node[at start, above, font=\tiny, anchor=south west] {silicaIntraFeed};
  \draw[->] (blender) -- (rollerCompactor) node[midway, above, font=\tiny] {intraBlend};
  \draw[->, dashed] (screen.south) .. controls +(0,-1.0) and +(0,-1.0) .. (rollerCompactor.south) node[midway, below, font=\tiny] {finesRecycle};
  \draw[->] (rollerCompactor) -- (mill) node[midway, above, font=\tiny] {ribbons};
  \draw[->] (mill) -- (screen) node[midway, above, font=\tiny] {milledGranules};
  \draw[->] (screen) -- (finalBlender) node[midway, above, font=\tiny] {granules};
  \draw[->] ($(finalBlender.west)+(-1.05,-0.51)$) -- ($(finalBlender.west)+(0,-0.51)$) node[at start, below, font=\tiny, anchor=north west] {smccFeed};
  \draw[->] ($(finalBlender.west)+(-1.05,-0.17)$) -- ($(finalBlender.west)+(0,-0.17)$) node[at start, below, font=\tiny, anchor=north west] {ccsExtraFeed};
  \draw[->] ($(finalBlender.west)+(-1.05,0.17)$) -- ($(finalBlender.west)+(0,0.17)$) node[at start, above, font=\tiny, anchor=south west] {silicaExtraFeed};
  \draw[->] ($(finalBlender.west)+(-1.05,0.51)$) -- ($(finalBlender.west)+(0,0.51)$) node[at start, above, font=\tiny, anchor=south west] {mgstExtraFeed};
  \draw[->] (finalBlender) -- (tabletPress) node[midway, above, font=\tiny] {finalBlend};
  \draw[->] ($(coatingSuspensionPrep.west)+(-1.05,-0.85)$) -- ($(coatingSuspensionPrep.west)+(0,-0.85)$) node[at start, below, font=\tiny, anchor=north west] {coatWater};
  \draw[->] ($(coatingSuspensionPrep.west)+(-1.05,-0.51)$) -- ($(coatingSuspensionPrep.west)+(0,-0.51)$) node[at start, below, font=\tiny, anchor=north west] {pvaFeed};
  \draw[->] ($(coatingSuspensionPrep.west)+(-1.05,-0.17)$) -- ($(coatingSuspensionPrep.west)+(0,-0.17)$) node[at start, below, font=\tiny, anchor=north west] {pegFeed};
  \draw[->] ($(coatingSuspensionPrep.west)+(-1.05,0.17)$) -- ($(coatingSuspensionPrep.west)+(0,0.17)$) node[at start, above, font=\tiny, anchor=south west] {talcFeed};
  \draw[->] ($(coatingSuspensionPrep.west)+(-1.05,0.51)$) -- ($(coatingSuspensionPrep.west)+(0,0.51)$) node[at start, above, font=\tiny, anchor=south west] {tio2Feed};
  \draw[->] ($(coatingSuspensionPrep.west)+(-1.05,0.85)$) -- ($(coatingSuspensionPrep.west)+(0,0.85)$) node[at start, above, font=\tiny, anchor=south west] {feOxFeed};
  \draw[->] ($(airHeater.west)+(-1.05,0.00)$) -- ($(airHeater.west)+(0,0.00)$) node[at start, above, font=\tiny, anchor=south west] {coatingAirIn};
  \draw[->] (tabletPress) -- (coatingPan) node[midway, above, font=\tiny] {tabletCores};
  \draw[->] (coatingSuspensionPrep) -- (coatingPan) node[midway, above, font=\tiny] {coatingSuspension};
  \draw[->] (coatingPan) -- (coatingDryer) node[midway, above, font=\tiny] {wetCoatedTablets};
  \draw[->] (airHeater) -- (coatingDryer) node[midway, above, font=\tiny] {coatingAir};
  \draw[->] (coatingDryer) -- (packagingLine) node[midway, above, font=\tiny] {coatedTablets};
  \draw[->] ($(purgeSplit.east)+(0,0.00)$) -- ($(purgeSplit.east)+(1.05,0.00)$) node[at end, above, font=\tiny, anchor=south east] {purgeGas};
  \draw[->] ($(recoveredTank.east)+(0,0.00)$) -- ($(recoveredTank.east)+(1.05,0.00)$) node[at end, above, font=\tiny, anchor=south east] {recoveredAcetone};
  \draw[->] ($(tabletPress.east)+(0,0.00)$) -- ($(tabletPress.east)+(1.05,0.00)$) node[at end, above, font=\tiny, anchor=south east] {tabletRejects};
  \draw[->] ($(coatingDryer.east)+(0,0.00)$) -- ($(coatingDryer.east)+(1.05,0.00)$) node[at end, above, font=\tiny, anchor=south east] {coatingExhaust};
  \draw[->] ($(packagingLine.east)+(0,0.00)$) -- ($(packagingLine.east)+(1.05,0.00)$) node[at end, above, font=\tiny, anchor=south east] {packagedTablets};
\end{tikzpicture}
\end{tutscheme}
\tutwhat{APALUTAMIDE 60 mg FILM-COATED TABLETS, a drug-product plant (the API is BOUGHT), a generic entrant at loss of exclusivity, on a SOURCED design basis (docs/design/apalutamide-design-basis-dossier.md): the originator's patent formula (US 12,303,493 B2: a 1:3 apalutamide:HPMCAS-LG dispersion, a 700 mg core, a 3 \% Opadry II coat), the dispersion spray-dried from 5 wt \% acetone (the patent's acetone example) through a pressure nozzle in a closed nitrogen loop (blower, heater, spray dryer, exhaust cooler, -10 C solvent condenser, purge and N2 make-up), the recovered acetone to a tank for QC; dry granulation (blending, roller compaction, milling, a screen with a fines recycle), final blend, tableting with rejects, film coating (an aqueous coat sprayed on and its water dried off by hot air) and packaging. The steady flowsheet is the material balance per unit of product (15 625 tablets/h pressed, the patent's PSD-3 dryer at 75 kg/h of solution, is a basis); a BATCH PLAN (system/postDict) sizes each step for a 400 kg lot, counts 25 lots a year for 14.0 million tablets (25 \% of the estimated European market), and finds the pace of each train -- the spray dryer on round-the-clock campaigns, the solid-dose suite lot by lot on two shifts. Blending, compaction, milling, screening, tableting and packaging are MASS-BALANCE PLACEHOLDERS (no equipment physics). Every number the dossier could not source is an author-set assumption, said so where it is declared (README).}
\begin{tutnarrative}
a drug-product plant: apalutamide 60 mg film-coated tablets. The API is BOUGHT; everything else is made here (DEV.md 4c C97, phase A). A BASE CASE: it runs, it is NOT optimised, and every number a group would change is a plain key, said to be an assumption where it is declared. The README is the narrative; this file is the topology.

THE BASIS (README section 2; the design-basis dossier, docs/design/apalutamide-design-basis-dossier.md): 15 625 tablets per hour PRESSED, the flowsheet basis, of the originator's patent formula (US 12,303,493 B2, Examples 3.1-3.3): a 700 mg core carrying 60 mg of apalutamide in 240 mg of a 1:3 apalutamide:HPMCAS-LG spray-dried dispersion, with a 21 mg (3.0 \%) film coat. 15 625 tablets/h is the patent's PSD-3 dryer at its 75 kg/h of 5 wt \% solution (3.75 kg/h of dispersion / 240 mg). THIS STEADY FLOWSHEET IS THE MATERIAL BALANCE PER UNIT OF PRODUCT (DEV.md 4c C99): the plant runs in BATCHES (lots of final blend, released one by one) and its spray dryer in CAMPAIGNS, and the BATCH PLAN in system/postDict sizes each step for one lot from these streams' masses, counts the lots a year from the declared production, and finds the step that paces each train. The absolute rate below is a basis -- except for the spray dryer, whose hardware this flowsheet models, so its rate IS the dryer's capacity.

SOLUTION AND SPRAY DRYING (the closed nitrogen loop) acetoneFeed $\to$ [acetoneTank] $\to$ acetoneToPrep --+ apiFeed (BOUGHT) -------------------------------+$\to$ [solutionPrep] $\to$ sprayFeed --+ hpmcasFeed -------------------------------------+ | v n2Makeup --+ [SD] sprayDryer +$\to$ [n2Mix] $\to$ [blower] $\to$ [gasHeater] $\to$ dryingGas ---------$\to$ | | gasRecycle-+ \textasciicircum{} 393 K asdPowder exhaustGas | | | | [exhaustCooler] 308 K (cooling water) | | | | [solventCondenser] 263 K (NH3 refrigeration) | | | | recoveredSolvent condenserGas | | | | [recoveredTank] [purgeSplit] $\to$ purgeGas (vent) | recoveredAcetone | +--------------------------------------------------------- gasRecycle (TEAR)

DRY GRANULATION, TABLETING, COATING, PACKAGING asdPowder + MCC + croscarmellose + silica (the patent's intragranular half) $\to$ [blender] $\to$ intraBlend $\to$ [rollerCompactor] $\to$ ribbons $\to$ [mill] $\to$ milledGranules \textasciicircum{} | +---------- finesRecycle (TEAR) <--- [screen] $\to$ granules granules + silicified MCC + croscarmellose + silica + MgSt $\to$ [finalBlender] $\to$ finalBlend $\to$ [tabletPress] $\to$ tabletCores +$\to$ tabletRejects water + PVA + PEG + talc + TiO2 + iron oxide $\to$ [coatingSuspensionPrep] $\to$ coatingSuspension tabletCores + coatingSuspension $\to$ [coatingPan] $\to$ wetCoatedTablets coatingAirIn $\to$ [airHeater] 333 K $\to$ coatingAir wetCoatedTablets + coatingAir $\to$ [coatingDryer] $\to$ coatedTablets + coatingExhaust coatedTablets $\to$ [packagingLine] $\to$ packagedTablets

TWO RECYCLES, both declared in system/solverDict: the drying gas (gasRecycle) and the compactor fines (finesRecycle). The recovered acetone is NOT recycled in this base case: it goes to a tank, for QC and reuse -- the README says why, and closing that loop is a group's question.

WHAT IS A PLACEHOLDER. The blenders, the roller compactor, the mill, the screen, the tablet press and the packaging line are FLOW PRIMITIVES (a `mixer`, a `splitter`): they conserve mass and energy and model NOTHING of the machine -- no blend uniformity, no ribbon density, no particle-size reduction, no compaction force, no tablet hardness. The flowsheet has no such equipment model; none is invented here. (The SIZER can check the press's dwell time and the compaction force a target tablet strength needs WHEN the case declares the formulation's compaction curves and dwell need -- `dwell \{\}` and `compaction \{\}`, DEV.md 4c C270, witness tabletPress01\_dwell\_and\_force; this case declares neither, and its run says the press was not checked for them.) The coating pan is TWO boxes, a `mixer` (the spray lands on the cores) and a `solidDryer` (the hot air takes the water off), which is the mass and energy balance of one machine.

\end{tutnarrative}
\tutreadme{A BASE CASE for a pharmaceutical-industry project group (DEV.md 4c C97, phase A; re-based on a sourced design basis in C110). The plant makes film-coated tablets of apalutamide: the active pharmaceutical ingredient (API) is bought, everything else is made on site --- the spray-dried solid dispersion, the granules, the tablets, the coat and the packs. It runs (a few seconds from a cold start), it converges, its mass and energy balances close, and it is deliberately NOT optimised --- every number a group would change is a plain key, and every one says where it came from.}
\tutreadme{The case is sealed: everything it reads is in its own constant/.}
\tutreadme{Its design basis is a dossier, not a guess: [docs/design/apalutamide-design-basis-dossier.md](https://github.com/choupo-admin/choupo/blob/main/docs/design/apalutamide-design-basis-dossier.md) ("the dossier" below) --- the EMA and FDA assessment reports, the label, the originator's patent, the originator's published sales, generic-price evidence and a red-team review, every number with its source, page and a confidence word ($\ast$stated$\ast$, $\ast$derived$\ast$, $\ast$estimate$\ast$, $\ast$not found$\ast$). Read it before you change anything: most of what you will be asked to defend is already argued there, and what is NOT is listed in its section 6 as the questions only a supplier or the originator can answer.}
\tutreadme{Read section 3 before anything else. Half of this plant is physics Choupo computes (the spray dryer, the nitrogen loop, the solvent recovery, the coating dryer, the heat exchangers); the other half is a mass balance with no equipment model at all (blending, compaction, milling, screening, tableting, packaging). The two halves look alike in the stream table. They are not alike.}
\tutreadme{The amorphous dispersion this plant makes --- apalutamide in HPMCAS, 1:3, spray dried --- is claimed by an active patent family of the originator (US 12,303,493 B2, Aragon Pharmaceuticals / Janssen: claim 1 the dispersion, claim 4 the 1:3 ratio, claim 10 spray drying, claims 18 and 20 the 60 mg and 240 mg tablets; anticipated expiry 3 Dec 2035, an assumption by Google Patents, not a legal conclusion; the European counterparts and their SPCs were NOT verified --- dossier section 3.6). A plant that copies this formulation and sells before then infringes, unless the patents fall.}
\tutreadme{So the case declares its business case, as Vitor ruled (DEV.md 4c C99 addendum): a generic entrant to the European market at loss of exclusivity. Three consequences, each in the case:}
\begin{tutgolden}
\goldrow{stream}{acetoneFeed}{F}{0.000340765610652}
\goldrow{stream}{acetoneFeed}{T}{298.15}
\goldrow{stream}{acetoneToPrep}{F}{0.000340765610652}
\goldrow{stream}{acetoneToPrep}{T}{298.15}
\goldrow{stream}{apiFeed}{F}{5.45455180278e-07}
\goldrow{stream}{apiFeed}{T}{298.15}
\goldrow{stream}{asdPowder}{F}{1.32670518028e-06}
\goldrow{stream}{asdPowder}{T}{319.261528059}
\goldrow{stream}{blowerOut}{F}{0.00469207665386}
\goldrow{stream}{blowerOut}{T}{271.540788837}
\goldrow{stream}{ccsExtraFeed}{F}{7.59548611111e-08}
\goldrow{stream}{ccsExtraFeed}{T}{298.15}
\goldrow{stream}{ccsIntraFeed}{F}{7.59548611111e-08}
\goldrow{stream}{ccsIntraFeed}{T}{298.15}
\goldmore{1500}
\end{tutgolden}

\clearpage
\tutentry{donepezil01\_tablets}{tutorials/batch/plant/donepezil01\_tablets}
\begin{tutscheme}
\begin{tikzpicture}[>={Stealth[length=2mm]}, every node/.style={font=\scriptsize}]
  \node[draw, rounded corners=2pt, align=center, minimum width=1.9cm, minimum height=0.85cm, fill=black!4] (binderPrep) at (0.0,-0.0) {\textbf{binderPrep}\\[-1pt]{\tiny mixer}};
  \node[draw, rounded corners=2pt, align=center, minimum width=1.9cm, minimum height=0.85cm, fill=black!4] (granulator) at (2.9,-0.0) {\textbf{granulator}\\[-1pt]{\tiny mixer}};
  \node[draw, rounded corners=2pt, align=center, minimum width=1.9cm, minimum height=0.85cm, fill=black!4] (fbdAirHeater) at (0.0,-1.7) {\textbf{fbdAirHeater}\\[-1pt]{\tiny phaseChanger}};
  \node[draw, rounded corners=2pt, align=center, minimum width=1.9cm, minimum height=0.85cm, fill=black!4] (fluidBedDryer) at (5.8,-0.0) {\textbf{fluidBedDryer}\\[-1pt]{\tiny solidDryer}};
  \node[draw, rounded corners=2pt, align=center, minimum width=1.9cm, minimum height=0.85cm, fill=black!4] (mill) at (8.7,-0.0) {\textbf{mill}\\[-1pt]{\tiny splitter}};
  \node[draw, rounded corners=2pt, align=center, minimum width=1.9cm, minimum height=0.85cm, fill=black!4] (finalBlender) at (11.6,-0.0) {\textbf{finalBlender}\\[-1pt]{\tiny mixer}};
  \node[draw, rounded corners=2pt, align=center, minimum width=1.9cm, minimum height=0.85cm, fill=black!4] (tabletPress) at (14.5,-0.0) {\textbf{tabletPress}\\[-1pt]{\tiny splitter}};
  \node[draw, rounded corners=2pt, align=center, minimum width=1.9cm, minimum height=0.85cm, fill=black!4] (coatingSuspensionPrep) at (0.0,-3.4) {\textbf{coatingSuspensionPrep}\\[-1pt]{\tiny mixer}};
  \node[draw, rounded corners=2pt, align=center, minimum width=1.9cm, minimum height=0.85cm, fill=black!4] (coatingAirHeater) at (0.0,-5.1) {\textbf{coatingAirHeater}\\[-1pt]{\tiny phaseChanger}};
  \node[draw, rounded corners=2pt, align=center, minimum width=1.9cm, minimum height=0.85cm, fill=black!4] (coatingPan) at (17.4,-0.0) {\textbf{coatingPan}\\[-1pt]{\tiny mixer}};
  \node[draw, rounded corners=2pt, align=center, minimum width=1.9cm, minimum height=0.85cm, fill=black!4] (coatingDryer) at (20.3,-0.0) {\textbf{coatingDryer}\\[-1pt]{\tiny solidDryer}};
  \node[draw, rounded corners=2pt, align=center, minimum width=1.9cm, minimum height=0.85cm, fill=black!4] (packagingLine) at (23.2,-0.0) {\textbf{packagingLine}\\[-1pt]{\tiny splitter}};
  \draw[->] ($(binderPrep.west)+(-1.05,-0.17)$) -- ($(binderPrep.west)+(0,-0.17)$) node[at start, below, font=\tiny, anchor=north west] {granWater};
  \draw[->] ($(binderPrep.west)+(-1.05,0.17)$) -- ($(binderPrep.west)+(0,0.17)$) node[at start, above, font=\tiny, anchor=south west] {hpcFeed};
  \draw[->] ($(granulator.west)+(-1.05,-0.51)$) -- ($(granulator.west)+(0,-0.51)$) node[at start, below, font=\tiny, anchor=north west] {apiFeed};
  \draw[->] ($(granulator.west)+(-1.05,-0.17)$) -- ($(granulator.west)+(0,-0.17)$) node[at start, below, font=\tiny, anchor=north west] {lactoseFeed};
  \draw[->] ($(granulator.west)+(-1.05,0.17)$) -- ($(granulator.west)+(0,0.17)$) node[at start, above, font=\tiny, anchor=south west] {starchFeed};
  \draw[->] ($(granulator.west)+(-1.05,0.51)$) -- ($(granulator.west)+(0,0.51)$) node[at start, above, font=\tiny, anchor=south west] {mccFeed};
  \draw[->] (binderPrep) -- (granulator) node[midway, above, font=\tiny] {binderSolution};
  \draw[->] ($(fbdAirHeater.west)+(-1.05,0.00)$) -- ($(fbdAirHeater.west)+(0,0.00)$) node[at start, above, font=\tiny, anchor=south west] {fbdAirIn};
  \draw[->] (granulator) -- (fluidBedDryer) node[midway, above, font=\tiny] {wetMass};
  \draw[->] (fbdAirHeater) -- (fluidBedDryer) node[midway, above, font=\tiny] {fbdAir};
  \draw[->] (fluidBedDryer) -- (mill) node[midway, above, font=\tiny] {driedGranules};
  \draw[->] (mill) -- (finalBlender) node[midway, above, font=\tiny] {milledGranules};
  \draw[->] ($(finalBlender.west)+(-1.05,0.00)$) -- ($(finalBlender.west)+(0,0.00)$) node[at start, above, font=\tiny, anchor=south west] {mgstFeed};
  \draw[->] (finalBlender) -- (tabletPress) node[midway, above, font=\tiny] {finalBlend};
  \draw[->] ($(coatingSuspensionPrep.west)+(-1.05,-0.85)$) -- ($(coatingSuspensionPrep.west)+(0,-0.85)$) node[at start, below, font=\tiny, anchor=north west] {coatWater};
  \draw[->] ($(coatingSuspensionPrep.west)+(-1.05,-0.51)$) -- ($(coatingSuspensionPrep.west)+(0,-0.51)$) node[at start, below, font=\tiny, anchor=north west] {hpmcFeed};
  \draw[->] ($(coatingSuspensionPrep.west)+(-1.05,-0.17)$) -- ($(coatingSuspensionPrep.west)+(0,-0.17)$) node[at start, below, font=\tiny, anchor=north west] {pegFeed};
  \draw[->] ($(coatingSuspensionPrep.west)+(-1.05,0.17)$) -- ($(coatingSuspensionPrep.west)+(0,0.17)$) node[at start, above, font=\tiny, anchor=south west] {talcFeed};
  \draw[->] ($(coatingSuspensionPrep.west)+(-1.05,0.51)$) -- ($(coatingSuspensionPrep.west)+(0,0.51)$) node[at start, above, font=\tiny, anchor=south west] {tio2Feed};
  \draw[->] ($(coatingSuspensionPrep.west)+(-1.05,0.85)$) -- ($(coatingSuspensionPrep.west)+(0,0.85)$) node[at start, above, font=\tiny, anchor=south west] {feOxFeed};
  \draw[->] ($(coatingAirHeater.west)+(-1.05,0.00)$) -- ($(coatingAirHeater.west)+(0,0.00)$) node[at start, above, font=\tiny, anchor=south west] {coatingAirIn};
  \draw[->] (tabletPress) -- (coatingPan) node[midway, above, font=\tiny] {tabletCores};
  \draw[->] (coatingSuspensionPrep) -- (coatingPan) node[midway, above, font=\tiny] {coatingSuspension};
  \draw[->] (coatingPan) -- (coatingDryer) node[midway, above, font=\tiny] {wetCoatedTablets};
  \draw[->] (coatingAirHeater) -- (coatingDryer) node[midway, above, font=\tiny] {coatingAir};
  \draw[->] (coatingDryer) -- (packagingLine) node[midway, above, font=\tiny] {coatedTablets};
  \draw[->] ($(fluidBedDryer.east)+(0,0.00)$) -- ($(fluidBedDryer.east)+(1.05,0.00)$) node[at end, above, font=\tiny, anchor=south east] {fbdExhaust};
  \draw[->] ($(tabletPress.east)+(0,0.00)$) -- ($(tabletPress.east)+(1.05,0.00)$) node[at end, above, font=\tiny, anchor=south east] {tabletRejects};
  \draw[->] ($(coatingDryer.east)+(0,0.00)$) -- ($(coatingDryer.east)+(1.05,0.00)$) node[at end, above, font=\tiny, anchor=south east] {coatingExhaust};
  \draw[->] ($(packagingLine.east)+(0,0.00)$) -- ($(packagingLine.east)+(1.05,0.00)$) node[at end, above, font=\tiny, anchor=south east] {packagedTablets};
\end{tikzpicture}
\end{tutscheme}
\tutwhat{DONEPEZIL HYDROCHLORIDE 10 mg FILM-COATED TABLETS, a drug-product plant (the API is BOUGHT): WET GRANULATION -- the API, lactose monohydrate, maize starch and microcrystalline cellulose wetted in a high-shear granulator with an aqueous hydroxypropyl-cellulose binder solution and dried in a fluid bed (`solidDryer`, hot air, the starch's sorption isotherm) -- milling, a final blend with magnesium stearate, tableting with rejects, an aqueous hypromellose film coat dried off by hot air, and packaging. The excipients are the reference product's (ARICEPT SmPC section 6.1) and the 174.33 mg of lactose per 10 mg tablet is its section 2; every other quantity is an author-set assumption. The steady flowsheet is the material balance per unit of product (720 000 tablets/h pressed, the fluid bed's own drying rate, is a basis); a BATCH PLAN (system/postDict) sizes each step for a 300 kg lot, counts the lots a year for 150 million tablets, and finds the pace of the granulation and tableting trains. Blending, granulation, milling, tableting and packaging are MASS-BALANCE PLACEHOLDERS (no equipment physics). No organic solvent, no spray drying, no hazardous area. One strength (10 mg); the two-strength year on the same equipment, campaigns and changeovers, is donepezil02\_two\_strengths.}
\begin{tutnarrative}
a drug-product plant: donepezil hydrochloride 10 mg film-coated tablets. The API is BOUGHT; everything else is made here (DEV.md 4c C104, built on the stated defaults of C110). A BASE CASE: it runs, it is NOT optimised, and every number a group would change is a plain key, said to be an assumption where it is declared. The README is the narrative; this file is the topology.

THE BASIS (README section 2): 720 000 tablets per hour PRESSED, the flowsheet basis, of a 280 mg core carrying 10 mg of donepezil hydrochloride and 174.33 mg of lactose monohydrate (the reference product's own figures), with a 3 \% film coat. THIS STEADY FLOWSHEET IS THE MATERIAL BALANCE PER UNIT OF PRODUCT (DEV.md 4c C99): the plant runs in BATCHES (lots released one by one), and the BATCH PLAN in system/postDict sizes each step for one lot from these streams' masses, counts the lots a year from the declared production, and finds the step that paces each train. The absolute rate is chosen as the fluid bed's own DRYING rate (one 300 kg lot dried in 1.5 h), so the fluid bed's air heater is sized for the duty it really has; the coater's heater is then sized at that same rate, which is about twice the coater's own -- README section 6.

WET GRANULATION granWater + hpcFeed $\to$ [binderPrep] $\to$ binderSolution --+ apiFeed (BOUGHT) + lactose + starch + MCC --------------+$\to$ [granulator] $\to$ wetMass fbdAirIn $\to$ [fbdAirHeater] 333 K $\to$ fbdAir | wetMass + fbdAir $\to$ [fluidBedDryer] $\to$ driedGranules + fbdExhaust driedGranules $\to$ [mill] $\to$ milledGranules

FINAL BLEND, TABLETING, COATING, PACKAGING milledGranules + mgstFeed $\to$ [finalBlender] $\to$ finalBlend $\to$ [tabletPress] $\to$ tabletCores +$\to$ tabletRejects coatWater + HPMC + macrogol + talc + TiO2 + yellow iron oxide $\to$ [coatingSuspensionPrep] tabletCores + coatingSuspension $\to$ [coatingPan] $\to$ wetCoatedTablets coatingAirIn $\to$ [coatingAirHeater] 333 K $\to$ coatingAir wetCoatedTablets + coatingAir $\to$ [coatingDryer] $\to$ coatedTablets + coatingExhaust coatedTablets $\to$ [packagingLine] $\to$ packagedTablets

NO RECYCLE, NO SOLVENT. Unlike apalutamide01\_tablets there is no spray dryer, no organic solvent, no nitrogen loop, no solvent recovery and no hazardous area: an aqueous wet granulation and an aqueous coat.

WHAT IS A PLACEHOLDER. The high-shear granulator, the mill, the final blender, the tablet press and the packaging line are FLOW PRIMITIVES (a `mixer`, a `splitter`): they conserve mass and energy and model NOTHING of the machine -- no granule growth, no granule size, no blend uniformity (THE critical quality question of a low-dose tablet, README section 8), no compaction force, no tablet hardness. The flowsheet has no such equipment model; none is invented here. (The SIZER can check the press's dwell time and the compaction force a target tablet strength needs WHEN the case declares the formulation's compaction curves and dwell need -- `dwell \{\}` and `compaction \{\}`, DEV.md 4c C270, witness tabletPress01\_dwell\_and\_force; this case declares neither, and its run says the press was not checked for them.) The two DRYERS are physics: `solidDryer` dries to the starch's sorption isotherm at the INLET air's humidity (a batch dried to its end point ends in equilibrium with the fresh air), with the energy balance on the package's enthalpy. The coating pan is TWO boxes, a `mixer` (the spray lands on the cores) and a `solidDryer` (the hot air takes the water off), which is the mass and energy balance of one machine.

\end{tutnarrative}
\tutreadme{A BASE CASE for a pharmaceutical-industry project group (DEV.md 4c C104, built on the stated defaults of C110). The plant makes film-coated tablets of donepezil hydrochloride by wet granulation: the active pharmaceutical ingredient (API) is bought, everything else is made on site --- the granules, the tablets, the coat and the packs. It runs (in seconds), its mass and energy balances close, and it is deliberately NOT optimised --- every number a group would change is a plain key, said to be an assumption where it is declared.}
\tutreadme{The case is sealed: everything it reads is in its own constant/.}
\tutreadme{It is the simpler sibling of apalutamide01\_tablets (the same folder, the same grammar, read that case's README for anything not repeated here). There is no spray drying, no organic solvent, no nitrogen loop, no solvent recovery and no hazardous area (ATEX): an aqueous binder, an aqueous coat, two hot-air dryers. What is harder here is not the chemistry but the dose: 10 mg of API in a 280 mg core (section 8).}
\tutreadme{Read section 3 before anything else. Two units of this plant are physics Choupo computes (the fluid-bed dryer and the coating dryer, with their air heaters); the rest is a mass balance with no equipment model at all (granulation, milling, blending, tableting, packaging). The two look alike in the stream table. They are not alike.}
\tutreadme{Two inferences, labelled as such: the lactose of the two strengths is in the ratio 2:1, which is what ONE common blend pressed at two tablet weights would give (a dose-proportional pair); and the yellow iron oxide is the only ingredient the two strengths do not share --- so the two products can share the whole granulation and differ only from the press on.}
\tutreadme{Where the group must look for what this case could not verify: ARICEPT was not authorised through the EMA's centralised procedure as far as this case could find, so there is no EPAR to read; the national SmPCs and public assessment reports (MHRA, HPRA, the other national agencies) and the public assessment reports of the generic donepezil products (written for each generic's authorisation, and which may describe its manufacturing route) are where the formulation and the process are described; the Ph. Eur. monographs of the excipients for their grades. Replace each assumption with what you find, and cite it.}
\begin{tutgolden}
\goldrow{stream}{apiFeed}{F}{4.80817775e-06}
\goldrow{stream}{apiFeed}{T}{298.15}
\goldrow{stream}{binderSolution}{F}{0.000740863846794}
\goldrow{stream}{binderSolution}{T}{298.15}
\goldrow{stream}{coatWater}{F}{0.000677035803497}
\goldrow{stream}{coatWater}{T}{298.15}
\goldrow{stream}{coatedTablets}{F}{0.000238322639252}
\goldrow{stream}{coatedTablets}{T}{317.588908471}
\goldrow{stream}{coatingAir}{F}{0.0702777777778}
\goldrow{stream}{coatingAir}{T}{333.15}
\goldrow{stream}{coatingAirIn}{F}{0.0702777777778}
\goldrow{stream}{coatingAirIn}{T}{298.15}
\goldrow{stream}{coatingExhaust}{F}{0.0709548135813}
\goldrow{stream}{coatingExhaust}{T}{317.588908471}
\goldmore{1103}
\end{tutgolden}

\clearpage
\tutentry{donepezil02\_two\_strengths}{tutorials/batch/plant/donepezil02\_two\_strengths}
\begin{tutscheme}
\begin{tikzpicture}[>={Stealth[length=2mm]}, every node/.style={font=\scriptsize}]
  \node[draw, rounded corners=2pt, align=center, minimum width=1.9cm, minimum height=0.85cm, fill=black!4] (binderPrep) at (0.0,-0.0) {\textbf{binderPrep}\\[-1pt]{\tiny mixer}};
  \node[draw, rounded corners=2pt, align=center, minimum width=1.9cm, minimum height=0.85cm, fill=black!4] (granulator) at (2.9,-0.0) {\textbf{granulator}\\[-1pt]{\tiny mixer}};
  \node[draw, rounded corners=2pt, align=center, minimum width=1.9cm, minimum height=0.85cm, fill=black!4] (fbdAirHeater) at (0.0,-1.7) {\textbf{fbdAirHeater}\\[-1pt]{\tiny phaseChanger}};
  \node[draw, rounded corners=2pt, align=center, minimum width=1.9cm, minimum height=0.85cm, fill=black!4] (fluidBedDryer) at (5.8,-0.0) {\textbf{fluidBedDryer}\\[-1pt]{\tiny solidDryer}};
  \node[draw, rounded corners=2pt, align=center, minimum width=1.9cm, minimum height=0.85cm, fill=black!4] (mill) at (8.7,-0.0) {\textbf{mill}\\[-1pt]{\tiny splitter}};
  \node[draw, rounded corners=2pt, align=center, minimum width=1.9cm, minimum height=0.85cm, fill=black!4] (finalBlender) at (11.6,-0.0) {\textbf{finalBlender}\\[-1pt]{\tiny mixer}};
  \node[draw, rounded corners=2pt, align=center, minimum width=1.9cm, minimum height=0.85cm, fill=black!4] (tabletPress) at (14.5,-0.0) {\textbf{tabletPress}\\[-1pt]{\tiny splitter}};
  \node[draw, rounded corners=2pt, align=center, minimum width=1.9cm, minimum height=0.85cm, fill=black!4] (coatingSuspensionPrep) at (0.0,-3.4) {\textbf{coatingSuspensionPrep}\\[-1pt]{\tiny mixer}};
  \node[draw, rounded corners=2pt, align=center, minimum width=1.9cm, minimum height=0.85cm, fill=black!4] (coatingAirHeater) at (0.0,-5.1) {\textbf{coatingAirHeater}\\[-1pt]{\tiny phaseChanger}};
  \node[draw, rounded corners=2pt, align=center, minimum width=1.9cm, minimum height=0.85cm, fill=black!4] (coatingPan) at (17.4,-0.0) {\textbf{coatingPan}\\[-1pt]{\tiny mixer}};
  \node[draw, rounded corners=2pt, align=center, minimum width=1.9cm, minimum height=0.85cm, fill=black!4] (coatingDryer) at (20.3,-0.0) {\textbf{coatingDryer}\\[-1pt]{\tiny solidDryer}};
  \node[draw, rounded corners=2pt, align=center, minimum width=1.9cm, minimum height=0.85cm, fill=black!4] (packagingLine) at (23.2,-0.0) {\textbf{packagingLine}\\[-1pt]{\tiny splitter}};
  \draw[->] ($(binderPrep.west)+(-1.05,-0.17)$) -- ($(binderPrep.west)+(0,-0.17)$) node[at start, below, font=\tiny, anchor=north west] {granWater};
  \draw[->] ($(binderPrep.west)+(-1.05,0.17)$) -- ($(binderPrep.west)+(0,0.17)$) node[at start, above, font=\tiny, anchor=south west] {hpcFeed};
  \draw[->] ($(granulator.west)+(-1.05,-0.51)$) -- ($(granulator.west)+(0,-0.51)$) node[at start, below, font=\tiny, anchor=north west] {apiFeed};
  \draw[->] ($(granulator.west)+(-1.05,-0.17)$) -- ($(granulator.west)+(0,-0.17)$) node[at start, below, font=\tiny, anchor=north west] {lactoseFeed};
  \draw[->] ($(granulator.west)+(-1.05,0.17)$) -- ($(granulator.west)+(0,0.17)$) node[at start, above, font=\tiny, anchor=south west] {starchFeed};
  \draw[->] ($(granulator.west)+(-1.05,0.51)$) -- ($(granulator.west)+(0,0.51)$) node[at start, above, font=\tiny, anchor=south west] {mccFeed};
  \draw[->] (binderPrep) -- (granulator) node[midway, above, font=\tiny] {binderSolution};
  \draw[->] ($(fbdAirHeater.west)+(-1.05,0.00)$) -- ($(fbdAirHeater.west)+(0,0.00)$) node[at start, above, font=\tiny, anchor=south west] {fbdAirIn};
  \draw[->] (granulator) -- (fluidBedDryer) node[midway, above, font=\tiny] {wetMass};
  \draw[->] (fbdAirHeater) -- (fluidBedDryer) node[midway, above, font=\tiny] {fbdAir};
  \draw[->] (fluidBedDryer) -- (mill) node[midway, above, font=\tiny] {driedGranules};
  \draw[->] (mill) -- (finalBlender) node[midway, above, font=\tiny] {milledGranules};
  \draw[->] ($(finalBlender.west)+(-1.05,0.00)$) -- ($(finalBlender.west)+(0,0.00)$) node[at start, above, font=\tiny, anchor=south west] {mgstFeed};
  \draw[->] (finalBlender) -- (tabletPress) node[midway, above, font=\tiny] {finalBlend};
  \draw[->] ($(coatingSuspensionPrep.west)+(-1.05,-0.85)$) -- ($(coatingSuspensionPrep.west)+(0,-0.85)$) node[at start, below, font=\tiny, anchor=north west] {coatWater};
  \draw[->] ($(coatingSuspensionPrep.west)+(-1.05,-0.51)$) -- ($(coatingSuspensionPrep.west)+(0,-0.51)$) node[at start, below, font=\tiny, anchor=north west] {hpmcFeed};
  \draw[->] ($(coatingSuspensionPrep.west)+(-1.05,-0.17)$) -- ($(coatingSuspensionPrep.west)+(0,-0.17)$) node[at start, below, font=\tiny, anchor=north west] {pegFeed};
  \draw[->] ($(coatingSuspensionPrep.west)+(-1.05,0.17)$) -- ($(coatingSuspensionPrep.west)+(0,0.17)$) node[at start, above, font=\tiny, anchor=south west] {talcFeed};
  \draw[->] ($(coatingSuspensionPrep.west)+(-1.05,0.51)$) -- ($(coatingSuspensionPrep.west)+(0,0.51)$) node[at start, above, font=\tiny, anchor=south west] {tio2Feed};
  \draw[->] ($(coatingSuspensionPrep.west)+(-1.05,0.85)$) -- ($(coatingSuspensionPrep.west)+(0,0.85)$) node[at start, above, font=\tiny, anchor=south west] {feOxFeed};
  \draw[->] ($(coatingAirHeater.west)+(-1.05,0.00)$) -- ($(coatingAirHeater.west)+(0,0.00)$) node[at start, above, font=\tiny, anchor=south west] {coatingAirIn};
  \draw[->] (tabletPress) -- (coatingPan) node[midway, above, font=\tiny] {tabletCores};
  \draw[->] (coatingSuspensionPrep) -- (coatingPan) node[midway, above, font=\tiny] {coatingSuspension};
  \draw[->] (coatingPan) -- (coatingDryer) node[midway, above, font=\tiny] {wetCoatedTablets};
  \draw[->] (coatingAirHeater) -- (coatingDryer) node[midway, above, font=\tiny] {coatingAir};
  \draw[->] (coatingDryer) -- (packagingLine) node[midway, above, font=\tiny] {coatedTablets};
  \draw[->] ($(fluidBedDryer.east)+(0,0.00)$) -- ($(fluidBedDryer.east)+(1.05,0.00)$) node[at end, above, font=\tiny, anchor=south east] {fbdExhaust};
  \draw[->] ($(tabletPress.east)+(0,0.00)$) -- ($(tabletPress.east)+(1.05,0.00)$) node[at end, above, font=\tiny, anchor=south east] {tabletRejects};
  \draw[->] ($(coatingDryer.east)+(0,0.00)$) -- ($(coatingDryer.east)+(1.05,0.00)$) node[at end, above, font=\tiny, anchor=south east] {coatingExhaust};
  \draw[->] ($(packagingLine.east)+(0,0.00)$) -- ($(packagingLine.east)+(1.05,0.00)$) node[at end, above, font=\tiny, anchor=south east] {packagedTablets};
\end{tikzpicture}
\end{tutscheme}
\tutwhat{DONEPEZIL HYDROCHLORIDE 10 mg AND 5 mg FILM-COATED TABLETS, TWO PRODUCTS in campaigns on one drug-product plant (the API is BOUGHT): WET GRANULATION -- the API, lactose monohydrate, maize starch and microcrystalline cellulose wetted in a high-shear granulator with an aqueous hydroxypropyl-cellulose binder solution and dried in a fluid bed (`solidDryer`, hot air, the starch's sorption isotherm) -- milling, a final blend with magnesium stearate, tableting with rejects, an aqueous hypromellose film coat dried off by hot air, and packaging. The excipients are the reference product's (ARICEPT SmPC section 6.1) and the 174.33 mg of lactose per 10 mg tablet is its section 2; every other quantity is an author-set assumption. The steady flowsheet is the material balance per unit of product (720 000 tablets/h pressed, the fluid bed's own drying rate, is a basis); a MULTIPRODUCT BATCH PLAN (system/postDict) sizes each step for a 300 kg lot of either strength, counts the lots a year for 90 million 10 mg and 60 million 5 mg tablets, runs them in single-product campaigns in a declared order with a direction-dependent clean between them (colour carry-over: white after yellow is the long one), and finds the pace of the granulation and tableting trains. Blending, granulation, milling, tableting and packaging are MASS-BALANCE PLACEHOLDERS (no equipment physics). No organic solvent, no spray drying, no hazardous area. The 5 mg tablet is the 10 mg tablet at half the weight (the same blend), so the one steady balance serves both.}
\begin{tutnarrative}
a drug-product plant: donepezil hydrochloride 10 mg AND 5 mg film-coated tablets, two products in campaigns on one plant (the steady balance below is both strengths': the 5 mg tablet is the 10 mg tablet at half the weight, README section 2). The API is BOUGHT; everything else is made here (DEV.md 4c C104, built on the stated defaults of C110). A BASE CASE: it runs, it is NOT optimised, and every number a group would change is a plain key, said to be an assumption where it is declared. The README is the narrative; this file is the topology.

THE BASIS (README section 2): 720 000 tablets per hour PRESSED, the flowsheet basis, of a 280 mg core carrying 10 mg of donepezil hydrochloride and 174.33 mg of lactose monohydrate (the reference product's own figures), with a 3 \% film coat. THIS STEADY FLOWSHEET IS THE MATERIAL BALANCE PER UNIT OF PRODUCT (DEV.md 4c C99): the plant runs in BATCHES (lots released one by one), and the BATCH PLAN in system/postDict sizes each step for one lot from these streams' masses, counts the lots a year from the declared production, and finds the step that paces each train. The absolute rate is chosen as the fluid bed's own DRYING rate (one 300 kg lot dried in 1.5 h), so the fluid bed's air heater is sized for the duty it really has; the coater's heater is then sized at that same rate, which is about twice the coater's own -- README section 6.

WET GRANULATION granWater + hpcFeed $\to$ [binderPrep] $\to$ binderSolution --+ apiFeed (BOUGHT) + lactose + starch + MCC --------------+$\to$ [granulator] $\to$ wetMass fbdAirIn $\to$ [fbdAirHeater] 333 K $\to$ fbdAir | wetMass + fbdAir $\to$ [fluidBedDryer] $\to$ driedGranules + fbdExhaust driedGranules $\to$ [mill] $\to$ milledGranules

FINAL BLEND, TABLETING, COATING, PACKAGING milledGranules + mgstFeed $\to$ [finalBlender] $\to$ finalBlend $\to$ [tabletPress] $\to$ tabletCores +$\to$ tabletRejects coatWater + HPMC + macrogol + talc + TiO2 + yellow iron oxide $\to$ [coatingSuspensionPrep] tabletCores + coatingSuspension $\to$ [coatingPan] $\to$ wetCoatedTablets coatingAirIn $\to$ [coatingAirHeater] 333 K $\to$ coatingAir wetCoatedTablets + coatingAir $\to$ [coatingDryer] $\to$ coatedTablets + coatingExhaust coatedTablets $\to$ [packagingLine] $\to$ packagedTablets

NO RECYCLE, NO SOLVENT. Unlike apalutamide01\_tablets there is no spray dryer, no organic solvent, no nitrogen loop, no solvent recovery and no hazardous area: an aqueous wet granulation and an aqueous coat.

WHAT IS A PLACEHOLDER. The high-shear granulator, the mill, the final blender, the tablet press and the packaging line are FLOW PRIMITIVES (a `mixer`, a `splitter`): they conserve mass and energy and model NOTHING of the machine -- no granule growth, no granule size, no blend uniformity (THE critical quality question of a low-dose tablet, README section 8), no compaction force, no tablet hardness. The flowsheet has no such equipment model; none is invented here. (The SIZER can check the press's dwell time and the compaction force a target tablet strength needs WHEN the case declares the formulation's compaction curves and dwell need -- `dwell \{\}` and `compaction \{\}`, DEV.md 4c C270, witness tabletPress01\_dwell\_and\_force; this case declares neither, and its run says the press was not checked for them.) The two DRYERS are physics: `solidDryer` dries to the starch's sorption isotherm at the INLET air's humidity (a batch dried to its end point ends in equilibrium with the fresh air), with the energy balance on the package's enthalpy. The coating pan is TWO boxes, a `mixer` (the spray lands on the cores) and a `solidDryer` (the hot air takes the water off), which is the mass and energy balance of one machine.

\end{tutnarrative}
\tutreadme{1. Explain the base result. The plant makes 150 million tablets (34 828 kg) in 115 lots and loses money: NPV -29.88 M EUR. Find the term of COM\_d that makes the loss (compare 2.73 C\_OL with 1.23 (C\_RM + C\_UT) in the Economic Appraisal), and use the Bases table of the batch plan so that 300 kg, 305.9 kg, 168 h, 928 h, 2127 h and 3840 h each mean one thing. 2. Find the constraint that matters. Two different words: the limiting stage (sets T\_L) of each train --- drying (4.25 h a lot) on granulation, packaging (11 h a 10 mg lot, 20 h a 5 mg lot) on tableting --- and the machine with the least capacity margin, the granulator at 92.8 \% of its bowl. Say which one limits the lot and which one the year. 3. Propose one improvement and justify it with a new run. Change ONE declaration (the crew, the lot, outOfPhase 2 on packaging, the campaign length), run again, and show the number that moved and the one that did not.}
\tutreadme{A BASE CASE for a pharmaceutical-industry project group (DEV.md 4c C249). It takes the plant of donepezil01\_tablets (wet granulation, an aqueous coat, the API bought) and asks the question that case could not answer: how is a year divided between two strengths made on the same equipment, and what does it cost to change from one to the other? The answer is a multiproduct batch plan --- single-product campaigns in a declared order, a clean between products that depends on the direction of the change, and a Gantt chart of the year. It runs (in seconds), it is deliberately NOT optimised, and every number a group would change is a plain key, said to be an assumption where it is declared.}
\tutreadme{The case is sealed: everything it reads is in its own constant/. The flowsheet, the components and the building are donepezil01\_tablets' own, unchanged (read that README for the chemistry, the dryers and the questions about content uniformity); this one is about the plan.}
\tutreadme{One steady flowsheet serves both, because the 5 mg tablet is the same blend at half the weight: a kilogram of either carries the same kilograms of every ingredient, and the batch plan reads its balance per kg of product (BatchPlanPass.H:189-193). The one thing that would differ --- the 10 mg coat's 1 \% yellow iron oxide, 0.03 \% of the tablet --- is ignored here and said so; a group that wants it declares a massPerProduct for the coating suspension and the plan announces ownBalance.}
\tutreadme{Two trains decoupled by the final-blend store, as in donepezil01\_tablets; lotsPerCampaign 8 and a 24 h full clean between campaigns of ONE product. Between the two products every stage has its own clean (productChangeover): the train's 24 h campaign clean PLUS the product's own part (a rinse check, the press tooling, the printed packaging), never less, so changing product never costs less time than continuing one; a product clean declared below the campaign changeover is announced. The train waits for the slowest. The Gantt of the year (campaignOrder ( strength10 strength5 )):}
\tutreadme{The tableting train is paced by the packaging line for both products, but at 11 h a lot for one and 20 h for the other: the limiting stage is a property of the product, not only of the plant. Equipment is sized for the product that needs the most (the plan names it, sizing by): here the 10 mg for every stage, because the 5 mg lot is the same mass and only takes longer.}
\begin{tutgolden}
\goldrow{stream}{apiFeed}{F}{4.80817775e-06}
\goldrow{stream}{apiFeed}{T}{298.15}
\goldrow{stream}{binderSolution}{F}{0.000740863846794}
\goldrow{stream}{binderSolution}{T}{298.15}
\goldrow{stream}{coatWater}{F}{0.000677035803497}
\goldrow{stream}{coatWater}{T}{298.15}
\goldrow{stream}{coatedTablets}{F}{0.000238322639252}
\goldrow{stream}{coatedTablets}{T}{317.588908471}
\goldrow{stream}{coatingAir}{F}{0.0702777777778}
\goldrow{stream}{coatingAir}{T}{333.15}
\goldrow{stream}{coatingAirIn}{F}{0.0702777777778}
\goldrow{stream}{coatingAirIn}{T}{298.15}
\goldrow{stream}{coatingExhaust}{F}{0.0709548135813}
\goldrow{stream}{coatingExhaust}{T}{317.588908471}
\goldmore{1477}
\end{tutgolden}

\clearpage
\tutentry{tabletPress01\_dwell\_and\_force}{tutorials/batch/plant/tabletPress01\_dwell\_and\_force}
\begin{tutscheme}
\begin{tikzpicture}[>={Stealth[length=2mm]}, every node/.style={font=\scriptsize}]
  \node[draw, rounded corners=2pt, align=center, minimum width=1.9cm, minimum height=0.85cm, fill=black!4] (tabletPress) at (0.0,-0.0) {\textbf{tabletPress}\\[-1pt]{\tiny splitter}};
  \draw[->] ($(tabletPress.west)+(-1.05,0.00)$) -- ($(tabletPress.west)+(0,0.00)$) node[at start, above, font=\tiny, anchor=south west] {finalBlend};
  \draw[->] ($(tabletPress.east)+(0,-0.17)$) -- ($(tabletPress.east)+(1.05,-0.17)$) node[at end, below, font=\tiny, anchor=north east] {tabletCores};
  \draw[->] ($(tabletPress.east)+(0,0.17)$) -- ($(tabletPress.east)+(1.05,0.17)$) node[at end, above, font=\tiny, anchor=south east] {tabletRejects};
\end{tikzpicture}
\end{tutscheme}
\tutwhat{THE TABLET PRESS, RATED ON DWELL TIME AND ON COMPACTION FORCE (DEV.md 4c C270). One rotary press makes placebo tablets of microcrystalline cellulose for a plant that runs in lots: the flowsheet is a mass split (99 \% cores, 1 \% rejects), a BATCH PLAN gives the lot and the press's operating time, and the `ratedMachine` sizer checks the press three ways -- its catalogue speed against the formulation's minimum dwell time (the dwell limit binds: 51 rpm, not the catalogue's 75), a declared overall equipment effectiveness, and the compaction pressure and punch force a 2 MPa tablet needs, found by inverting the Heckel and Ryshkewitch-Duckworth curves that tutorials/props/powder/compaction01\_mcc FITTED to open measurements of MCC. The press geometry, speed, maximum force and the plant's year are author-set assumptions.}
\begin{tutnarrative}
ONE rotary tablet press. The README is the narrative; this file is the topology.

finalBlend $\to$ [tabletPress] $\to$ tabletCores +$\to$ tabletRejects

THE PRESS IS A PLACEHOLDER IN THE FLOWSHEET: a `splitter` conserves mass and models nothing of the machine. The machine is CHECKED by the sizing pass (system/postDict, `type ratedMachine`, kind tabletPress), which reads the lot from the batch plan and the formulation's needs from the case.

THE BASIS: 100 kg/h of final blend (pure microcrystalline cellulose, a placebo), the steady material balance per hour of pressing; the batch plan scales it to the lot. An author-set round number.

\end{tutnarrative}
\tutreadme{One rotary tablet press makes placebo tablets of microcrystalline cellulose (MCC) for a plant that runs in lots (DEV.md 4c C270). The case asks the two questions a press catalogue does not answer:}
\tutreadme{1. How fast can this press really turn for this powder? The catalogue gives a turret speed. The powder needs a minimum time under the compression roll, the dwell time, and that time sets the speed. 2. How hard must the press push to make a tablet that is strong enough? The answer comes from the two compaction curves the compaction bench fitted to open measurements of MCC, run backwards.}
\tutreadme{The case is sealed: everything it reads is in its own constant/.}
\tutreadme{- The flowsheet (system/flowsheetDict) is one splitter: 100 kg/h of final blend in, 99 \% tablet cores and 1 \% rejects out. It is a mass balance and nothing else: the splitter knows nothing about punches, speeds or forces. - The batch plan (system/postDict, batchPlan) gives the press a lot of 200 kg of blend (500 000 tablets of 0.40 g) and 7 hours to press it. - The sizing pass (system/postDict, sizing, type ratedMachine, kind tabletPress) checks the press. This is where all the arithmetic below happens, and the run prints every step of it.}
\tutreadme{Every number that no source gave is an author-set assumption, and the dictionary says so beside it: the press (36 stations, 75 rpm, an 8 mm punch flat, a 0.30 m pitch circle, an 80 kN maximum force), the year, the lot and the effectiveness factors. Replace them with your press's data sheet and your own plan.}
\tutreadme{The punch head has a flat. While the flat passes under the roll, the punch stays at its lowest point and the powder is held at full compression. At a turret speed of n revolutions per second, a point on the pitch circle of diameter D moves at pi D n, so the flat of length L stays under the roll for}
\begin{tutgolden}
\goldrow{stream}{finalBlend}{F}{0.000171318824722}
\goldrow{stream}{finalBlend}{T}{298.15}
\goldrow{stream}{tabletCores}{F}{0.000169605636475}
\goldrow{stream}{tabletCores}{T}{298.15}
\goldrow{stream}{tabletRejects}{F}{1.71318824722e-06}
\goldrow{stream}{tabletRejects}{T}{298.15}
\goldrow{equipment}{tabletPress}{basis}{load\_=\_tablets\_per\_hour\_=\_the\_mass\_the\_press's\_inlet\_carries\_per\_lot\_/\_the\_pressed\_tablet's\_mass\_/\_the\_stage's\_operating\_time\_(read\_from\_the\_batch\_plan);\_capacity\_=\_machines\_x\_stations\_x\_turret\_rpm\_x\_60\_x\_compressions\_per\_revolution\_(nameplate;\_the\_DWELL\_LIMIT\_BINDS:\_rated\_at\_n\_=\_L\_flat\_/\_(pi\_D\_pitch\_t\_min),\_below\_the\_catalogue\_speed;\_x\_the\_declared\_overall\_equipment\_effectiveness\_(availability\_x\_performance\_x\_quality));\_COMPACTION\_(one\_fitted\_material,\_MCC\_Avicel\_PH200\_(the\_compaction\_bench,\_op\_mcc\_inDie),\_no\_mixture\_rule):\_porosity\_from\_Ryshkewitch-Duckworth\_inverted\_at\_the\_target\_strength,\_pressure\_from\_Heckel\_inverted\_at\_its\_solid\_fraction,\_force\_=\_pressure\_x\_punch\_tip\_area.\_The\_rating\_is\_the\_CASE's\_(tabletPress,\_declared\_with\_its\_source);\_the\_flowsheet's\_press\_is\_a\_mass\_split\_and\_models\_no\_compaction:\_the\_sizer\_checks\_the\_dwell\_time\_and\_the\_compaction\_force\_ONLY\_where\_the\_case\_declares\_the\_formulation's\_needs\_(`dwell\_\{\_\}`,\_`compaction\_\{\_\}`),\_and\_an\_overall\_equipment\_effectiveness\_only\_where\_it\_declares\_one\_(`effectiveness\_\{\_\}`)}
\goldrow{equipment}{tabletPress}{tagNo}{TP-101}
\goldrow{equipment}{tabletPress}{values.availability}{0.9}
\goldrow{equipment}{tabletPress}{values.capacity\_dwellLimited}{110007.896665}
\goldrow{equipment}{tabletPress}{values.capacity\_effective}{94056.7516487}
\goldrow{equipment}{tabletPress}{values.capacity\_machine}{162000}
\goldrow{equipment}{tabletPress}{values.compactionPressure}{29.7460137671}
\goldrow{equipment}{tabletPress}{values.compressionsPerRevolution}{1}
\goldmore{72}
\end{tutgolden}

\subsection{reactor}
\grouptheory{reactor}
\clearpage
\tutentry{batch01\_first\_order}{tutorials/batch/reactor/batch01\_first\_order}
\begin{tutscheme}
\begin{tikzpicture}[>={Stealth[length=2mm]}, every node/.style={font=\scriptsize}]
  \node[draw, rounded corners=2pt, align=center, minimum width=1.9cm, minimum height=0.85cm, fill=black!4] (reactor) at (0.0,-0.0) {\textbf{reactor}\\[-1pt]{\tiny batchReactor}};
\end{tikzpicture}
\end{tutscheme}
\tutwhat{Isothermal batch Fischer esterification (AcOH + EtOH <$\to$ EtOAc + H2O) at 350 K, pseudo-first-order in ethanol; analytical exp-decay comparison}
\begin{tutnarrative}
isothermal pseudo-first-order batch reactor.

Single vessel, pseudo-first-order esterification (AcOH + EtOH $\to$ EtOAc + H2O, rate = k.[ethanol]). Initial charge is an equimolar ethanol + acetic acid mix; kinetics from the reactions library (Arrhenius A=1e8 1/s, Ea=70 kJ/mol). At T = 350 K:

k(T) = 1e8 $\cdot$ exp(-70000 / (8.314 $\cdot$ 350)) $\approx$ 3.57e-3 s\textasciicircum{}-1 $\tau$ = 1/k $\approx$ 280 s

Analytical solution for first-order in a constant-volume batch:

n\_A(t) = n\_A(0) $\cdot$ exp(-k t)

The simulator integrates dn\_A/dt = -k $\cdot$ (n\_A / V) $\cdot$ V = -k n\_A with RK4 at dt = 1 s; final t = 600 s gives X\_A $\approx$ 1 - exp(-k$\cdot$600) $\approx$ 1 - exp(-2.14) $\approx$ 0.882

\end{tutnarrative}
\tutreadme{18.5 mol of an equimolar ethanol / acetic-acid charge (about a litre) is held at 350 K in a closed vessel for 600 s and the Fischer esterification runs --- the kinetics of cstr01, now integrated in time by RK4 at 1 s steps. The golden: ethanol falls from 9.25 mol to 1.083 mol (X = 88.3 \%), 8.167 mol of ethyl acetate made, and the jacket has removed only 17.18 kJ to hold the temperature --- exactly -2.10 kJ/mol of reaction at 350 K, the elements-datum heat every reactor in Choupo prices on. The liquid esterification is only mildly exothermic (calorimetry: about -3 to -4 kJ/mol), so the jacket duty here is small, not the cooling load an exotherm of tens of kJ/mol would need.}
\tutreadme{1. choupoBatch, not choupoSolve. A batch vessel has no steady state; the case declares endTime 600; deltaT 1.0; writeInterval 30; in controlDict and the engine integrates dn/dt = $\nu$ r V. The Log tab prints a snapshot every 30 s: four inventories and T. 2. First order means an exponential, and you can check it. With k = 3.575e-3 s?$^1$ at 350 K (the same Arrhenius pair as cstr01), n(600 s) = 9.25$\cdot$e\textasciicircum{}(-2.145) = 1.083 mol. The table is that curve, sampled. 3. The campaign ledger closes at machine precision. [campaign] mass balance ? closure 4.5e-16, energy balance: dH\_vessels = -17.18 kJ ? = Q\_ledger -17.18 kJ, closure 0 --- every joule the jacket took is a state difference on the charge, never a quadrature, which is why the closure is zero and not merely small. 4. A CSTR at the same $\tau$ does worse. 600 s here gives 88 \%; the continuous reactor at $\tau$ = 288 s gave 51 \% and a PFR 64 \%. Three reactors, one rate law, three answers --- the batch one is the upper envelope, because nothing leaves early.}
\tutreadme{Set T\_setpoint 370 K;: k roughly triples, the exponential steepens, and the ledger's Q follows the extent (small, because the reaction is only mildly exothermic). Then batch02\_adiabatic removes the jacket, and T becomes one more state on the trajectory.}
\begin{tutgolden}
\goldrow{kpi}{campaign}{element\_C\_closure\_rel}{3.75075346157e-16}
\goldrow{kpi}{campaign}{element\_H\_closure\_rel}{3.00060276926e-16}
\goldrow{kpi}{campaign}{element\_O\_closure\_rel}{2.50050230771e-16}
\goldrow{kpi}{campaign}{element\_worst\_closure\_rel}{3.75075346157e-16}
\goldrow{kpi}{campaign}{energy\_H\_external\_kJ}{0}
\goldrow{kpi}{campaign}{energy\_Q\_ledger\_kJ}{-17.1846198661}
\goldrow{kpi}{campaign}{energy\_balance\_available}{1}
\goldrow{kpi}{campaign}{energy\_closure\_rel}{0}
\goldrow{kpi}{campaign}{energy\_dH\_vessels\_kJ}{-17.1846198661}
\goldrow{kpi}{campaign}{energy\_records\_invalid}{0}
\goldrow{kpi}{campaign}{energy\_records\_valid}{1}
\goldrow{kpi}{campaign}{energy\_residual\_kJ}{0}
\goldrow{kpi}{campaign}{mass\_closure\_rel}{4.5240474843e-16}
\goldrow{kpi}{campaign}{mass\_kg\_external\_out}{0}
\goldmore{26}
\end{tutgolden}

\clearpage
\tutentry{batch01b\_first\_order\_eulersi}{tutorials/batch/reactor/batch01b\_first\_order\_eulersi}
\begin{tutscheme}
\begin{tikzpicture}[>={Stealth[length=2mm]}, every node/.style={font=\scriptsize}]
  \node[draw, rounded corners=2pt, align=center, minimum width=1.9cm, minimum height=0.85cm, fill=black!4] (reactor) at (0.0,-0.0) {\textbf{reactor}\\[-1pt]{\tiny batchReactor}};
\end{tikzpicture}
\end{tutscheme}
\tutwhat{GREEN WITNESS: semi-implicit Euler (EulerSI) on batch01 mild isothermal esterification -- first order suffices off-stiffness; twin of the ignition01\_h2o2\_eulersi expected failure.}
\begin{tutnarrative}
GREEN WITNESS for the EulerSI integrator: batch01's mild isothermal esterification, integrated by semi-implicit (linearly-implicit) Euler instead of the RK4 default. On a non-stiff problem the first-order method simply works, and the golden pins that its answer agrees with the physics batch01 already pins. Its EXPECTED-FAILURE twin is ignition01\_h2o2\_eulersi (stiff nonlinear ignition, where first order goes unphysical and the engine refuses by name). Registered model no case selected until 2026-08-22 (invariant I8).

\end{tutnarrative}
\begin{tutgolden}
\goldrow{kpi}{campaign}{element\_C\_closure\_rel}{3.75075346157e-16}
\goldrow{kpi}{campaign}{element\_H\_closure\_rel}{4.50090415389e-16}
\goldrow{kpi}{campaign}{element\_O\_closure\_rel}{2.50050230771e-16}
\goldrow{kpi}{campaign}{element\_worst\_closure\_rel}{4.50090415389e-16}
\goldrow{kpi}{campaign}{energy\_H\_external\_kJ}{0}
\goldrow{kpi}{campaign}{energy\_Q\_ledger\_kJ}{-17.1758877049}
\goldrow{kpi}{campaign}{energy\_balance\_available}{1}
\goldrow{kpi}{campaign}{energy\_closure\_rel}{0}
\goldrow{kpi}{campaign}{energy\_dH\_vessels\_kJ}{-17.1758877049}
\goldrow{kpi}{campaign}{energy\_records\_invalid}{0}
\goldrow{kpi}{campaign}{energy\_records\_valid}{1}
\goldrow{kpi}{campaign}{energy\_residual\_kJ}{0}
\goldrow{kpi}{campaign}{mass\_closure\_rel}{5.65505935537e-16}
\goldrow{kpi}{campaign}{mass\_kg\_external\_out}{0}
\goldmore{26}
\end{tutgolden}

\clearpage
\tutentry{batch02\_adiabatic}{tutorials/batch/reactor/batch02\_adiabatic}
\begin{tutscheme}
\begin{tikzpicture}[>={Stealth[length=2mm]}, every node/.style={font=\scriptsize}]
  \node[draw, rounded corners=2pt, align=center, minimum width=1.9cm, minimum height=0.85cm, fill=black!4] (reactor) at (0.0,-0.0) {\textbf{reactor}\\[-1pt]{\tiny batchReactor}};
\end{tikzpicture}
\end{tutscheme}
\tutwhat{Adiabatic batch esterification (AcOH + EtOH $\to$ EtOAc + H2O), starting cold (320 K); the nearly thermoneutral liquid reaction warms the batch only \textasciitilde{}2 K in 600 s and is \textasciitilde{}22 \% complete (thermal runaway needs an exothermic reaction; replacement case planned)}
\begin{tutnarrative}
pseudo-first-order esterification (AcOH + EtOH $\to$ EtOAc + H2O) in an adiabatic batch reactor, starting cold.

T\_0 = 320 K (k\_0 \textasciitilde{} 3.7e-4 s\textasciicircum{}-1, tau\_0 \textasciitilde{} 2700 s). dH\_rxn = -2.43 kJ/mol of ethanol at 320 K, computed from the species' standardThermochemistry (elements datum; the engine's own `[reaction heat]` line). It is about ten times smaller than the -26.06 kJ/mol this case used to print: the acetic-acid correction (docs/design/acetic-acid-latent-heat-closure.md) made the liquid esterification nearly thermoneutral. Equimolar ethanol + acetic acid charge; ethyl acetate + water are the products. Mass-conserving, so the energy balance is clean. Result: T rises only 320 $\to$ 322.2 K in the 600 s window, and the reaction is only \textasciitilde{}22 \% complete at the end (it used to reach 428 K and complete conversion). There is no runaway and no plateau; the rate barely accelerates because T barely moves. The adiabatic ODE integrates H = const on the canonical enthalpy surface (dH\_rxn re-read at the vessel's T at every step: -2.428 kJ/mol at 320 K, -2.416 at 321.1 K), so a restart from a written instant ends on the uninterrupted answer.

The original teaching point of this case (thermal runaway: induction, then a sharp temperature rise, then a plateau) needs a genuinely exothermic reaction; a replacement case on such a reaction is planned. No parameter was changed to recover the runaway.

The energy balance is just:

dT/dt = - r $\ast$ V $\ast$ dH\_rxn / (n\_A $\ast$ Cp\_A + n\_B $\ast$ Cp\_B)

Constant-Cp here (polynomial degree 0 for both ethanol and water in the database), so the result extrapolates linearly above the nominal liquidHeatCapacity Trange (280-373 K) without numerical pathology.

\end{tutnarrative}
\begin{tutgolden}
\goldrow{kpi}{reactor}{P\_final}{101300}
\goldrow{kpi}{reactor}{T\_final}{322.248040531}
\goldrow{kpi}{reactor}{n\_aceticAcid\_final}{0.00722392816212}
\goldrow{kpi}{reactor}{n\_ethanol\_final}{0.00722392816212}
\goldrow{kpi}{reactor}{n\_ethylAcetate\_final}{0.00202607183788}
\goldrow{kpi}{reactor}{n\_water\_final}{0.00202607183788}
\goldrow{kpi}{reactor}{t\_end}{600}
\goldrow{kpi}{reactor}{totalMoles\_final}{0.0185}
\goldrow{kpi}{campaign}{element\_C\_closure\_rel}{3.75075346157e-16}
\goldrow{kpi}{campaign}{element\_H\_closure\_rel}{3.00060276926e-16}
\goldrow{kpi}{campaign}{element\_O\_closure\_rel}{3.75075346157e-16}
\goldrow{kpi}{campaign}{element\_worst\_closure\_rel}{3.75075346157e-16}
\goldrow{kpi}{campaign}{energy\_balance\_available}{0}
\goldrow{kpi}{campaign}{energy\_records\_invalid}{0}
\goldmore{12}
\end{tutgolden}

\clearpage
\tutentry{batch03\_consecutive}{tutorials/batch/reactor/batch03\_consecutive}
\begin{tutscheme}
\begin{tikzpicture}[>={Stealth[length=2mm]}, every node/.style={font=\scriptsize}]
  \node[draw, rounded corners=2pt, align=center, minimum width=1.9cm, minimum height=0.85cm, fill=black!4] (reactor) at (0.0,-0.0) {\textbf{reactor}\\[-1pt]{\tiny batchReactor}};
\end{tikzpicture}
\end{tutscheme}
\tutwhat{Isothermal batch reactor: consecutive A $\to$ B $\to$ C with k1 > k2; intermediate B peaks then decays}
\begin{tutnarrative}
two first-order reactions in series:

with k1 > k2. This is the textbook series-reaction setup: B is produced from A faster than it is consumed to C, so n\_B grows at first, reaches a maximum, then decays as A runs out.

Analytical solution (constant V, isothermal, first-order):

n\_A(t) = n\_A0 $\ast$ exp(-k1 t) n\_B(t) = n\_A0 $\ast$ k1 / (k2 - k1) $\ast$ [exp(-k1 t) - exp(-k2 t)] n\_C(t) = n\_A0 - n\_A(t) - n\_B(t)

t\_max = ln(k1 / k2) / (k1 - k2)

With k1 = 0.01 1/s and k2 = 0.002 1/s:

t\_max = ln(5) / 0.008 \textasciitilde{} 201 s n\_B\_max / n\_A0 = (k1/k2)\textasciicircum{}(k2/(k2-k1)) = 5\textasciicircum{}(-0.25) \textasciitilde{} 0.669

The case demonstrates: 1. multi-reaction support in BatchReactor (v0.21); 2. time-dependent selectivity --- harvesting B at t \textasciitilde{} 200 s maximises yield, waiting longer wastes B to C.

No species supplies activation energy: both A\_pre are quoted with Ea = 0 so that the rate constants are exactly the values above regardless of T (the case is isothermal anyway).

\end{tutnarrative}
\begin{tutgolden}
\goldrow{kpi}{campaign}{energy\_balance\_available}{0}
\goldrow{kpi}{campaign}{energy\_records\_invalid}{1}
\goldrow{kpi}{campaign}{energy\_records\_valid}{0}
\goldrow{kpi}{campaign}{mass\_closure\_rel}{5.55111512313e-16}
\goldrow{kpi}{campaign}{mass\_kg\_external\_out}{0}
\goldrow{kpi}{campaign}{mass\_kg\_final}{0.5}
\goldrow{kpi}{campaign}{mass\_kg\_initial}{0.5}
\goldrow{kpi}{campaign}{mass\_residual\_kg}{2.77555756156e-16}
\goldrow{kpi}{campaign}{transfers\_H\_invalid}{0}
\goldrow{kpi}{campaign}{transfers\_H\_valid}{0}
\goldrow{kpi}{campaign}{transfers\_logged}{0}
\goldrow{kpi}{reactor}{P\_final}{101300}
\goldrow{kpi}{reactor}{T\_final}{350}
\goldrow{kpi}{reactor}{n\_compA\_final}{3.05902320887e-09}
\goldmore{8}
\end{tutgolden}

\clearpage
\tutentry{batch04\_reversible\_wgs}{tutorials/batch/reactor/batch04\_reversible\_wgs}
\begin{tutscheme}
\begin{tikzpicture}[>={Stealth[length=2mm]}, every node/.style={font=\scriptsize}]
  \node[draw, rounded corners=2pt, align=center, minimum width=1.9cm, minimum height=0.85cm, fill=black!4] (reactor) at (0.0,-0.0) {\textbf{reactor}\\[-1pt]{\tiny batchReactor}};
\end{tikzpicture}
\end{tutscheme}
\tutwhat{Batch reversible water-gas shift (CO + H2O <$\to$ CO2 + H2) at 600 K; conversion relaxes to X\_eq}
\begin{tutnarrative}
the reversible water-gas shift in a closed BATCH reactor (choupoBatch), the time-domain sibling of cstr02 / pfr02.

Equimolar CO + H2O charged at 600 K. The forward Arrhenius rate drives conversion up; the reverse term (k\_rev = k\_fwd / K\_eq, K\_eq from the four species' g\_pure\_ig, re-evaluated each step) brakes it. The batch trajectory relaxes onto the SAME thermodynamic ceiling the CSTR and PFR reached --- X $\to$ X\_eq \textasciitilde{} 0.84 --- then sits there (net rate $\to$ 0), instead of running to X = 1.

The forward A is the illustrative value tuned for stable RK4 (see constant/reactions); the equilibrium X\_eq is A-independent.

\end{tutnarrative}
\tutreadme{CO + H2O <$\to$ CO2 + H2, equimolar charge, 600 K, 10 bar (4.989 litres: V = n R T / P for 1 mol), isothermal, reversible Arrhenius kinetics (k\_rev = k\_fwd / K\_eq from the species' ideal-gas Gibbs functions). The conversion relaxes onto the thermodynamic ceiling X\_eq \textasciitilde{} 0.84 (n\_CO 5e-4 $\to$ 7.93e-5 kmol) and stays there.}
\tutreadme{At 600 K and 10 bar all four species are vapours (water's saturation pressure at 600 K is above 10 bar; CO, CO2 and H2 are far above their critical temperatures). The vessel therefore declares energy gasConstantV;, which puts its heat of reaction, its jacket-duty record, its campaign vessel enthalpies and its impulse pricing on the ideal-gas enthalpy surface. Before C448 the case ran without it, and the campaign heat ledger was priced on liquid surfaces for gases (a "liquid" CO or H2 at 600 K), giving -0.299 kJ and then +0.001 kJ for an extent of 0.42 mol of a reaction whose heat is tens of kJ/mol: non-physical both times.}
\tutreadme{dH\_rxn(T) = dH\_rxn(298.15 K) + integral from 298.15 K to T of dCp dT, from the catalogue's own records (data/standards/components):}
\tutreadme{$\ast$ formation enthalpies (J/mol): CO -110530, H2O(g) -241826, CO2 -393520, H2 0, so dH\_rxn(298.15 K) = -393520 - (-110530 - 241826) = -41164 J/mol. $\ast$ ideal-gas Cp polynomials a + bT + cT\textasciicircum{}2 + dT\textasciicircum{}3 (J/mol/K); dCp = Cp(CO2) + Cp(H2) - Cp(CO) - Cp(H2O) has a = -16.173, b = 0.09364, c = -1.08275e-4, d = 4.1111e-8. $\ast$ integrating to 600 K: dH\_rxn(600 K) = -38941.07 J/mol. $\ast$ extent = n\_CO(0) - n\_CO(end) = 0.5 mol - 0.0792842 mol = 0.4207158 mol (from the trajectory below).}
\tutreadme{Q = extent x dH\_rxn(600 K) = 0.4207158 x (-38.94107) = -16.38312 kJ (heat removed by the jacket; the reaction is exothermic).}
\tutreadme{The engine's campaign ledger (energy\_Q\_ledger\_kJ, E\_reaction\_total\_kJ) prints -16.3831222515 kJ: the hand value to the digits above.}
\begin{tutgolden}
\goldrow{kpi}{reactor}{P\_final}{999999.994147}
\goldrow{kpi}{reactor}{T\_final}{600}
\goldrow{kpi}{reactor}{n\_CO2\_final}{0.000420715762123}
\goldrow{kpi}{reactor}{n\_CO\_final}{7.92842378768e-05}
\goldrow{kpi}{reactor}{n\_H2\_final}{0.000420715762123}
\goldrow{kpi}{reactor}{n\_water\_final}{7.92842378768e-05}
\goldrow{kpi}{reactor}{t\_end}{1200}
\goldrow{kpi}{reactor}{totalMoles\_final}{0.001}
\goldrow{kpi}{campaign}{element\_C\_closure\_rel}{2.16840434497e-15}
\goldrow{kpi}{campaign}{element\_H\_closure\_rel}{2.16840434497e-15}
\goldrow{kpi}{campaign}{element\_O\_closure\_rel}{2.16840434497e-15}
\goldrow{kpi}{campaign}{element\_worst\_closure\_rel}{2.16840434497e-15}
\goldrow{kpi}{campaign}{energy\_H\_external\_kJ}{0}
\goldrow{kpi}{campaign}{energy\_Q\_ledger\_kJ}{-16.3831222515}
\goldmore{26}
\end{tutgolden}

\clearpage
\tutentry{batch04\_transient\_dirs}{tutorials/batch/reactor/batch04\_transient\_dirs}
\begin{tutscheme}
\begin{tikzpicture}[>={Stealth[length=2mm]}, every node/.style={font=\scriptsize}]
  \node[draw, rounded corners=2pt, align=center, minimum width=1.9cm, minimum height=0.85cm, fill=black!4] (reactor) at (0.0,-0.0) {\textbf{reactor}\\[-1pt]{\tiny batchReactor}};
\end{tikzpicture}
\end{tutscheme}
\tutwhat{Adiabatic batch esterification with OpenFOAM-style REAL-TIME instant dirs: scrub 0/ 30/ 60/ ... s to watch conversion climb to \textasciitilde{}0.22 while T rises only 320 $\to$ 322 K (the runaway of the old exotherm is gone; replacement case planned)}
\begin{tutnarrative}
the batch02\_adiabatic case re-run with REAL-TIME instant directories: the integrator dumps 0/ 30/ 60/ ... s snapshots (OpenFOAM-style) so the trajectory -- conversion 0 $\to$ \textasciitilde{}0.22, T 320 $\to$ 322.2 K -- can be scrubbed frame by frame in the GUI.

The physics below is inherited from batch02\_adiabatic, unchanged: pseudo-first-order esterification (AcOH + EtOH $\to$ EtOAc + H2O) in an adiabatic batch reactor.

T\_0 = 320 K (k\_0 \textasciitilde{} 3.7e-4 s\textasciicircum{}-1, tau\_0 \textasciitilde{} 2700 s). dH\_rxn = -2.43 kJ/mol of ethanol at 320 K, computed from the species' standardThermochemistry (elements datum; the engine's own `[reaction heat]` line). It is about ten times smaller than the -26.06 kJ/mol this case used to print: the acetic-acid correction (docs/design/acetic-acid-latent-heat-closure.md) made the liquid esterification nearly thermoneutral. Equimolar ethanol + acetic acid charge; ethyl acetate + water are the products. Mass-conserving, so the energy balance is clean. Result: T rises only 320 $\to$ 322.2 K in the 600 s window, and the reaction is only \textasciitilde{}22 \% complete at the end (it used to reach 428 K and complete conversion). There is no runaway and no plateau; the rate barely accelerates because T barely moves. The adiabatic ODE integrates H = const on the canonical enthalpy surface (dH\_rxn re-read at the vessel's T at every step: -2.428 kJ/mol at 320 K, -2.416 at 321.1 K), so a restart from a written instant ends on the uninterrupted answer.

The original teaching point of this case (thermal runaway: induction, then a sharp temperature rise, then a plateau) needs a genuinely exothermic reaction; a replacement case on such a reaction is planned. No parameter was changed to recover the runaway.

The energy balance is just:

dT/dt = - r $\ast$ V $\ast$ dH\_rxn / (n\_A $\ast$ Cp\_A + n\_B $\ast$ Cp\_B)

Constant-Cp here (polynomial degree 0 for both ethanol and water in the database), so the result extrapolates linearly above the nominal liquidHeatCapacity Trange (280-373 K) without numerical pathology.

\end{tutnarrative}
\begin{tutgolden}
\goldrow{kpi}{reactor}{P\_final}{101300}
\goldrow{kpi}{reactor}{T\_final}{322.248040531}
\goldrow{kpi}{reactor}{n\_aceticAcid\_final}{0.00722392816212}
\goldrow{kpi}{reactor}{n\_ethanol\_final}{0.00722392816212}
\goldrow{kpi}{reactor}{n\_ethylAcetate\_final}{0.00202607183788}
\goldrow{kpi}{reactor}{n\_water\_final}{0.00202607183788}
\goldrow{kpi}{reactor}{t\_end}{600}
\goldrow{kpi}{reactor}{totalMoles\_final}{0.0185}
\goldrow{kpi}{campaign}{element\_C\_closure\_rel}{3.75075346157e-16}
\goldrow{kpi}{campaign}{element\_H\_closure\_rel}{3.00060276926e-16}
\goldrow{kpi}{campaign}{element\_O\_closure\_rel}{3.75075346157e-16}
\goldrow{kpi}{campaign}{element\_worst\_closure\_rel}{3.75075346157e-16}
\goldrow{kpi}{campaign}{energy\_balance\_available}{0}
\goldrow{kpi}{campaign}{energy\_records\_invalid}{0}
\goldmore{12}
\end{tutgolden}

\clearpage
\tutentry{batch05\_crystalliser}{tutorials/batch/reactor/batch05\_crystalliser}
\begin{tutscheme}
\begin{tikzpicture}[>={Stealth[length=2mm]}, every node/.style={font=\scriptsize}]
  \node[draw, rounded corners=2pt, align=center, minimum width=1.9cm, minimum height=0.85cm, fill=black!4] (cryst) at (0.0,-0.0) {\textbf{cryst}\\[-1pt]{\tiny batchCrystalliser}};
\end{tikzpicture}
\end{tutscheme}
\tutwhat{Batch crystalliser: supersaturated sugar relaxes, PBE by method of moments}
\begin{tutnarrative}
Single closed vessel. Charge: \textasciitilde{}500 kg water + \textasciitilde{}1250 kg dissolved sucrose (c0 \textasciitilde{} 2.5 kg/kg, above c\_sat(20 C) \textasciitilde{} 2.03 kg/kg, so S0 \textasciitilde{} 1.23). Unseeded: the moments start at zero and primary nucleation (j=0) bootstraps the population from the initial supersaturation. Kinetics from constant/crystallisation.

\end{tutnarrative}
\begin{tutgolden}
\goldrow{kpi}{campaign}{element\_C\_closure\_rel}{6.81588841762e-15}
\goldrow{kpi}{campaign}{element\_H\_closure\_rel}{4.18663266355e-15}
\goldrow{kpi}{campaign}{element\_O\_closure\_rel}{4.18663266355e-15}
\goldrow{kpi}{campaign}{element\_worst\_closure\_rel}{6.81588841762e-15}
\goldrow{kpi}{campaign}{energy\_balance\_available}{1}
\goldrow{kpi}{campaign}{energy\_records\_invalid}{0}
\goldrow{kpi}{campaign}{energy\_records\_valid}{1}
\goldrow{kpi}{campaign}{mass\_closure\_rel}{5.07045231785e-15}
\goldrow{kpi}{campaign}{mass\_kg\_external\_out}{0}
\goldrow{kpi}{campaign}{mass\_kg\_final}{1748.87224776}
\goldrow{kpi}{campaign}{mass\_kg\_initial}{1748.87224776}
\goldrow{kpi}{campaign}{mass\_residual\_kg}{8.86757334229e-12}
\goldrow{kpi}{campaign}{transfers\_H\_invalid}{0}
\goldrow{kpi}{campaign}{transfers\_H\_valid}{0}
\goldmore{22}
\end{tutgolden}

\clearpage
\tutentry{batch06\_adaptive\_runaway}{tutorials/batch/reactor/batch06\_adaptive\_runaway}
\begin{tutscheme}
\begin{tikzpicture}[>={Stealth[length=2mm]}, every node/.style={font=\scriptsize}]
  \node[draw, rounded corners=2pt, align=center, minimum width=1.9cm, minimum height=0.85cm, fill=black!4] (reactor) at (0.0,-0.0) {\textbf{reactor}\\[-1pt]{\tiny batchReactor}};
\end{tikzpicture}
\end{tutscheme}
\tutwhat{Adiabatic esterification (AcOH+EtOH), nearly thermoneutral after the acetic-acid correction; ADAPTIVE Rosenbrock23 time-stepping --- the step grows as the transient settles (glass-box step history; the thermal runaway of the old exotherm is gone, replacement case planned)}
\begin{tutnarrative}
the batch02\_adiabatic esterification integrated with ADAPTIVE Rosenbrock23 time-stepping, the step history printed glass-box so the student sees the controller at work. (The directory name keeps "runaway" from the old exotherm; there is no runaway any more.)

The physics below is inherited from batch02\_adiabatic, unchanged: pseudo-first-order esterification (AcOH + EtOH $\to$ EtOAc + H2O) in an adiabatic batch reactor.

T\_0 = 320 K (k\_0 \textasciitilde{} 3.7e-4 s\textasciicircum{}-1, tau\_0 \textasciitilde{} 2700 s). dH\_rxn = -2.43 kJ/mol of ethanol at 320 K, computed from the species' standardThermochemistry (elements datum; the engine's own `[reaction heat]` line). It is about ten times smaller than the -26.06 kJ/mol this case used to print: the acetic-acid correction (docs/design/acetic-acid-latent-heat-closure.md) made the liquid esterification nearly thermoneutral. Equimolar ethanol + acetic acid charge; ethyl acetate + water are the products. Mass-conserving, so the energy balance is clean. Result: T rises only 320 $\to$ 322.2 K in the 600 s window, and the reaction is only \textasciitilde{}22 \% complete at the end (it used to reach 428 K and complete conversion). There is no runaway and no plateau; the rate barely accelerates because T barely moves. The adiabatic ODE integrates H = const on the canonical enthalpy surface (dH\_rxn re-read at the vessel's T at every step: -2.428 kJ/mol at 320 K, -2.416 at 321.1 K), so a restart from a written instant ends on the uninterrupted answer.

The original teaching point of this case (thermal runaway: induction, then a sharp temperature rise, then a plateau) needs a genuinely exothermic reaction; a replacement case on such a reaction is planned. No parameter was changed to recover the runaway.

The energy balance is just:

dT/dt = - r $\ast$ V $\ast$ dH\_rxn / (n\_A $\ast$ Cp\_A + n\_B $\ast$ Cp\_B)

Constant-Cp here (polynomial degree 0 for both ethanol and water in the database), so the result extrapolates linearly above the nominal liquidHeatCapacity Trange (280-373 K) without numerical pathology.

\end{tutnarrative}
\begin{tutgolden}
\goldrow{kpi}{reactor}{P\_final}{101300}
\goldrow{kpi}{reactor}{T\_final}{322.248042017}
\goldrow{kpi}{reactor}{n\_aceticAcid\_final}{0.00722392679638}
\goldrow{kpi}{reactor}{n\_ethanol\_final}{0.00722392679638}
\goldrow{kpi}{reactor}{n\_ethylAcetate\_final}{0.00202607320362}
\goldrow{kpi}{reactor}{n\_water\_final}{0.00202607320362}
\goldrow{kpi}{reactor}{t\_end}{600}
\goldrow{kpi}{reactor}{totalMoles\_final}{0.0185}
\goldrow{kpi}{campaign}{element\_C\_closure\_rel}{3.75075346157e-16}
\goldrow{kpi}{campaign}{element\_H\_closure\_rel}{3.00060276926e-16}
\goldrow{kpi}{campaign}{element\_O\_closure\_rel}{2.50050230771e-16}
\goldrow{kpi}{campaign}{element\_worst\_closure\_rel}{3.75075346157e-16}
\goldrow{kpi}{campaign}{energy\_balance\_available}{0}
\goldrow{kpi}{campaign}{energy\_records\_invalid}{0}
\goldmore{12}
\end{tutgolden}

\clearpage
\tutentry{batch07\_lhhw\_methylAcetate}{tutorials/batch/reactor/batch07\_lhhw\_methylAcetate}
\begin{tutscheme}
\begin{tikzpicture}[>={Stealth[length=2mm]}, every node/.style={font=\scriptsize}]
  \node[draw, rounded corners=2pt, align=center, minimum width=1.9cm, minimum height=0.85cm, fill=black!4] (reactor) at (0.0,-0.0) {\textbf{reactor}\\[-1pt]{\tiny batchReactor}};
\end{tikzpicture}
\end{tutscheme}
\tutwhat{LHHW in a BATCH reactor: methyl-acetate synthesis over Amberlyst 15 (Popken 2000, Eq 16)}
\begin{tutnarrative}
the SAME rate law as cstr07, in a BATCH vessel.

CH3OH + CH3COOH <$\to$ CH3COOCH3 + H2O over Amberlyst 15

Popken, Gotze \& Gmehling, Ind. Eng. Chem. Res. 39 (2000) 2601, Eq 16: the adsorption-based LHHW on UNIQUAC activities. Nothing about the rate law is restated here -- it lives in constant/reactions, exactly as the steady CSTR reads it, because the SAME `RateLaw` evaluates it in both. A batch reactor and a continuous one cannot disagree about what an adsorption denominator means.

What changes is only what a batch vessel does with a rate: it integrates it. The charge is equimolar methanol + acetic acid; the trajectory climbs to the ACTIVITY equilibrium, X = 76.6 \%, and stops. That plateau is asserted nowhere: it is where k1 a'\_HOAc a'\_MeOH = k-1 a'\_MeOAc a'\_H2O, i.e. the K\_a = 25.50 the authors' independent Eq 17 recovers from the same eight constants -- and the same number the steady CSTR climbs to as its volume grows. Two reactors, two integration schemes, one equilibrium.

Water is the strongest adsorber (K/M = 0.29 against 0.052 for the acid), so as the reaction MAKES water it crowds the acid off the resin and the denominator grows. That self-inhibition is why reactive distillation, which pulls the water out as it forms, is worth its capital.

Do NOT try to see the inhibition by deleting the `adsorption` block and keeping k1. It tells you nothing: the constants were REGRESSED for this rate law, on these activities, and a rate constant without its rate law is a number without a meaning. (Delete the denominator and the reaction gets FASTER, because the adsorbed activities a' are small and dividing by their square is a multiplication. That is an arithmetic accident, not physics.) To SEE the inhibition, spike the charge with water at the same acid and alcohol fractions:

methanol 0.4, aceticAcid 0.4, water 0.0 $\to$ initial rate 58.3 mmol/(L.s) methanol 0.4, aceticAcid 0.4, water 0.2 $\to$ initial rate 29.0 mmol/(L.s)

Half the rate, with every reactant fraction unchanged. The water is not a reactant here; it is sitting on the sites.

`catalystLoading` (kg of dry resin per m\textasciicircum{}3) is what turns the paper's per-gram rate constants into a volumetric rate. With concentrations in kmol/m\textasciicircum{}3 the conversion is exact and unit-free: mol/(g.s) x kg/m\textasciicircum{}3 = kmol/(m\textasciicircum{}3.s).

The vessel runs isothermal at 333.15 K, the top of the authors' data range. Do not extrapolate: their activation energies are APPARENT values, fitted over 303-333 K.

\end{tutnarrative}
\begin{tutgolden}
\goldrow{kpi}{reactor}{P\_final}{101300}
\goldrow{kpi}{reactor}{T\_final}{333.15}
\goldrow{kpi}{reactor}{n\_aceticAcid\_final}{0.00272280695052}
\goldrow{kpi}{reactor}{n\_methanol\_final}{0.00272280695052}
\goldrow{kpi}{reactor}{n\_methylAcetate\_final}{0.00652719304948}
\goldrow{kpi}{reactor}{n\_water\_final}{0.00652719304948}
\goldrow{kpi}{reactor}{t\_end}{3600}
\goldrow{kpi}{reactor}{totalMoles\_final}{0.0185}
\goldrow{kpi}{campaign}{element\_C\_closure\_rel}{2.12542696156e-15}
\goldrow{kpi}{campaign}{element\_H\_closure\_rel}{2.06291440386e-15}
\goldrow{kpi}{campaign}{element\_O\_closure\_rel}{2.00040184617e-15}
\goldrow{kpi}{campaign}{element\_worst\_closure\_rel}{2.12542696156e-15}
\goldrow{kpi}{campaign}{energy\_H\_external\_kJ}{0}
\goldrow{kpi}{campaign}{energy\_Q\_ledger\_kJ}{-26.3966571762}
\goldmore{26}
\end{tutgolden}

\clearpage
\tutentry{batch08\_isothermal\_duty}{tutorials/batch/reactor/batch08\_isothermal\_duty}
\begin{tutscheme}
\begin{tikzpicture}[>={Stealth[length=2mm]}, every node/.style={font=\scriptsize}]
  \node[draw, rounded corners=2pt, align=center, minimum width=1.9cm, minimum height=0.85cm, fill=black!4] (reactor) at (0.0,-0.0) {\textbf{reactor}\\[-1pt]{\tiny batchReactor}};
\end{tikzpicture}
\end{tutscheme}
\tutwhat{Isothermal jacket duty from a DECLARED dH\_rxn: E = dH$\ast$dXi via extents (lumped species, no formation data)}
\begin{tutnarrative}
one isothermal batch reactor, one declared dH\_rxn, one hand-checkable jacket duty. See system/controlDict for what this case tests and constant/reactions for the analytic anchor (E = -50 kJ/mol x 0.9975 mol = -49.88 kJ over 600 s).

\end{tutnarrative}
\begin{tutgolden}
\goldrow{kpi}{campaign}{energy\_balance\_available}{0}
\goldrow{kpi}{campaign}{energy\_records\_invalid}{0}
\goldrow{kpi}{campaign}{energy\_records\_valid}{1}
\goldrow{kpi}{campaign}{mass\_closure\_rel}{1.11022302463e-15}
\goldrow{kpi}{campaign}{mass\_kg\_external\_out}{0}
\goldrow{kpi}{campaign}{mass\_kg\_final}{0.05}
\goldrow{kpi}{campaign}{mass\_kg\_initial}{0.05}
\goldrow{kpi}{campaign}{mass\_residual\_kg}{5.55111512313e-17}
\goldrow{kpi}{campaign}{transfers\_H\_invalid}{0}
\goldrow{kpi}{campaign}{transfers\_H\_valid}{0}
\goldrow{kpi}{campaign}{transfers\_logged}{0}
\goldrow{kpi}{reactor}{E\_reaction\_seg0\_kJ}{-49.8760623911}
\goldrow{kpi}{reactor}{E\_reaction\_total\_kJ}{-49.8760623911}
\goldrow{kpi}{reactor}{P\_final}{101300}
\goldmore{13}
\end{tutgolden}

\clearpage
\tutentry{batch26\_cooling\_failure}{tutorials/batch/reactor/batch26\_cooling\_failure}
\begin{tutscheme}
\begin{tikzpicture}[>={Stealth[length=2mm]}, every node/.style={font=\scriptsize}]
  \node[draw, rounded corners=2pt, align=center, minimum width=1.9cm, minimum height=0.85cm, fill=black!4] (reactor) at (0.0,-0.0) {\textbf{reactor}\\[-1pt]{\tiny batchReactor}};
\end{tikzpicture}
\end{tutscheme}
\tutwhat{Cooling failure on a batch (fictitious A + B $\to$ P in solvent S): isothermal at 330 K, the jacket stops at t\_f, the accumulated B drives T to T\_cf}
\begin{tutnarrative}
a BATCH exothermic reaction held at T\_p by its jacket, and what happens when the jacket stops.

THE REACTION IS FICTITIOUS: A + B $\to$ P in an inert solvent S, an ideal liquid with no vapour phase in the reactor (constant/reactions and constant/components/reactantA.dat say how every number was set). It is a teaching witness for the process-safety EduTool, not a model of any real chemistry, and its records are coherent by construction: the heat of reaction, -60.000 kJ per mol of B, comes from formation data whose differences are the only thing that carries meaning, and it does not depend on temperature (dCp = 0). The whole charge -- 5 mol A, 5 mol B, 15 mol S -- is in the vessel at t = 0. That is what makes a batch the WORST case for thermal safety: at t = 0 all of the reactant is unconverted, so a cooling failure at the start releases the whole adiabatic temperature rise.

THE SCENARIO (the recipe below):

t < t\_f isothermal at T\_p = 330 K; the jacket removes the heat of reaction (the energy ledger records that duty) t = t\_f `setParameter reactor adiabatic 1` -- the jacket stops t > t\_f adiabatic: the B still unconverted at t\_f (the ACCUMULATION, X\_ac) reacts and heats the mass to T\_cf

The thermal-safety estimate is T\_cf = T\_p + X\_ac $\ast$ dT\_ad (a batch's mass does not change, so Stoessel's semi-batch mass correction is 1), with dT\_ad the adiabatic rise of the full charge from T\_p (run this case with the failure at t = 0 to have the engine compute it). Because dCp = 0 here, the total heat capacity of the charge does not change as it reacts, and the estimate is EXACT for this batch: the engine's T\_cf (T\_final of a failure run) and the estimate agree to the integration error. The MTSR is the largest T\_cf over every failure instant: for a batch, the failure at t = 0.

`value 0` on the same event keeps the jacket working: the same case then runs the normal isothermal batch, which is how the EduTool draws X\_ac(t).

WHAT IS NOT MODELLED: the vessel has no vapour phase -- past the mixture's boiling point a real open vessel would boil and a closed one would pressurise. The MTSR is defined on exactly that hypothetical, which is why it is compared with the boiling point (MTT) rather than trusted beyond it. No decomposition is modelled either: T\_D24 is an assumption the EduTool's reader sets.

\end{tutnarrative}
\begin{tutgolden}
\goldrow{kpi}{campaign}{energy\_balance\_available}{0}
\goldrow{kpi}{campaign}{energy\_records\_invalid}{0}
\goldrow{kpi}{campaign}{energy\_records\_valid}{1}
\goldrow{kpi}{campaign}{mass\_closure\_rel}{1.97372982156e-15}
\goldrow{kpi}{campaign}{mass\_kg\_external\_out}{0}
\goldrow{kpi}{campaign}{mass\_kg\_final}{2.25}
\goldrow{kpi}{campaign}{mass\_kg\_initial}{2.25}
\goldrow{kpi}{campaign}{mass\_residual\_kg}{4.4408920985e-15}
\goldrow{kpi}{campaign}{moles\_kmol\_external\_out}{0}
\goldrow{kpi}{campaign}{moles\_kmol\_final}{0.02}
\goldrow{kpi}{campaign}{moles\_kmol\_initial}{0.025}
\goldrow{kpi}{campaign}{transfers\_H\_invalid}{0}
\goldrow{kpi}{campaign}{transfers\_H\_valid}{0}
\goldrow{kpi}{campaign}{transfers\_logged}{0}
\goldmore{17}
\end{tutgolden}

\clearpage
\tutentry{batch27\_semibatch\_dosing}{tutorials/batch/reactor/batch27\_semibatch\_dosing}
\begin{tutscheme}
\begin{tikzpicture}[>={Stealth[length=2mm]}, every node/.style={font=\scriptsize}]
  \node[draw, rounded corners=2pt, align=center, minimum width=1.9cm, minimum height=0.85cm, fill=black!4] (reactor) at (0.0,-0.0) {\textbf{reactor}\\[-1pt]{\tiny batchReactor}};
  \node[draw, rounded corners=2pt, align=center, minimum width=1.9cm, minimum height=0.85cm, fill=black!4] (dosingTank) at (0.0,-1.7) {\textbf{dosingTank}\\[-1pt]{\tiny batchAccumulator}};
  \node[draw, rounded corners=2pt, align=center, minimum width=1.9cm, minimum height=0.85cm, fill=black!4] (heldFeed) at (0.0,-3.4) {\textbf{heldFeed}\\[-1pt]{\tiny batchAccumulator}};
\end{tikzpicture}
\end{tutscheme}
\tutwhat{Semi-batch (fictitious A + B $\to$ P in solvent S): B dosed in 16 portions over 4 h at 330 K, the jacket fails after the last portion, the small accumulation keeps T\_cf low}
\begin{tutnarrative}
the SAME fictitious reaction as batch26 (A + B $\to$ P in the inert solvent S), the same totals, the same T\_p, run SEMI-BATCH: A and the solvent are in the reactor and B is DOSED, a sixteenth at a time every 15 minutes, over four hours. Then the jacket fails.

This is the inherently safer answer to batch26 in Kletz's sense of limiting what is there: a batch reactor holds all of its reactant at t = 0; a dosed reactor holds only what it has been given and has not yet converted. That unconverted inventory -- the ACCUMULATION -- is what a cooling failure turns into temperature, so limiting it limits the maximum temperature of the synthesis reaction (MTSR). The lever is the DOSING RATE against the REACTION RATE: dose faster than the reaction consumes and the reactor fills with unreacted B, which is a batch again.

THE DOSING is written as portions, each a recipe `transfer` from the dosing tank, because a recipe transfer is what choupoBatch has: a batchReactor carries no continuous feed. Portion-wise addition is also how a hazardous reagent is often added in practice. Each transfer moves a FRACTION of what is left in the tank (1/16, then 1/15, ... then 1), so every portion is 1/16 of B. The dosing tank is at T\_p, so the feed brings no sensible-heat effect into the isothermal reactor (into a hot adiabatic one the dose is mixed by enthalpy equality and cools it a little; a cold feed would take more of the reaction's heat -- left out on purpose). With discrete doses the accumulation jumps at each portion, so the worst instant for a failure is RIGHT AFTER a portion; the trajectory row at a portion's instant is written after the transfer.

THE TWO SAFETY EVENTS come first in the recipe:

1. the COOLING FAILURE at t\_f (`adiabatic 1`; `0` keeps the jacket); 2. what happens to the FEED at the same instant, one transfer with three readings: `to heldFeed; fraction 1;` the DOSING INTERLOCK works: the methodology's assumption that the feed stops when the cooling does. A recipe transfer has no valve, so the interlock is written as what it achieves -- everything left in the dosing tank is moved to `heldFeed`, a vessel nothing is ever dosed from, and every later portion moves a fraction of an empty tank, which is nothing. `fraction 0;` the interlock FAILS: nothing is held back and the pump keeps dosing into an adiabatic reactor, at its rate. `to reactor; fraction 1;` the WHOLE dosing tank enters the reactor at once (a valve left open, a pump run away): the worst a feed failure can do, bounded only by what the tank holds. A flow-limited pump limits the RATE, not the TOTAL: with the whole of B in the tank, that total is the batch's. In a plant an interlock is a sensor, a logic solver and a valve, each of which can fail; here it is one line, and the EduTool runs it all three ways. batch28\_dosing\_vessel\_limit is the same reactor with a dosing vessel that holds only a quarter of B.

Normal operation, for the accumulation curve: the failure at `value 0` and the feed event at `fraction 0`.

The shipped t\_f is 1 s after the last portion, where the accumulation is close to its largest -- NOT where the temperature after a failure is largest: early in the dosing the accumulation heats a smaller mass (Stoessel's M\_f/M(t) correction), so the MTSR comes earlier; the EduTool shows where, and confirms it with the engine at every post-portion instant.

WHAT IS NOT MODELLED, besides what batch26 names: the reactor's volume is held at its declared value while the dosing adds liquid (a batchReactor does not grow its volume on a charge). The rate law is first order in B and zero order in A, so the extent rate k n\_B does not depend on V and this changes no number here; a rate law with a concentration term of another order would be wrong.

\end{tutnarrative}
\begin{tutgolden}
\goldrow{kpi}{campaign}{energy\_balance\_available}{0}
\goldrow{kpi}{campaign}{energy\_records\_invalid}{0}
\goldrow{kpi}{campaign}{energy\_records\_valid}{16}
\goldrow{kpi}{campaign}{mass\_closure\_rel}{5.92118946467e-16}
\goldrow{kpi}{campaign}{mass\_kg\_external\_out}{0}
\goldrow{kpi}{campaign}{mass\_kg\_final}{2.25}
\goldrow{kpi}{campaign}{mass\_kg\_initial}{2.25}
\goldrow{kpi}{campaign}{mass\_residual\_kg}{1.33226762955e-15}
\goldrow{kpi}{campaign}{moles\_kmol\_external\_out}{0}
\goldrow{kpi}{campaign}{moles\_kmol\_final}{0.0200009153077}
\goldrow{kpi}{campaign}{moles\_kmol\_initial}{0.025}
\goldrow{kpi}{campaign}{transfers\_H\_invalid}{0}
\goldrow{kpi}{campaign}{transfers\_H\_valid}{17}
\goldrow{kpi}{campaign}{transfers\_logged}{17}
\goldmore{53}
\end{tutgolden}

\clearpage
\tutentry{batch28\_dosing\_vessel\_limit}{tutorials/batch/reactor/batch28\_dosing\_vessel\_limit}
\begin{tutscheme}
\begin{tikzpicture}[>={Stealth[length=2mm]}, every node/.style={font=\scriptsize}]
  \node[draw, rounded corners=2pt, align=center, minimum width=1.9cm, minimum height=0.85cm, fill=black!4] (reactor) at (0.0,-0.0) {\textbf{reactor}\\[-1pt]{\tiny batchReactor}};
  \node[draw, rounded corners=2pt, align=center, minimum width=1.9cm, minimum height=0.85cm, fill=black!4] (dayTank) at (0.0,-1.7) {\textbf{dayTank}\\[-1pt]{\tiny batchAccumulator}};
  \node[draw, rounded corners=2pt, align=center, minimum width=1.9cm, minimum height=0.85cm, fill=black!4] (storage) at (0.0,-3.4) {\textbf{storage}\\[-1pt]{\tiny batchAccumulator}};
  \node[draw, rounded corners=2pt, align=center, minimum width=1.9cm, minimum height=0.85cm, fill=black!4] (heldFeed) at (0.0,-5.1) {\textbf{heldFeed}\\[-1pt]{\tiny batchAccumulator}};
\end{tikzpicture}
\end{tutscheme}
\tutwhat{Semi-batch with a day tank that holds a quarter of B (fictitious A + B $\to$ P in solvent S): the cooling fails and the whole day tank enters the reactor at once; T\_cf stays low}
\begin{tutnarrative}
batch27 with the DOSING VESSEL SIZED so that a feed failure cannot take the reactor past its limits.

The same fictitious reaction (A + B $\to$ P in the inert solvent S), the same totals, the same T\_p and the same sixteen portions, one every 15 minutes. What changes is where B waits. In batch27 the dosing tank holds ALL of B, so the worst a feed failure can do -- the whole tank entering at once, or the pump running on with no interlock to stop it -- delivers the batch's charge: a flow-limited pump limits the RATE, not the TOTAL. Here the only vessel connected to the reactor is a DAY TANK that holds a QUARTER of B (four portions). It is refilled from storage three times, each refill made only once the day tank is empty, between two portions; storage has no line to the reactor.

THE FAILURE SCENARIO (the first three recipe events, all at t\_f):

1. the COOLING FAILURE (`adiabatic 1`; `0` keeps the jacket); 2. the WORST FEED FAILURE: the WHOLE day tank enters the reactor at once (`to reactor; fraction 1;`). `to heldFeed; fraction 1;` is the interlock holding it back, `fraction 0;` the pump running on through its remaining portions -- which delivers the same total, only later; 3. NO REFILL after a failure: what is left in storage is moved to `heldFeed` (`fraction 1`; `fraction 0` in normal operation).

Event 3 is the honest limit of the claim. The day tank's size removes the dependence on the DOSING interlock; what remains is the refill, which is a discrete, deliberate act taken when the day tank is empty and the reactor's state is known -- a procedure, or an interlock of its own. The page says so; it does not hide it inside the word "inherent".

With the whole day tank entering, the worst instant is right after a refill (or at t = 0), when the tank is full: the B that can reach the reactor is then the accumulation plus a quarter of the charge. The shipped t\_f is 1 s after the first refill.

WHAT IS NOT MODELLED: as batch26 and batch27 (no vapour phase, no decomposition, the reactor's volume held at its declared value -- the rate is first order in B, so V does not enter the extent rate).

\end{tutnarrative}
\begin{tutgolden}
\goldrow{kpi}{campaign}{energy\_balance\_available}{0}
\goldrow{kpi}{campaign}{energy\_records\_invalid}{0}
\goldrow{kpi}{campaign}{energy\_records\_valid}{4}
\goldrow{kpi}{campaign}{mass\_closure\_rel}{1.7763568394e-15}
\goldrow{kpi}{campaign}{mass\_kg\_external\_out}{0}
\goldrow{kpi}{campaign}{mass\_kg\_final}{2.25}
\goldrow{kpi}{campaign}{mass\_kg\_initial}{2.25}
\goldrow{kpi}{campaign}{mass\_residual\_kg}{3.99680288865e-15}
\goldrow{kpi}{campaign}{moles\_kmol\_external\_out}{0}
\goldrow{kpi}{campaign}{moles\_kmol\_final}{0.0225}
\goldrow{kpi}{campaign}{moles\_kmol\_initial}{0.025}
\goldrow{kpi}{campaign}{transfers\_H\_invalid}{0}
\goldrow{kpi}{campaign}{transfers\_H\_valid}{21}
\goldrow{kpi}{campaign}{transfers\_logged}{21}
\goldmore{53}
\end{tutgolden}

\clearpage
\tutentry{batch30\_fast\_equilibrium\_full}{tutorials/batch/reactor/batch30\_fast\_equilibrium\_full}
\begin{tutscheme}
\begin{tikzpicture}[>={Stealth[length=2mm]}, every node/.style={font=\scriptsize}]
  \node[draw, rounded corners=2pt, align=center, minimum width=1.9cm, minimum height=0.85cm, fill=black!4] (reactor) at (0.0,-0.0) {\textbf{reactor}\\[-1pt]{\tiny batchReactor}};
\end{tikzpicture}
\end{tutscheme}
\tutwhat{A fast reversible step feeding a slow one, A <$\to$ B fast and B $\to$ C slow, ALL THREE rates integrated: the full model batch31 is compared with}
\begin{tutnarrative}
a fast reversible step feeding a slow one, every rate integrated:

A <-$\to$ B kf = 0.1 1/s, kr = 0.025 1/s (K = kf/kr = 4, fast) B --$\to$ C k = 0.001 1/s (slow)

Isothermal at 350 K, first order throughout, Ea = 0 so each rate constant is exactly the number written in constant/reactions.

THE QUESTION. A and B interconvert a hundred times faster than B leaves for C. After a short first moment they sit at equilibrium with each other, n\_B = K n\_A, and the pair behaves as ONE pool that drains into C at the rate of the B it contains:

d(n\_A + n\_B)/dt = -k n\_B = -k K/(1+K) (n\_A + n\_B)

That is partial equilibrium, and batch31\_fast\_equilibrium\_reduced runs exactly that law, on one pool, with no kf and no kr. The two answers for n\_C differ by the delay of the first moment, while B is still forming: measured with these numbers the largest gap is 0.62 \% of the charge, and it falls as kf/k grows. The EduTool page "When the fast obeys the slow" changes kf and kr together (K held at 4) and draws the gap.

THE COST. The fast pair sets a rate constant of kf + kr = 0.125 1/s. An explicit fixed RK4 step would have to stay below 2.785/0.125 = 22 s, and below 0.0028 s at the page's largest ratio; the adaptive Rosenbrock23 integrator declared in controlDict prints every step it takes.

No species carries formation data and none is needed: isothermal mode holds T, so the heat of reaction is not used (the run says so), the same posture as batch03\_consecutive.

\end{tutnarrative}
\tutreadme{An isothermal batch reactor with A <$\to$ B (kf = 0.1, kr = 0.025 1/s, K = 4) and B $\to$ C (k = 0.001 1/s), every rate integrated by the adaptive Rosenbrock23 integrator. 10 mol of A, 6000 s.}
\tutreadme{1. The first moment. A and B reach equilibrium with each other within tens of seconds; after that n\_B / n\_A stays at 4 while both drain to C. 2. The reduced twin. batch31\_fast\_equilibrium\_reduced runs one pool decaying at k K/(1+K). The largest gap in n\_C between the two is 0.62 \% of the charge, and it shrinks as kf/k grows. 3. The step history. At verbosity 3 the log prints every step the integrator accepts; count them against the reduced twin's.}
\begin{tutgolden}
\goldrow{kpi}{campaign}{energy\_balance\_available}{0}
\goldrow{kpi}{campaign}{energy\_records\_invalid}{1}
\goldrow{kpi}{campaign}{energy\_records\_valid}{0}
\goldrow{kpi}{campaign}{mass\_closure\_rel}{3.66373598126e-15}
\goldrow{kpi}{campaign}{mass\_kg\_external\_out}{0}
\goldrow{kpi}{campaign}{mass\_kg\_final}{0.5}
\goldrow{kpi}{campaign}{mass\_kg\_initial}{0.5}
\goldrow{kpi}{campaign}{mass\_residual\_kg}{1.83186799063e-15}
\goldrow{kpi}{campaign}{moles\_kmol\_external\_out}{0}
\goldrow{kpi}{campaign}{moles\_kmol\_final}{0.01}
\goldrow{kpi}{campaign}{moles\_kmol\_initial}{0.01}
\goldrow{kpi}{campaign}{transfers\_H\_invalid}{0}
\goldrow{kpi}{campaign}{transfers\_H\_valid}{0}
\goldrow{kpi}{campaign}{transfers\_logged}{0}
\goldmore{8}
\end{tutgolden}

\clearpage
\tutentry{batch31\_fast\_equilibrium\_reduced}{tutorials/batch/reactor/batch31\_fast\_equilibrium\_reduced}
\begin{tutscheme}
\begin{tikzpicture}[>={Stealth[length=2mm]}, every node/.style={font=\scriptsize}]
  \node[draw, rounded corners=2pt, align=center, minimum width=1.9cm, minimum height=0.85cm, fill=black!4] (reactor) at (0.0,-0.0) {\textbf{reactor}\\[-1pt]{\tiny batchReactor}};
\end{tikzpicture}
\end{tutscheme}
\tutwhat{The same reactor reduced by partial equilibrium: A and B as ONE pool decaying at k K/(1+K), written as one expression law}
\begin{tutnarrative}
the REDUCED twin of batch30\_fast\_equilibrium\_full.

THE REDUCTION. When A <$\to$ B is fast against B $\to$ C, A and B stay at equilibrium with each other, n\_B = K n\_A, and only their sum moves:

d(pool)/dt = -k K/(1+K) pool, pool = n\_A + n\_B

So this case runs ONE reaction, pool $\to$ C, at that rate, written as an `expression` law in constant/reactions. Here compA STANDS FOR THE POOL: at t = 0 it holds the whole charge, and at any time the share that is B is K/(1+K) = 0.8 of it. compB stays declared (the two cases share one component list) and stays at zero.

WHAT IS KEPT. Two numbers, k and K. The fast rate constants kf and kr do not appear at all: any pair with the same ratio gives this answer. That is the reduction's real gain -- it needs one parameter fewer, and it is the parameter that is hardest to measure. What is lost is the first moment, while B forms from A: batch30 measures it.

\end{tutnarrative}
\tutreadme{The reduced twin of batch30\_fast\_equilibrium\_full. A and B are lumped into one pool (carried as compA) that drains to C at k K/(1+K), written in constant/reactions as a type expression law. No kf and no kr appear.}
\tutreadme{1. Two parameters instead of three. k and K are all the law needs. 2. What is lost. The moment while B forms from A: against the full twin the largest gap in n\_C is 0.62 \% of the charge. 3. compA is the pool, not A alone: 0.8 of it is B at every instant.}
\begin{tutgolden}
\goldrow{kpi}{campaign}{energy\_balance\_available}{0}
\goldrow{kpi}{campaign}{energy\_records\_invalid}{1}
\goldrow{kpi}{campaign}{energy\_records\_valid}{0}
\goldrow{kpi}{campaign}{mass\_closure\_rel}{1.11022302463e-15}
\goldrow{kpi}{campaign}{mass\_kg\_external\_out}{0}
\goldrow{kpi}{campaign}{mass\_kg\_final}{0.5}
\goldrow{kpi}{campaign}{mass\_kg\_initial}{0.5}
\goldrow{kpi}{campaign}{mass\_residual\_kg}{5.55111512313e-16}
\goldrow{kpi}{campaign}{moles\_kmol\_external\_out}{0}
\goldrow{kpi}{campaign}{moles\_kmol\_final}{0.01}
\goldrow{kpi}{campaign}{moles\_kmol\_initial}{0.01}
\goldrow{kpi}{campaign}{transfers\_H\_invalid}{0}
\goldrow{kpi}{campaign}{transfers\_H\_valid}{0}
\goldrow{kpi}{campaign}{transfers\_logged}{0}
\goldmore{8}
\end{tutgolden}

\clearpage
\tutentry{batch32\_fast\_intermediate\_full}{tutorials/batch/reactor/batch32\_fast\_intermediate\_full}
\begin{tutscheme}
\begin{tikzpicture}[>={Stealth[length=2mm]}, every node/.style={font=\scriptsize}]
  \node[draw, rounded corners=2pt, align=center, minimum width=1.9cm, minimum height=0.85cm, fill=black!4] (reactor) at (0.0,-0.0) {\textbf{reactor}\\[-1pt]{\tiny batchReactor}};
\end{tikzpicture}
\end{tutscheme}
\tutwhat{A $\to$ I $\to$ P with the intermediate I consumed fast, BOTH rates integrated: the full model batch33 is compared with}
\begin{tutnarrative}
a reactive intermediate, every rate integrated:

A --k1-$\to$ I --k2-$\to$ P k1 = 0.001 1/s, k2 = 0.1 1/s

(I is compB and P is compC.) Isothermal at 350 K, first order, Ea = 0.

THE QUESTION. I is used up a hundred times faster than it is made, so it never accumulates: after a short first moment its amount is whatever makes its two rates equal, n\_I = (k1/k2) n\_A, small and following A. Then P is made as fast as A is consumed:

That is the quasi-steady intermediate, and batch33\_fast\_intermediate\_reduced runs exactly that law, A $\to$ P, with no k2. The two answers for n\_P differ by the I that is in transit: measured with these numbers the largest gap is 0.95 \% of the charge. For a closed system the difference in P is the amount of I itself, so the gap and the largest n\_I are the same number. The EduTool page changes k2 and draws it.

THE SAME EQUATION AS TWO TANKS. The gap here, at k1/k2 = eps, is the gap between two mixing tanks in series at tau\_f/tau\_s = eps (unsteady07/unsteady08): eps\textasciicircum{}(1/(1-eps)) in both, because both are a fast first-order system driven by a slow one.

No species carries formation data and none is needed: isothermal mode holds T, so the heat of reaction is not used, as in batch03\_consecutive.

\end{tutnarrative}
\tutreadme{An isothermal batch reactor with A $\to$ I $\to$ P (k1 = 0.001, k2 = 0.1 1/s; I is compB, P is compC), both rates integrated. 10 mol of A, 5000 s.}
\tutreadme{1. I never accumulates. After a short first moment n\_I stays near (k1/k2) n\_A, under 1 \% of the charge. 2. The reduced twin. batch33\_fast\_intermediate\_reduced makes P at k1 n\_A directly. The largest gap in n\_P is 0.95 \% of the charge, and it equals the largest n\_I: in a closed vessel the missing P is the I in transit. 3. Two tanks in disguise. The gap at k1/k2 = eps is the gap of unsteady07/unsteady08 at the same eps: the same equation.}
\begin{tutgolden}
\goldrow{kpi}{campaign}{energy\_balance\_available}{0}
\goldrow{kpi}{campaign}{energy\_records\_invalid}{1}
\goldrow{kpi}{campaign}{energy\_records\_valid}{0}
\goldrow{kpi}{campaign}{mass\_closure\_rel}{4.4408920985e-16}
\goldrow{kpi}{campaign}{mass\_kg\_external\_out}{0}
\goldrow{kpi}{campaign}{mass\_kg\_final}{0.5}
\goldrow{kpi}{campaign}{mass\_kg\_initial}{0.5}
\goldrow{kpi}{campaign}{mass\_residual\_kg}{2.22044604925e-16}
\goldrow{kpi}{campaign}{moles\_kmol\_external\_out}{0}
\goldrow{kpi}{campaign}{moles\_kmol\_final}{0.01}
\goldrow{kpi}{campaign}{moles\_kmol\_initial}{0.01}
\goldrow{kpi}{campaign}{transfers\_H\_invalid}{0}
\goldrow{kpi}{campaign}{transfers\_H\_valid}{0}
\goldrow{kpi}{campaign}{transfers\_logged}{0}
\goldmore{8}
\end{tutgolden}

\clearpage
\tutentry{batch33\_fast\_intermediate\_reduced}{tutorials/batch/reactor/batch33\_fast\_intermediate\_reduced}
\begin{tutscheme}
\begin{tikzpicture}[>={Stealth[length=2mm]}, every node/.style={font=\scriptsize}]
  \node[draw, rounded corners=2pt, align=center, minimum width=1.9cm, minimum height=0.85cm, fill=black!4] (reactor) at (0.0,-0.0) {\textbf{reactor}\\[-1pt]{\tiny batchReactor}};
\end{tikzpicture}
\end{tutscheme}
\tutwhat{The same reactor reduced by the quasi-steady intermediate: A $\to$ P at the rate A makes I, written as one expression law}
\begin{tutnarrative}
the REDUCED twin of batch32\_fast\_intermediate\_full.

THE REDUCTION. When I is consumed much faster than it is made, its two rates balance, k1 n\_A = k2 n\_I, and P is made at the rate A is used:

This case runs that one law, A $\to$ P, written as an `expression` law in constant/reactions. compB (the intermediate) stays declared and stays at zero: the reduced model does not carry it.

WHAT IS KEPT. One rate constant, k1, and the two stable species. k2 does not appear: it only decides how little I there is, never how fast P is made. What is lost is the I in transit, which batch32 measures.

\end{tutnarrative}
\tutreadme{The reduced twin of batch32\_fast\_intermediate\_full: one law, A $\to$ P at k1 c\_A, written in constant/reactions as a type expression law. The intermediate is not carried; k2 does not appear.}
\tutreadme{1. One rate constant in the stable species alone. 2. What is lost. The I in transit: 0.95 \% of the charge at most, against the full twin.}
\begin{tutgolden}
\goldrow{kpi}{campaign}{energy\_balance\_available}{0}
\goldrow{kpi}{campaign}{energy\_records\_invalid}{1}
\goldrow{kpi}{campaign}{energy\_records\_valid}{0}
\goldrow{kpi}{campaign}{mass\_closure\_rel}{4.4408920985e-16}
\goldrow{kpi}{campaign}{mass\_kg\_external\_out}{0}
\goldrow{kpi}{campaign}{mass\_kg\_final}{0.5}
\goldrow{kpi}{campaign}{mass\_kg\_initial}{0.5}
\goldrow{kpi}{campaign}{mass\_residual\_kg}{2.22044604925e-16}
\goldrow{kpi}{campaign}{moles\_kmol\_external\_out}{0}
\goldrow{kpi}{campaign}{moles\_kmol\_final}{0.01}
\goldrow{kpi}{campaign}{moles\_kmol\_initial}{0.01}
\goldrow{kpi}{campaign}{transfers\_H\_invalid}{0}
\goldrow{kpi}{campaign}{transfers\_H\_valid}{0}
\goldrow{kpi}{campaign}{transfers\_logged}{0}
\goldmore{8}
\end{tutgolden}

\clearpage
\tutentry{batch34\_fractional\_order\_at\_zero}{tutorials/batch/reactor/batch34\_fractional\_order\_at\_zero}
\begin{tutscheme}
\begin{tikzpicture}[>={Stealth[length=2mm]}, every node/.style={font=\scriptsize}]
  \node[draw, rounded corners=2pt, align=center, minimum width=1.9cm, minimum height=0.85cm, fill=black!4] (reactor) at (0.0,-0.0) {\textbf{reactor}\\[-1pt]{\tiny batchReactor}};
\end{tikzpicture}
\end{tutscheme}
\tutwhat{Isothermal batch reactor: A $\to$ B $\to$ C, B absent at t = 0 and consumed at order 1/2 (the analytic Jacobian at c = 0, C468)}
\begin{tutnarrative}
the witness of C468 defect 1:

A --k1-$\to$ B --k2 c\_B\textasciicircum{}0.5-$\to$ C

B is ABSENT at t = 0 and is consumed at the fractional order 1/2. The slope d(c\textasciicircum{}0.5)/dc = 0.5 c\textasciicircum{}-0.5 is infinite at c = 0, so a Newton matrix cannot hold it. Until 2026-11-22 the analytic Jacobian differenced such a reaction at its own step and the audit, differencing at ITS step, read a different number across the singularity and refused the run (d(n\_compB)/d(n\_compB) -2.678 against -2.671). The Jacobian now takes the slope at the smallest concentration it resolves, c\_floor = sqrt(eps) s / V, and the audit takes that reaction's share of the column out of the analytic Jacobian and out of its reference (the share at the singularity has no reference), auditing the rest (Theory Guide sec:batch-analytic-jacobian). The answer is the same within the error estimate: a Rosenbrock method has no Newton iteration, the Jacobian enters the step formula itself, so a different J moves the step a little, never the solution beyond the tolerance.

Isothermal 350 K, constant V; mass and element closure are gated as for every batch case. k2 = 2e-3 (m3/kmol)\textasciicircum{}0.5 / s.

\end{tutnarrative}
\begin{tutgolden}
\goldrow{kpi}{campaign}{energy\_balance\_available}{0}
\goldrow{kpi}{campaign}{energy\_records\_invalid}{1}
\goldrow{kpi}{campaign}{energy\_records\_valid}{0}
\goldrow{kpi}{campaign}{mass\_closure\_rel}{1.7763568394e-15}
\goldrow{kpi}{campaign}{mass\_kg\_external\_out}{0}
\goldrow{kpi}{campaign}{mass\_kg\_final}{0.5}
\goldrow{kpi}{campaign}{mass\_kg\_initial}{0.5}
\goldrow{kpi}{campaign}{mass\_residual\_kg}{8.881784197e-16}
\goldrow{kpi}{campaign}{moles\_kmol\_external\_out}{0}
\goldrow{kpi}{campaign}{moles\_kmol\_final}{0.01}
\goldrow{kpi}{campaign}{moles\_kmol\_initial}{0.01}
\goldrow{kpi}{campaign}{transfers\_H\_invalid}{0}
\goldrow{kpi}{campaign}{transfers\_H\_valid}{0}
\goldrow{kpi}{campaign}{transfers\_logged}{0}
\goldmore{8}
\end{tutgolden}

\subsection{recipes}
\grouptheory{recipes}
\clearpage
\tutentry{recipe01\_react\_then\_distill}{tutorials/batch/recipes/recipe01\_react\_then\_distill}
\begin{tutscheme}
\begin{tikzpicture}[>={Stealth[length=2mm]}, every node/.style={font=\scriptsize}]
  \node[draw, rounded corners=2pt, align=center, minimum width=1.9cm, minimum height=0.85cm, fill=black!4] (reactor) at (0.0,-0.0) {\textbf{reactor}\\[-1pt]{\tiny batchReactor}};
  \node[draw, rounded corners=2pt, align=center, minimum width=1.9cm, minimum height=0.85cm, fill=black!4] (receiver) at (0.0,-1.7) {\textbf{receiver}\\[-1pt]{\tiny batchAccumulator}};
  \node[draw, rounded corners=2pt, align=center, minimum width=1.9cm, minimum height=0.85cm, fill=black!4] (still) at (0.0,-3.4) {\textbf{still}\\[-1pt]{\tiny batchStill}};
\end{tikzpicture}
\end{tutscheme}
\tutwhat{Two-vessel batch recipe: esterify (AcOH + EtOH $\to$ EtOAc + H2O) in a reactor, transfer the charge to a still, then distil}
\begin{tutnarrative}
the canonical 2-stage batch recipe.

Stage 1 (t = 0 $\to$ 400 s): A batch reactor at 350 K esterifies ethanol + acetic acid $\to$ ethyl acetate + water (the same first-order Arrhenius reaction as batch01). k $\approx$ 3.58e-3 s?$^1$ at 350 K, tau $\approx$ 280 s. Over 400 s the ethanol conversion reaches \textasciitilde{}76 \%.

Stage 2 (t = 400 $\to$ 1200 s): The recipe TRANSFERs the reactor contents to an initially empty batch still. The still then runs a Rayleigh distillation at F\_vap = 2.0e-5 kmol/s for 800 s, vaporising 0.016 kmol (\textasciitilde{}85 \% of the charge) -- the condensed distillate is routed, step by step, into a passive batchAccumulator (`dischargeTo receiver;`), so the campaign ends with the ester-rich distillate IN a vessel while the ethanol-depleted residue stays in the still's own pot (the pot IS the bottoms tank; no third transfer needed). Mass is conserved across the three vessels at every accepted step. Starting from the converted mix (ethyl acetate + water + unreacted ethanol/acid), the still progressively strips the more volatile ethyl acetate + ethanol; pot composition follows the Raoult y-x curve. (The thermoPhysPropDict here uses ideal Raoult for pedagogical simplicity.)

This case demonstrates the recipe layer (v0.23): a list of time-triggered events acting on named units. Each event in the `recipe` block below has a `time`, an `action`, and action-specific keys. Actions implemented in v0.23:

transfer move all contents of `from` into `to` setParameter change a named operating parameter on a unit

Events fire at the next time-loop step boundary after t\_trigger; the trajectory CSV thus shows a single-step discontinuity in the transferred quantities.

\end{tutnarrative}
\begin{tutgolden}
\goldrow{kpi}{campaign}{element\_C\_closure\_rel}{1.31276371155e-15}
\goldrow{kpi}{campaign}{element\_H\_closure\_rel}{1.2002411077e-15}
\goldrow{kpi}{campaign}{element\_O\_closure\_rel}{1.12522603847e-15}
\goldrow{kpi}{campaign}{element\_worst\_closure\_rel}{1.31276371155e-15}
\goldrow{kpi}{campaign}{energy\_H\_external\_kJ}{0}
\goldrow{kpi}{campaign}{energy\_Q\_ledger\_kJ}{27.0151986826}
\goldrow{kpi}{campaign}{energy\_balance\_available}{1}
\goldrow{kpi}{campaign}{energy\_closure\_rel}{4.21630243446e-11}
\goldrow{kpi}{campaign}{energy\_dH\_vessels\_kJ}{27.0151992663}
\goldrow{kpi}{campaign}{energy\_records\_invalid}{0}
\goldrow{kpi}{campaign}{energy\_records\_valid}{4}
\goldrow{kpi}{campaign}{energy\_residual\_kJ}{5.83673795518e-07}
\goldrow{kpi}{campaign}{mass\_closure\_rel}{1.13101187107e-15}
\goldrow{kpi}{campaign}{mass\_kg\_external\_out}{0}
\goldmore{66}
\end{tutgolden}

\clearpage
\tutentry{recipe02\_setpoint\_ramp}{tutorials/batch/recipes/recipe02\_setpoint\_ramp}
\begin{tutscheme}
\begin{tikzpicture}[>={Stealth[length=2mm]}, every node/.style={font=\scriptsize}]
  \node[draw, rounded corners=2pt, align=center, minimum width=1.9cm, minimum height=0.85cm, fill=black!4] (reactor) at (0.0,-0.0) {\textbf{reactor}\\[-1pt]{\tiny batchReactor}};
\end{tikzpicture}
\end{tutscheme}
\tutwhat{Single-reactor recipe with T\_setpoint steps: 320 K $\to$ 350 K @ 200 s $\to$ 380 K @ 500 s}
\begin{tutnarrative}
isothermal batch reactor with a staged temperature programme controlled by the recipe layer.

The reactor starts at 320 K, where the first-order Arrhenius reaction is very slow (k\_320 $\approx$ 3.7e-4 s?$^1$, $\tau$ $\approx$ 2700 s). Two recipe events ramp the setpoint upward:

t = 200 s $\to$ 320 K $\to$ 350 K (k jumps \textasciitilde{}10x to 3.58e-3 s?$^1$) t = 500 s $\to$ 350 K $\to$ 380 K (k jumps another \textasciitilde{}7x to 0.025 s?$^1$)

In the trajectory CSV the temperature column shows a clean step-change at each event, and the conversion curve shows the corresponding three regimes:

0 $\le$ t < 200 s barely-moving plateau (slow) 200 $\le$ t < 500 s visible drop in n\_ethanol (moderate) 500 $\le$ t rapid quench to full conversion (fast)

Demonstrates the recipe layer's `setParameter` action without inter-unit transfers, so this is the cleanest possible test of the parameter-change pathway.

\end{tutnarrative}
\begin{tutgolden}
\goldrow{kpi}{campaign}{element\_C\_closure\_rel}{3.75075346157e-16}
\goldrow{kpi}{campaign}{element\_H\_closure\_rel}{3.00060276926e-16}
\goldrow{kpi}{campaign}{element\_O\_closure\_rel}{2.50050230771e-16}
\goldrow{kpi}{campaign}{element\_worst\_closure\_rel}{3.75075346157e-16}
\goldrow{kpi}{campaign}{energy\_H\_external\_kJ}{0}
\goldrow{kpi}{campaign}{energy\_Q\_ledger\_kJ}{114.207859118}
\goldrow{kpi}{campaign}{energy\_balance\_available}{1}
\goldrow{kpi}{campaign}{energy\_closure\_rel}{0}
\goldrow{kpi}{campaign}{energy\_dH\_vessels\_kJ}{114.207859118}
\goldrow{kpi}{campaign}{energy\_records\_invalid}{0}
\goldrow{kpi}{campaign}{energy\_records\_valid}{5}
\goldrow{kpi}{campaign}{energy\_residual\_kJ}{0}
\goldrow{kpi}{campaign}{mass\_closure\_rel}{4.5240474843e-16}
\goldrow{kpi}{campaign}{mass\_kg\_external\_out}{0}
\goldmore{35}
\end{tutgolden}

\clearpage
\tutentry{recipe03\_conditional\_partial}{tutorials/batch/recipes/recipe03\_conditional\_partial}
\begin{tutscheme}
\begin{tikzpicture}[>={Stealth[length=2mm]}, every node/.style={font=\scriptsize}]
  \node[draw, rounded corners=2pt, align=center, minimum width=1.9cm, minimum height=0.85cm, fill=black!4] (reactor) at (0.0,-0.0) {\textbf{reactor}\\[-1pt]{\tiny batchReactor}};
  \node[draw, rounded corners=2pt, align=center, minimum width=1.9cm, minimum height=0.85cm, fill=black!4] (still) at (0.0,-1.7) {\textbf{still}\\[-1pt]{\tiny batchStill}};
\end{tikzpicture}
\end{tutscheme}
\tutwhat{Conditional + partial recipe: react until x\_ethanol < 0.30, then transfer 50\% to a still}
\begin{tutnarrative}
recipe extensions (v0.50).

Two recipe features beyond the time-triggered transfer/setParameter of v0.23:

1. CONDITION-TRIGGERED event (`when \{ ... \}` instead of `time`): the event fires the first step its condition becomes true, evaluated on the live vessel state. Quantities: T, total, n\_<comp>, x\_<comp>; operators: gt, ge, lt, le.

2. PARTIAL transfer (`fraction`): a fraction (0..1) of the source moves; the source keeps the rest. Default 1.0 == empty the source.

Here the reactor runs the esterification (ethanol + acetic acid $\to$ ethyl acetate + water; the batch01 first-order Arrhenius test, k \textasciitilde{} 3.58e-3 s\textasciicircum{}-1 at 350 K, so n\_ethanol = n\_ethanol(0).exp(-k t)). The charge is equimolar (x\_ethanol starts at 0.5); when the ethanol mole fraction drops below 0.30, the recipe transfers HALF the reactor contents to a still; the reactor keeps the other half and reacts on, while the still distils its charge.

\end{tutnarrative}
\begin{tutgolden}
\goldrow{kpi}{reactor}{P\_final}{101300}
\goldrow{kpi}{reactor}{T\_final}{350}
\goldrow{kpi}{reactor}{n\_aceticAcid\_final}{6.33853236066e-05}
\goldrow{kpi}{reactor}{n\_ethanol\_final}{6.33853236066e-05}
\goldrow{kpi}{reactor}{n\_ethylAcetate\_final}{0.00456161467639}
\goldrow{kpi}{reactor}{n\_water\_final}{0.00456161467639}
\goldrow{kpi}{reactor}{t\_end}{1200}
\goldrow{kpi}{reactor}{totalMoles\_final}{0.00925}
\goldrow{kpi}{still}{P\_final}{101300}
\goldrow{kpi}{still}{T\_final}{390.250742415}
\goldrow{kpi}{still}{n\_aceticAcid\_final}{0}
\goldrow{kpi}{still}{n\_ethanol\_final}{0}
\goldrow{kpi}{still}{n\_ethylAcetate\_final}{0}
\goldrow{kpi}{still}{n\_water\_final}{0}
\goldmore{58}
\end{tutgolden}

\clearpage
\tutentry{recipe04\_charge\_enthalpy}{tutorials/batch/recipes/recipe04\_charge\_enthalpy}
\begin{tutscheme}
\begin{tikzpicture}[>={Stealth[length=2mm]}, every node/.style={font=\scriptsize}]
  \node[draw, rounded corners=2pt, align=center, minimum width=1.9cm, minimum height=0.85cm, fill=black!4] (hot) at (0.0,-0.0) {\textbf{hot}\\[-1pt]{\tiny batchAccumulator}};
  \node[draw, rounded corners=2pt, align=center, minimum width=1.9cm, minimum height=0.85cm, fill=black!4] (cold) at (0.0,-1.7) {\textbf{cold}\\[-1pt]{\tiny batchAccumulator}};
\end{tikzpicture}
\end{tutscheme}
\tutwhat{H-equality recipe charge: T\_mix from the first law, not the molar average; balance CLOSES across a transfer}
\begin{tutnarrative}
mixing is a FIRST-LAW act.

Two passive tanks. `hot` holds 0.5 mol of ethanol at 340 K; `cold` holds 1.0 mol of water at 300 K. At t = 100 s the recipe empties the hot tank into the cold one. The charge solves the mix temperature from enthalpy equality on the elements datum,

H(n\_EtOH + n\_W, T\_mix) = H(n\_EtOH, 340 K) + H(n\_W, 300 K),

so heat capacity DIFFERENCES matter: ethanol carries \textasciitilde{}112 J/(mol K) against water's \textasciitilde{}75, and the enthalpy-true T\_mix (\textasciitilde{}317 K) sits \textasciitilde{}4 K above the equal-Cp molar average (313.3 K) a naive mixing rule gives.

Because both species price on the elements datum, no jacket acts, and nothing leaves, the campaign ENERGY BALANCE is AVAILABLE and closes: dH\_vessels = 0 to the mix-solve tolerance --- the transfer moved enthalpy between tanks without creating or destroying any.

\end{tutnarrative}
\begin{tutgolden}
\goldrow{kpi}{campaign}{element\_C\_closure\_rel}{0}
\goldrow{kpi}{campaign}{element\_H\_closure\_rel}{0}
\goldrow{kpi}{campaign}{element\_O\_closure\_rel}{0}
\goldrow{kpi}{campaign}{element\_worst\_closure\_rel}{0}
\goldrow{kpi}{campaign}{energy\_H\_external\_kJ}{0}
\goldrow{kpi}{campaign}{energy\_Q\_ledger\_kJ}{0}
\goldrow{kpi}{campaign}{energy\_balance\_available}{1}
\goldrow{kpi}{campaign}{energy\_closure\_rel}{1.53779787259e-09}
\goldrow{kpi}{campaign}{energy\_dH\_vessels\_kJ}{-4.20121068601e-07}
\goldrow{kpi}{campaign}{energy\_records\_invalid}{0}
\goldrow{kpi}{campaign}{energy\_records\_valid}{0}
\goldrow{kpi}{campaign}{energy\_residual\_kJ}{4.20121068601e-07}
\goldrow{kpi}{campaign}{mass\_closure\_rel}{0}
\goldrow{kpi}{campaign}{mass\_kg\_external\_out}{0}
\goldmore{24}
\end{tutgolden}

\subsection{still}
\grouptheory{still}
\clearpage
\tutentry{still01\_benzene\_toluene}{tutorials/batch/still/still01\_benzene\_toluene}
\begin{tutscheme}
\begin{tikzpicture}[>={Stealth[length=2mm]}, every node/.style={font=\scriptsize}]
  \node[draw, rounded corners=2pt, align=center, minimum width=1.9cm, minimum height=0.85cm, fill=black!4] (still) at (0.0,-0.0) {\textbf{still}\\[-1pt]{\tiny batchStill}};
\end{tikzpicture}
\end{tutscheme}
\tutwhat{Rayleigh batch distillation: benzene/toluene 50/50, 1 atm, ideal Raoult}
\begin{tutnarrative}
Rayleigh batch distillation, ideal Raoult.

Pot starts with 1 mol of equimolar benzene/toluene at 1 atm. A constant vapour offtake F\_vap = 1.5e-6 kmol/s strips the pot over 600 s --- by linear mass balance the pot should hold

i.e. L(600 s) = 1 mol - 0.9 mol = 0.1 mol (90 \% distilled).

Benzene is the more volatile component (alpha \textasciitilde{} 2.4 at 1 atm), so the still pot progressively concentrates in toluene as benzene flashes off preferentially. Classical Rayleigh integration with constant alpha gives the closed form

ln(L / L\_0) = (1 / (alpha - 1)) [ ln(x / x\_0) + alpha $\ast$ ln((1 - x\_0) / (1 - x)) ]

where x = x\_benzene in the still pot. Plotting (L/L\_0) vs x\_pot from trajectory.csv should track this analytical curve to RK4 precision.

\end{tutnarrative}
\tutreadme{One mole of 50/50 benzene/toluene in a pot at 1.013 bar, boiled off for 600 s with the vapour taken away as it forms --- no reflux, no trays: the Rayleigh equation integrated in time. The pot starts at its bubble point, 365.00 K, and ends at 380.04 K with 0.1 mol left in it at 91.8 \% toluene; 0.9 mol (0.0760 kg) has left as distillate. The reboiler supplied 30.32 kJ, the condenser took 29.38 kJ away.}
\tutreadme{1. The residue gets heavier, and hotter, because the light one leaves first. The Log tab prints the pot every 20 s: benzene falls faster than toluene, T climbs from 365 to 380 K as the pot approaches toluene's boiling point. That is the Rayleigh curve as numbers. 2. The whole campaign is a ledger. [campaign] mass balance: m0 = 0.0851 kg = mF 0.0091 + external 0.0760 kg, closure 0 --- what left the domain is counted, not lost. The energy side closes at 4.6e-16 with two records, reboiler and condenser, each priced at its own service temperature (380.0 K and 365.3 K). 3. The thermodynamics is Raoult, and the caveats say where it is stretched. Benzene's Antoine window ends at 354 K and the pot starts at 365 K; toluene's ends at 380 K and the pot ends at 380.04. Both extrapolations are announced at the moment they happen, on the console, inside the table. 4. Utilities are scheduled in time. reports/utilities/ utilityDemand.csv carries the staircase of steam and cooling-water demand across the 600 s --- a peak, not just a total, which is what a batch plant is actually sized on.}
\tutreadme{Halve endTime: the pot is still half full and barely enriched --- most of the separation happens late, when there is little left. Then still04\_rectifier\_benzene\_toluene puts three stages and a reflux over the same pot, and the distillate purity that a Rayleigh still cannot reach appears.}
\begin{tutgolden}
\goldrow{kpi}{campaign}{element\_C\_closure\_rel}{2.66880534766e-16}
\goldrow{kpi}{campaign}{element\_H\_closure\_rel}{1.23908819713e-16}
\goldrow{kpi}{campaign}{element\_worst\_closure\_rel}{2.66880534766e-16}
\goldrow{kpi}{campaign}{energy\_H\_external\_kJ}{38.7602049735}
\goldrow{kpi}{campaign}{energy\_Q\_ledger\_kJ}{0.943056487975}
\goldrow{kpi}{campaign}{energy\_balance\_available}{1}
\goldrow{kpi}{campaign}{energy\_closure\_rel}{6.41611564455e-16}
\goldrow{kpi}{campaign}{energy\_dH\_vessels\_kJ}{-37.8171484855}
\goldrow{kpi}{campaign}{energy\_records\_invalid}{0}
\goldrow{kpi}{campaign}{energy\_records\_valid}{2}
\goldrow{kpi}{campaign}{energy\_residual\_kJ}{2.48689957516e-14}
\goldrow{kpi}{campaign}{mass\_closure\_rel}{0}
\goldrow{kpi}{campaign}{mass\_kg\_external\_out}{0.0760287317788}
\goldrow{kpi}{campaign}{mass\_kg\_final}{0.00909876822121}
\goldmore{35}
\end{tutgolden}

\clearpage
\tutentry{still02\_ethanol\_water\_azeo}{tutorials/batch/still/still02\_ethanol\_water\_azeo}
\begin{tutscheme}
\begin{tikzpicture}[>={Stealth[length=2mm]}, every node/.style={font=\scriptsize}]
  \node[draw, rounded corners=2pt, align=center, minimum width=1.9cm, minimum height=0.85cm, fill=black!4] (still) at (0.0,-0.0) {\textbf{still}\\[-1pt]{\tiny batchStill}};
\end{tikzpicture}
\end{tutscheme}
\tutwhat{Rayleigh batch distillation: ethanol/water 40/60 with NRTL, residue concentrates in water}
\begin{tutnarrative}
Rayleigh batch distillation of a non-ideal binary with an azeotrope.

System: ethanol / water with NRTL, pressure 1 atm. Minimum-boiling positive azeotrope at x\_eth \textasciitilde{} 0.894 (mole), T \textasciitilde{} 351.3 K.

Pot starts at x\_eth = 0.40 (a "wine"-like feed --- well BELOW the azeotrope). Boil-off rate F\_vap = 1.5e-6 kmol/s for 600 s consumes 90 \% of the pot, leaving a 10 \% residue.

Below the azeotrope, gamma\_ethanol > 1 and gamma\_water \textasciitilde{} 1, so the vapour is markedly richer in ethanol than the liquid (y\_eth \textasciitilde{} 0.61 at x\_eth = 0.40). Ethanol therefore vaporises preferentially and the still pot progressively concentrates in WATER --- i.e. the residue at end-of-run is mostly water.

The pedagogical contrast with still01 (ideal Raoult, benzene/toluene) is the SHAPE of the enrichment curve:

Raoult: smooth, alpha-driven (alpha \textasciitilde{} 2.4, roughly constant). NRTL : curvature changes as x crosses concentration regimes in which gamma\_ethanol shrinks toward 1; the still composition curve x\_pot(t) is noticeably non-linear vs the linear pot-volume curve L(t) = L\_0 - F\_vap$\ast$t.

Crucially, this case demonstrates that a single Rayleigh still cannot cross the azeotrope. Starting below 0.894 in ethanol, the pot can only move DOWNWARD in ethanol fraction; no amount of boiling will lift the pot composition past 0.894. (To break the azeotrope one needs entrainer, pressure swing, or extractive distillation --- topics for later tutorials.)

\end{tutnarrative}
\begin{tutgolden}
\goldrow{kpi}{still}{P\_final}{101300}
\goldrow{kpi}{still}{T\_final}{373.117313962}
\goldrow{kpi}{still}{n\_ethanol\_final}{2.08003362152e-11}
\goldrow{kpi}{still}{n\_water\_final}{9.99999650927e-05}
\goldrow{kpi}{still}{t\_end}{600}
\goldrow{kpi}{still}{totalMoles\_final}{9.99999989764e-05}
\goldrow{kpi}{campaign}{element\_C\_closure\_rel}{1.35525271561e-15}
\goldrow{kpi}{campaign}{element\_H\_closure\_rel}{9.63735264432e-16}
\goldrow{kpi}{campaign}{element\_O\_closure\_rel}{6.50521303491e-16}
\goldrow{kpi}{campaign}{element\_worst\_closure\_rel}{1.35525271561e-15}
\goldrow{kpi}{campaign}{energy\_H\_external\_kJ}{-249.365160058}
\goldrow{kpi}{campaign}{energy\_Q\_ledger\_kJ}{0.522318304026}
\goldrow{kpi}{campaign}{energy\_balance\_available}{1}
\goldrow{kpi}{campaign}{energy\_closure\_rel}{1.13740334008e-16}
\goldmore{36}
\end{tutgolden}

\clearpage
\tutentry{still03\_with\_receiver}{tutorials/batch/still/still03\_with\_receiver}
\begin{tutscheme}
\begin{tikzpicture}[>={Stealth[length=2mm]}, every node/.style={font=\scriptsize}]
  \node[draw, rounded corners=2pt, align=center, minimum width=1.9cm, minimum height=0.85cm, fill=black!4] (still) at (0.0,-0.0) {\textbf{still}\\[-1pt]{\tiny batchStill}};
  \node[draw, rounded corners=2pt, align=center, minimum width=1.9cm, minimum height=0.85cm, fill=black!4] (receiver) at (0.0,-1.7) {\textbf{receiver}\\[-1pt]{\tiny batchAccumulator}};
\end{tikzpicture}
\end{tutscheme}
\tutwhat{Rayleigh still distilling INTO a receiver tank (dischargeTo accumulator)}
\begin{tutnarrative}
Rayleigh batch distillation INTO a receiver (accumulator vessel + continuous dischargeTo).

Same pot as still01 (1 mol equimolar benzene/toluene, F\_vap = 1.5e-6 kmol/s, ideal Raoult), but the still's vapour is no longer lost: the `dischargeTo receiver;` line routes the condensed distillate, step by step, into a passive `batchAccumulator`. So the trajectory shows the POT depleting AND the RECEIVER filling --- the receiver collects the total distillate (amount + composition). Mass is conserved between the two vessels: pot + receiver = 1 mol at all times. Benzene is the more volatile component, so the receiver is benzene-rich early and the cumulative receiver composition drifts toward 0.5 as distillation nears completion (it ends holding everything the pot lost).

\end{tutnarrative}
\begin{tutgolden}
\goldrow{kpi}{campaign}{element\_C\_closure\_rel}{2.66880534766e-16}
\goldrow{kpi}{campaign}{element\_H\_closure\_rel}{1.23908819713e-16}
\goldrow{kpi}{campaign}{element\_worst\_closure\_rel}{2.66880534766e-16}
\goldrow{kpi}{campaign}{energy\_H\_external\_kJ}{0}
\goldrow{kpi}{campaign}{energy\_Q\_ledger\_kJ}{0.943056487975}
\goldrow{kpi}{campaign}{energy\_balance\_available}{1}
\goldrow{kpi}{campaign}{energy\_closure\_rel}{2.96217392064e-11}
\goldrow{kpi}{campaign}{energy\_dH\_vessels\_kJ}{0.943056485707}
\goldrow{kpi}{campaign}{energy\_records\_invalid}{0}
\goldrow{kpi}{campaign}{energy\_records\_valid}{2}
\goldrow{kpi}{campaign}{energy\_residual\_kJ}{2.26835439321e-09}
\goldrow{kpi}{campaign}{mass\_closure\_rel}{0}
\goldrow{kpi}{campaign}{mass\_kg\_external\_out}{0}
\goldrow{kpi}{campaign}{mass\_kg\_final}{0.0851275}
\goldmore{41}
\end{tutgolden}

\clearpage
\tutentry{still04\_rectifier\_benzene\_toluene}{tutorials/batch/still/still04\_rectifier\_benzene\_toluene}
\begin{tutscheme}
\begin{tikzpicture}[>={Stealth[length=2mm]}, every node/.style={font=\scriptsize}]
  \node[draw, rounded corners=2pt, align=center, minimum width=1.9cm, minimum height=0.85cm, fill=black!4] (still) at (0.0,-0.0) {\textbf{still}\\[-1pt]{\tiny batchStill}};
  \node[draw, rounded corners=2pt, align=center, minimum width=1.9cm, minimum height=0.85cm, fill=black!4] (receiver) at (0.0,-1.7) {\textbf{receiver}\\[-1pt]{\tiny batchAccumulator}};
\end{tikzpicture}
\end{tutscheme}
\tutwhat{batch rectification: 3 stages + total condenser, R=3, with receiver tank (dischargeTo accumulator)}
\begin{tutnarrative}
the CLASSICAL batch rectification (model rectifier, design forum \#85): the Rayleigh pot now feeds a column of 3 ideal stages topped by a total condenser at constant reflux R = 3.

F\_vap is the internal BOILUP V = 1.5e-6 kmol/s; the distillate PRODUCT is D = V/(R+1) = 3.75e-7 kmol/s and the reflux L = R.D returns to the column. The pot loses PRODUCT only: dn\_i/dt = -D.x\_D\_i, with x\_D(t) solved at every instant from the quasi-steady cascade (stage balances + bubble-T equilibrium + total-condenser closure -- a declared residual, printed on failure).

What to watch (trajectory.csv): still.xD\_benzene starts \textasciitilde{}0.949 (vs \textasciitilde{}0.71 for the raw pot vapour -- the column working) and DRIFTS DOWN as the pot strips: the constant-R policy trades purity for time. The receiver's ACCUMULATED benzene fraction sits between the instantaneous x\_D and the feed -- two different compositions, deliberately both visible. The pot loses only D.dt = 0.225 mmol over 600 s (vs 0.9 mmol in still01: same boilup, but 3/4 of it refluxes).

Physics anchors (verified at authoring): x\_D(0) grows with nStages and R (N=1,R=1: 0.806; N=2,R=3: 0.909; N=3,R=3: 0.949) and the R $\to$ infinity limit (R=200: 0.9765) lands inside the Fenske band alpha\textasciicircum{}(N+1) for alpha = 2.35-2.6. With nStages 0 and refluxRatio 0 the model reproduces still01's Rayleigh trajectory to the last digit.

\end{tutnarrative}
\begin{tutgolden}
\goldrow{kpi}{campaign}{element\_C\_closure\_rel}{0}
\goldrow{kpi}{campaign}{element\_H\_closure\_rel}{0}
\goldrow{kpi}{campaign}{element\_worst\_closure\_rel}{0}
\goldrow{kpi}{campaign}{energy\_H\_external\_kJ}{0}
\goldrow{kpi}{campaign}{energy\_Q\_ledger\_kJ}{0.534521976282}
\goldrow{kpi}{campaign}{energy\_balance\_available}{1}
\goldrow{kpi}{campaign}{energy\_closure\_rel}{7.05914882685e-10}
\goldrow{kpi}{campaign}{energy\_dH\_vessels\_kJ}{0.534521958795}
\goldrow{kpi}{campaign}{energy\_records\_invalid}{0}
\goldrow{kpi}{campaign}{energy\_records\_valid}{2}
\goldrow{kpi}{campaign}{energy\_residual\_kJ}{1.74866894298e-08}
\goldrow{kpi}{campaign}{mass\_closure\_rel}{0}
\goldrow{kpi}{campaign}{mass\_kg\_external\_out}{0}
\goldrow{kpi}{campaign}{mass\_kg\_final}{0.0851275}
\goldmore{46}
\end{tutgolden}

\clearpage
\tutentry{still05\_rectifier\_constant\_xD}{tutorials/batch/still/still05\_rectifier\_constant\_xD}
\begin{tutscheme}
\begin{tikzpicture}[>={Stealth[length=2mm]}, every node/.style={font=\scriptsize}]
  \node[draw, rounded corners=2pt, align=center, minimum width=1.9cm, minimum height=0.85cm, fill=black!4] (still) at (0.0,-0.0) {\textbf{still}\\[-1pt]{\tiny batchStill}};
  \node[draw, rounded corners=2pt, align=center, minimum width=1.9cm, minimum height=0.85cm, fill=black!4] (receiver) at (0.0,-1.7) {\textbf{receiver}\\[-1pt]{\tiny batchAccumulator}};
\end{tikzpicture}
\end{tutscheme}
\tutwhat{batch rectification, constant-composition policy: R(t) holds x\_D until refluxMax, then D=0}
\begin{tutnarrative}
the SECOND classical reflux policy (design forum \#85-3/4): hold the distillate composition constant by RAISING the reflux as the pot strips, and STOP when the declared ceiling can no longer reach the target.

Same pot and column as still04 (1 mol equimolar benzene/toluene, V = 1.5e-6 kmol/s, 3 ideal stages + total condenser), but instead of a fixed R the policy solves R(t) each step -- a BRACKETED solve over [0, refluxMax] on a feasibility map, never a bare Newton -- so that x\_D[benzene] = 0.93 exactly. As benzene depletes, R climbs (\textasciitilde{}2 at the start); when even refluxMax = 8 cannot reach the target the run announces `onRefluxLimit stopDistillation`: D = 0, the pot stops losing mass, and the campaign idles VISIBLY to endTime with still.refluxLimitReached = 1 in the trajectory -- never a magic termination.

What to watch (trajectory.csv): still.xD\_benzene pinned at 0.9300 while R climbs -- the mirror image of still04, which held R and let purity drift; still.R rising 2.0 $\to$ 8 (the ceiling), then the limit event; the receiver's accumulated composition \textasciitilde{}0.93 (the whole cut is on-spec -- THE selling point of this policy); pot inventory constant after t\_refluxLimit.

Declared refusals (each with its own diagnostic, try them): x\_D\_target 0.5 $\to$ BELOW the R=0 floor (even zero reflux over-purifies); x\_D\_target 0.995 with refluxMax 5 $\to$ UNREACHABLE from the charge; ethanol/water NRTL with x\_D\_target 0.95 $\to$ the 0.894 azeotrope caps the cascade at ANY reflux (the barrier is real VLE, not a message).

\end{tutnarrative}
\begin{tutgolden}
\goldrow{kpi}{campaign}{element\_C\_closure\_rel}{0}
\goldrow{kpi}{campaign}{element\_H\_closure\_rel}{0}
\goldrow{kpi}{campaign}{element\_worst\_closure\_rel}{0}
\goldrow{kpi}{campaign}{energy\_H\_external\_kJ}{0}
\goldrow{kpi}{campaign}{energy\_Q\_ledger\_kJ}{0.690487726727}
\goldrow{kpi}{campaign}{energy\_balance\_available}{1}
\goldrow{kpi}{campaign}{energy\_closure\_rel}{1.51572395207e-08}
\goldrow{kpi}{campaign}{energy\_dH\_vessels\_kJ}{0.690487258335}
\goldrow{kpi}{campaign}{energy\_records\_invalid}{0}
\goldrow{kpi}{campaign}{energy\_records\_valid}{2}
\goldrow{kpi}{campaign}{energy\_residual\_kJ}{4.68392697073e-07}
\goldrow{kpi}{campaign}{mass\_closure\_rel}{0}
\goldrow{kpi}{campaign}{mass\_kg\_external\_out}{0}
\goldrow{kpi}{campaign}{mass\_kg\_final}{0.0851275}
\goldmore{48}
\end{tutgolden}

\clearpage
\tutentry{still06\_ledger\_mixed\_validity}{tutorials/batch/still/still06\_ledger\_mixed\_validity}
\begin{tutscheme}
\begin{tikzpicture}[>={Stealth[length=2mm]}, every node/.style={font=\scriptsize}]
  \node[draw, rounded corners=2pt, align=center, minimum width=1.9cm, minimum height=0.85cm, fill=black!4] (potA) at (0.0,-0.0) {\textbf{potA}\\[-1pt]{\tiny batchStill}};
  \node[draw, rounded corners=2pt, align=center, minimum width=1.9cm, minimum height=0.85cm, fill=black!4] (recA) at (0.0,-1.7) {\textbf{recA}\\[-1pt]{\tiny batchAccumulator}};
  \node[draw, rounded corners=2pt, align=center, minimum width=1.9cm, minimum height=0.85cm, fill=black!4] (dump) at (0.0,-3.4) {\textbf{dump}\\[-1pt]{\tiny batchAccumulator}};
  \node[draw, rounded corners=2pt, align=center, minimum width=1.9cm, minimum height=0.85cm, fill=black!4] (fresh) at (0.0,-5.1) {\textbf{fresh}\\[-1pt]{\tiny batchAccumulator}};
  \node[draw, rounded corners=2pt, align=center, minimum width=1.9cm, minimum height=0.85cm, fill=black!4] (potB) at (0.0,-6.8) {\textbf{potB}\\[-1pt]{\tiny batchStill}};
  \node[draw, rounded corners=2pt, align=center, minimum width=1.9cm, minimum height=0.85cm, fill=black!4] (recB) at (0.0,-8.5) {\textbf{recB}\\[-1pt]{\tiny batchAccumulator}};
  \node[draw, rounded corners=2pt, align=center, minimum width=1.9cm, minimum height=0.85cm, fill=black!4] (spike) at (0.0,-10.2) {\textbf{spike}\\[-1pt]{\tiny batchAccumulator}};
\end{tikzpicture}
\end{tutscheme}
\tutwhat{Ledger H-validity is MONOTONIC: mixed valid/invalid packages never resurrect a poisoned record}
\begin{tutnarrative}
MONOTONIC enthalpy validity on the material ledger (the adversarial case of design forum \#99/\#101).

The campaign ledger integrates each transfer's enthalpy PACKAGE BY PACKAGE at each package's own temperature. A package containing a species with NO formation datum (here `compA`, a lumped teaching species that deliberately carries no standardThermochemistry block) cannot be priced --- the record must mark H invalid and NAME the species. The contract under test is the temporal one:

once ONE package is unpriceable, the record's integrated H has lost a parcel FOREVER. Validity is MONOTONIC: a later priceable package must never resurrect it.

Two stills exercise the two orders on ONE accumulated record each:

potA (invalid $\to$ valid): starts with 10 \% compA, so its distillate record potA$\to$recA begins UNPRICEABLE. At t = 200 s the recipe dumps the whole pot to `dump` and recharges clean benzene/toluene from `fresh`; every later distillate package is priceable. The record must STAY invalid (H\_missing = compA), and the discrete potA$\to$dump record is invalid too (it carried the compA).

potB (valid $\to$ invalid): starts clean, so its distillate record potB$\to$recB begins priceable and accumulates real enthalpy. At t = 200 s the recipe spikes it with pure compA from `spike`; every later package is unpriceable. The record flips invalid and the partial integral is WITHHELD --- an emitted H\_kJ that silently lost its tail would be a lie.

Five ledger records result: potA$\to$recA (continuous, poisoned early), potB$\to$recB (continuous, poisoned late), potA$\to$dump and spike$\to$potB (discrete, carrying compA: unpriceable), fresh$\to$potA (discrete, clean: the ONE priceable record). The golden pins the campaign census transfers\_H\_valid = 1 / transfers\_H\_invalid = 4; the results JSON carries H\_missing = [ compA ] and NO H\_kJ key on every poisoned record.

\end{tutnarrative}
\begin{tutgolden}
\goldrow{kpi}{campaign}{mass\_closure\_rel}{2.1198000621e-16}
\goldrow{kpi}{campaign}{mass\_kg\_external\_out}{0}
\goldrow{kpi}{campaign}{mass\_kg\_final}{0.26186975}
\goldrow{kpi}{campaign}{mass\_kg\_initial}{0.26186975}
\goldrow{kpi}{campaign}{mass\_residual\_kg}{5.55111512313e-17}
\goldrow{kpi}{campaign}{transfers\_H\_invalid}{4}
\goldrow{kpi}{campaign}{transfers\_H\_valid}{1}
\goldrow{kpi}{campaign}{transfers\_logged}{5}
\goldrow{kpi}{dump}{P\_final}{101300}
\goldrow{kpi}{dump}{T\_final}{372.484229602}
\goldrow{kpi}{dump}{n\_benzene\_final}{0.000249727144577}
\goldrow{kpi}{dump}{n\_compA\_final}{9.64864338217e-05}
\goldrow{kpi}{dump}{n\_toluene\_final}{0.000353786421588}
\goldrow{kpi}{dump}{t\_end}{600}
\goldmore{90}
\end{tutgolden}

